\documentclass[11pt]{amsart}
\usepackage[margin=2.5cm]{geometry}
\usepackage{amssymb,amsmath}
\usepackage{booktabs,adjustbox,array}
\usepackage{cite}
\usepackage{xcolor}
\usepackage[colorlinks=true,linkcolor=magenta,urlcolor=magenta,citecolor=magenta]{hyperref}
\newtheorem{theorem}{Theorem}[section]
\newtheorem{proposition}[theorem]{Proposition}
\newtheorem{lemma}[theorem]{Lemma}
\newtheorem*{thmA}{Theorem A}
\newtheorem*{thmB}{Theorem B}
\newtheorem*{thmC}{Theorem C}
\newcommand{\T}{\mathbb T}
\setcounter{MaxMatrixCols}{30}
\newcommand{\arx}[1]{{\hypersetup{allcolors=black}\href{https://arxiv.org/pdf/#1}{arXiv:#1}}}
\title[Complex Hadamard matrices]{Higher invariants of complex Hadamard matrices}
\author{T. Banica}
\author{C. Opus}
\begin{document}
\subjclass[2020]{46L37 (05B20, 46L65)}
\keywords{Complex Hadamard matrix, quantum permutation group, subfactor}
\begin{abstract}
Associated to any complex Hadamard matrix $H\in M_N(\mathbb T)$ is a quantum permutation group $G\subset_uS_N^+$, the Hopf image of the representation associated to the magic matrix formed by the rank 1 projections $P_{ij}=Proj(H_i/H_j)$, and its invariants $c_k=\dim(Fix(u^{\otimes k}))$ are the dimensions of the relative commutants of the subfactor associated to $H$. We compute these invariants well beyond the levels $k=2,3$ previously known, up to $k=9$ in the best cases, by using Jones--Wenzl and annular decompositions. For the Di\c t\u a deformations of the Fourier matrices we compute $G$ for all values of the parameters, and we study the order 8 families $F_{8,5}$, in connection with free wreath products. For the matrices of small order with $c_2=2$, such as the Karlsson ones at generic parameters, we find $c_k=Cat_k$ at all the levels computed. Finally, via an integrality principle and character sums, we prove that the Paley II matrices, and more generally those coming from symmetric conference matrices, have infinite depth, as well as the Paley I matrices with $11\leq q<10^5$.
\end{abstract}
\maketitle
\setcounter{tocdepth}{1}
\tableofcontents

\section*{Introduction}

For a complex Hadamard matrix $H\in M_N(\T)$ with rows $H_1,\ldots,H_N$, the rank 1 projections $P_{ij}=Proj(H_i/H_j)$ form a magic matrix, and the quantum group $G=G_H\subset S_N^+$ is the Hopf image of the representation associated to this magic matrix \cite{bb1,bbs}. The planar algebra of the subfactor associated to $H$ by Popa and Jones \cite{jon,pop} is the fixed point planar algebra of $G$ \cite{bbs}, so the numbers
$$c_k=\dim Fix(u^{\otimes k})=\int_G\chi^k$$
are the dimensions of the relative commutants, up to an index shift. The subfactor has finite depth precisely when $G$ is finite, and the simplest possible answer, the ``only Temperley--Lieb'' one, is $G=S_N^+$, which amounts to $c_k=Cat_k$ for all $k$. We would like to compute $G$, or as much as possible about it, for as many $H$ as possible. The previous literature stops at $k=2,3$, rarely $4$ \cite{cn,ni}, with the notable exception of the infinite depth results of Montgomery for the Paley II matrices and for $P_7$ \cite{mo2,mo3}, obtained with the angle operator \cite{mo1,swa}.

We compute here the invariants $c_k$ well beyond these levels, for a large basis of complex and real matrices, up to $k=9$ in the best cases. The computations themselves use a number of tools, such as rotation sectors, probabilistic counts with noncrossing seeds, the duality between $H$ and $H^t$, and above all the decomposition of the truncated operators $T_k$ into Jones--Wenzl pieces. Then we use this data as a guide in 3 directions: the Di\c t\u a deformations, where $G$ can be computed, the wild matrices, those of small order like the Karlsson matrices \cite{kar}, which appear to have $G=S_N^+$, and the real matrices, mainly the Paley ones, where infinite depth can be proved.

\medskip

\noindent\textbf{Di\c t\u a deformations.} For the deformed Fourier matrices $F_M\otimes_QF_N$ the quantum group was computed for generic $Q$ in \cite{bb3}, and Bichon developed in \cite{bic} the tools for the non-generic case. Let $X,Y$ be finite abelian groups, let $\Gamma_{X,Y}=\mathbb Z^{(|X|-1)(|Y|-1)}\rtimes Y$ be the group from \cite{bb3}, with the lattice part identified with $I_X\otimes I_Y$, the product of the augmentation ideals, and encode a dephased parameter matrix $Q\in M_{X\times Y}(\T)$ into an element $\theta\in(\mathbb R/\mathbb Z)[X\times Y]$. We obtain:

\begin{thmA}
For any $X,Y$ and any $Q$, the quantum group of $F_X\otimes_QF_Y$ is $C^*(\Gamma_{X,Y}/N_Q)\rtimes_\Phi C(X)$, where $N_Q$ consists of the elements $r\in I_X\otimes I_Y$ such that $r\theta$ is a sum of a function of $X$ and a function of $Y$. In particular this quantum group is the generic one if and only if, for each Galois orbit of pairs of nontrivial characters of $X,Y$, the vector of character sums of the exponents of the parameters which are not roots of unity is nonzero, and when $Q$ has roots of unity as entries, it is finite, of explicit dimension.
\end{thmA}

This is Theorem \ref{thm:cyclic} below, which is made explicit for $F_2\otimes_QF_n$ and $F_n\otimes_QF_2$, and in particular at orders 6 and 8, where the generic locus has a nontrivial structure, even away from the roots of unity. It also covers the real case, as for instance the matrix $\mathrm{had}_{16.1}$ from McKay's library \cite{mck}, whose quantum group is finite, of dimension 256. For the Di\c t\u a matrices which are not of deformed Fourier type, our main examples are the regular families $F_{8,5a},F_{8,5b},F_{8,5c}$ of order 8, built from 2 matrices $F_4^s$. Here we show that the quantum groups are contained in the free wreath products $S_4^+\wr_*\mathbb Z_2$ and $H_2^+\wr_*\mathbb Z_2$, the latter being the quantum automorphism group of the binary tree \cite{dkr,dfk}, with equality at all the levels computed, we reduce the equality to a freeness question for 2 copies of a quantum group, and we exclude the known intermediate objects, the quotients of Izumi, Morrison and Penneys type \cite{imp}, and the free wreath products with amalgamation of Freslon \cite{fre}. In subfactor terms, the question is whether the planar algebra is a Bisch--Jones free composition \cite{bjo,twa}.

\medskip

\noindent\textbf{Wild matrices.} Let us call wild the complex Hadamard matrices with $G_H=S_N^+$. The matrices of Di\c t\u a type have $c_2\geq3$, and for all the matrices that we tested with $c_2=2$, we find $c_k=Cat_k$ at all the levels computed. This extends the observations of Camp and Nicoara at levels 2 and 3 \cite{cn}. The main tool for the computations is as follows:

\begin{thmB}
For $H\in M_N(\T)$ with $N\geq5$, we have $T_k\simeq\bigoplus_sm_k(s)\Phi_s$, where $m_k(s)$ are the multiplicities of the irreducible corepresentations $r_s$ of $S_N^+$ in $u^{\otimes k}$, and $\Phi_s\simeq\Psi_2\oplus\ldots\oplus\Psi_s\oplus[s\ {\rm odd}]\cdot0_{N-1}$ for $s\geq1$, where $\Psi_t$ is the restriction of $T_t$ to the annular lowest weight space $\Lambda_t\subset(\mathbb C^N)^{\otimes t}$. In particular $G_H=S_N^+$ if and only if $\|\Psi_t\|<1$ for all $t\geq2$.
\end{thmB}

This is Theorem \ref{thm:jw} and Theorem \ref{thm:annular} below, the latter being the spectral version, for the tracial states with $\psi(u_{ij})=\frac{1}{N}$, of the annular decomposition of Jones \cite{jo2}. In practice, if $c_{k-1}=Cat_{k-1}$, then $c_k$ is $Cat_k$ plus the multiplicity of the eigenvalue 1 of $\Psi_k$, and this can be decided by a probabilistic count on $\Lambda_k$, without storing any fixed vectors. We obtain in this way the following values:

\medskip

\begin{center}
\begin{tabular}{llc}
\toprule
order & matrices & $c_k=Cat_k$, numerically, up to\\
\midrule
6 & $T_6$, $BF_6$, $H_6^q$, $BN_6^q$, $K_6$ & $k=9$\\
7 & $P_7^q$, $R_7^q$ & $k=8$\\
8 & $F_{8,2}$, $F_{8,4}$, $K_{8,2}$, $M_{8,1}$, $I_{8,12}$, $K_{8,30}$, $L_{8,30}$ & $k=7$\\
9, 10 & $H_9$, $BN_9$, $BN_{10}$ & $k=7$\\
11, 12 & $N_{11}$, $N_{11}'$, Paley $P_{12}$ & $k=6$\\
20, 24 & Paley $P_{20}^1$, $P_{20}^2$, $P_{24}^1$ & $k=5$\\
20, 24 & $\mathrm{had}_{20.1}$, the 45 classes of order 24 with $c_2=2$ & $k=4$\\
32 to 48 & Paley $P_{32}^1$, $P_{36}^2$, $P_{44}^1$, $P_{48}^1$ & $k=4$\\
52 to 108 & Paley $P_{52}^2$, $P_{100}^2$, $P_{104}^1$, $P_{108}^1$, $P_{108}^2$ & $k=3$\\
28 & all the 487 classes, including $P_{28}^1$, $P_{28}^2$ & $k=3$\\
6 to 8 & 13,600 special points of the families & $k=4$\\
\bottomrule
\end{tabular}
\end{center}

\medskip

\noindent\textbf{Real matrices.} The real Hadamard matrices of orders 12 to 108 with $c_2=2$ that we tested, including the Paley matrices up to order 108, are all Temperley--Lieb as far as computed, while at order 16 all the 5 classes are of Di\c t\u a type, with sequences $c_k$ decreasing from one class to the next. A first step towards wildness is infinite depth, and we have here the following general principle, and its application to the Paley I matrices:

\begin{thmC}
If $G_H$ is finite, there is a monic polynomial $p\in\mathbb Z[x]$ such that $p(NT_k)=0$ for any $k$, so an eigenvalue of $NT_k$ which is not an algebraic integer forces infinite depth. For the matrices of order $2(q+1)$ built from symmetric conference matrices, including all the Paley II matrices, an explicit eigenvector of $NT_2$, with eigenvalue $1+(2q+1)(q-3)^2/(q+1)^2$, gives infinite depth for any $q\geq5$. For the Paley I matrix $P^1_{q+1}$, the matrices $X_\psi(x,y)=\psi(y-x)$ on $\mathbb F_q\times\mathbb F_q$, with $\psi$ a nontrivial character of $\mathbb F_q^*$ of odd order, are eigenvectors of $T_2$, with eigenvalues given by character sums, and this gives infinite depth for all the prime powers $q\equiv3$ mod 4 with $11\leq q<10^5$.
\end{thmC}

This is Theorem \ref{thm:integral}, Theorem \ref{thm:conference} and Theorem \ref{thm:paley} below. The first assertion recovers the infinite depth results of Montgomery for Paley II and $P_7$ \cite{mo2,mo3}, whose eigenvalues, such as $16q/(q+1)^2$ for Paley II, appear as eigenvalues of $NT_2$, and has as well a version for the traces, the numbers $Tr((NT_k)^r)$, which are partition functions on tori, having to be algebraic integers. The second assertion gives a new proof for Paley II, which uses only the conference matrix axioms, and so applies as well to the non-Paley symmetric conference matrices. For Paley II we also find, besides Montgomery's eigenvalue, an eigenvalue $\alpha=32((a^2-b^2)^2+1)$, where $q=a^2+b^2$ with $a$ odd, when $q$ is prime, coming from the value at 1 of Greene's hypergeometric function ${}_3F_2$ over $\mathbb F_q$ \cite{gre,ono}, and for Paley I the representation theory of $PSL_2(\mathbb F_q)$ produces, when $q\equiv3$ mod 8, a rational eigenvalue which numerically equals $(U-2)^2+16q$, with $U$ coming from a curve with complex multiplication by $\mathbb Z[\sqrt{-2}]$. For Paley I, the eigenvalues involve character sums attached to genus 1 curves over $\mathbb F_q$, which have no closed form, and this is why the second assertion is stated up to a bound, with exact integer certificates coming from Galois orbits of characters. The cases $q=3,7$, excluded, correspond to the Walsh matrices $W_4,W_8$, whose quantum groups are finite.

\medskip

\noindent\textbf{Motivations.} The original motivation is statistical mechanics in the sense of Jones \cite{jon,jo2}: the $c_k$ are the partition function side, and families of matrices with $N\to\infty$, such as the Paley ones, are the first place where thermodynamic limits can appear. On the other hand, the data suggests a general relation between the algebraic invariants of $H$, those of $G$, the geometric ones, such as the defect \cite{ba1,tz2}, and the analytic ones, such as the glow \cite{banicaHM}. Let us mention as well that there are other notions of quantum symmetries for the Hadamard matrices, coming from the associated Hadamard graphs, which lead to a different picture, with all the real Hadamard matrices of a given order being quantum isomorphic, see Gromada \cite{gro}, and Brannan, Gromada, Hern\'andez Palomares and Priebe \cite{bghp} for the complex case.

\medskip

\noindent\textbf{Computations.} Many results of this paper are of numerical nature. The values of $c_k$, and the norms, are floating point results, the ones coming from probabilistic counts being marked with a star, unless confirmed by another method, and only the bounds $c_k\geq Cat_k$, the closed formulae and the integer computations of section \ref{sec:infdepth} are exact, see section \ref{sec:conv}. All the programs used, written in Python, are available with the paper, together with their output files. The heaviest computations, at the highest levels, were done on 2 laptops, by 18 programs named after felines, from \texttt{lynx.py} to \texttt{sandcat.py}, and fighters, from \texttt{fedor.py} to \texttt{sefo.py}, which are listed in section \ref{sec:conv}, the other programs being mentioned in the text, at the place where they are used.

\medskip

\noindent\textbf{Organization.} Sections 1-2 contain preliminaries, on the complex Hadamard matrices, the quantum permutation groups, the Temperley--Lieb algebra, and the Hopf algebra tools that we need, written for readers coming from combinatorics, as well as from quantum algebra. Sections 3-4 deal with the deformed Fourier matrices, numerically and then algebraically, and section 5 deals with the Di\c t\u a families $F_{8,5}$ and the free wreath products. Section 6 deals with the real matrices, and sections 7-8 with the other wild matrices, of orders 6 and 7, and of orders 8 to 11. Finally, section 9 discusses the Jones--Wenzl and annular decompositions.

\section{Complex Hadamard matrices}

A complex Hadamard matrix is a square matrix $H\in M_N(\mathbb C)$ whose entries are on the unit circle, $|H_{ij}|=1$, and whose rows are pairwise orthogonal, with respect to the scalar product $\langle x,y\rangle=\sum_ix_i\bar{y}_i$ of $\mathbb C^N$. Equivalently, $H\in M_N(\T)$ and $H/\sqrt N$ is unitary. Two such matrices are called equivalent, $H\sim K$, if one can pass from one to the other by permuting the rows and columns, and by multiplying the rows and columns by numbers in $\T$. Any matrix is equivalent to a dephased one, having only 1 entries on the first row and the first column. See \cite{banicaHM,tz1}.

\medskip

\noindent\textbf{Fourier matrices.} The basic example is the Fourier matrix, $(F_N)_{ij}=w^{ij}$ with $w=e^{2\pi i/N}$, with indices $i,j\in\{0,\ldots,N-1\}$. More generally, for a finite abelian group $G=\mathbb Z_{N_1}\times\ldots\times\mathbb Z_{N_r}$ we have the generalized Fourier matrix $F_G=F_{N_1}\otimes\ldots\otimes F_{N_r}$, and we have $F_G\sim F_{G'}$ iff $G\simeq G'$. At $N=2,3,5$ the Fourier matrix is the only complex Hadamard matrix, up to equivalence \cite{ha1}, while at $N\geq4$ there are many other examples, which are the object of this paper. Already at $N=6$ the classification is a well-known open problem, see \cite{cwt,kar,sz4}.

\medskip

\noindent\textbf{Real and Butson matrices.} The real Hadamard matrices, with $\pm1$ entries, exist only at $N=1,2$ and at $N\in4\mathbb N$, and the Hadamard conjecture states that they exist at all $N\in4\mathbb N$. The basic examples are the Walsh matrices $W_{2^n}=F_2^{\otimes n}$, which are Fourier matrices, and the Paley and Williamson matrices, constructed in section \ref{sec:real} below. More generally, the Butson matrices are the complex Hadamard matrices whose entries are roots of unity of a given order.

\medskip

\noindent\textbf{Di\c t\u a deformations.} Given Hadamard matrices $H\in M_M(\T)$ and $K\in M_N(\T)$, and a parameter matrix $Q\in M_{M\times N}(\T)$, the Di\c t\u a deformation of $H\otimes K$ is the matrix \cite{dit}
$$(H\otimes_QK)_{ia,jb}=Q_{ib}H_{ij}K_{ab}$$
which is Hadamard, for any choice of $Q$. For $Q$ the all-one matrix we recover $H\otimes K$. The basic case is that of the deformed Fourier matrices $F_M\otimes_QF_N$, with for instance all the $4\times4$ complex Hadamard matrices being of the form $F_2\otimes_QF_2$, up to equivalence. More generally, and this is Di\c t\u a's original construction \cite{dit}, the inner matrix can depend on the outer row: given $K_1,\ldots,K_M\in M_N(\T)$, the matrix
$$(H\otimes_Q(K_1,\ldots,K_M))_{ia,jb}=Q_{ib}H_{ij}(K_i)_{ab}$$
is Hadamard too, because the scalar product of the rows $ia$ and $kc$ is
$$\sum_bQ_{ib}\bar{Q}_{kb}(K_i)_{ab}\overline{(K_k)}_{cb}\sum_jH_{ij}\bar{H}_{kj}=MN\delta_{ik}\delta_{ac}$$ We say that a matrix is of Di\c t\u a type when it is equivalent to some $H\otimes_Q(K_1,\ldots,K_M)$ with $M,N\geq2$. In such a matrix, the quotient of 2 rows $ia,ic$ with the same outer index depends only on the inner column index $b$, and this is what our block structure searches test, as a first step, a candidate structure being then subject to the condition that the outer blocks, obtained by restricting one row of each group to one class of columns, are all equivalent to $H$. Note that in Di\c t\u a's original array the inner matrix is attached to the block column, and that our row convention is its transpose. Since transposition matters, a matrix can be of Di\c t\u a type while its transpose is not. The general form is genuinely more general: at order 8, the matrices $F_{8,5b}$ of section \ref{sec:f85} are of the form $F_2\otimes(K,L)$, with inner matrices $K,L$ which are in general not equivalent, and in this case no single inner matrix can be used, as explained in section \ref{sec:nondita}, where an example with $c_2=3$ which is not of Di\c t\u a type, even in this general sense, is discussed too. We will see many more examples in what follows, of Di\c t\u a type or not.

\medskip

\noindent\textbf{Defect, regularity.} The defect of $H$, introduced by Tadej and \.Zyczkowski \cite{tz2}, is the number
$$d(H)=\dim\left\{A\in M_N(\mathbb R)\,\Big|\,\sum_kH_{ik}\bar{H}_{jk}(A_{ik}-A_{jk})=0,\ \forall i,j\right\}-(2N-1)$$
which is an upper bound for the dimension of the smooth families of dephased matrices passing through $H$. When $d(H)=0$, the matrix $H$ is isolated, as is for instance the case for the Fourier matrix $F_p$, with $p$ prime. Also, following \cite{gt}, a matrix is called regular when the scalar products between its rows decompose as sums of cycles, a cycle being a sum of the form $q(1+\zeta+\ldots+\zeta^{p-1})$, with $\zeta$ a primitive root of unity of prime order $p$, and $q\in\T$.

\medskip

\noindent\textbf{Commuting squares.} The interest in the complex Hadamard matrices from the operator algebra viewpoint comes from Popa's observation \cite{pop}, that up to conjugation by a unitary, the pairs of orthogonal maximal abelian subalgebras of $M_N(\mathbb C)$ are the pairs $(\Delta,H\Delta H^*)$, with $\Delta\subset M_N(\mathbb C)$ being the diagonal matrices, and with $H$ Hadamard \cite{ha1}. Indeed, for $D,E\in\Delta$ and $U\in U_N$ we have $tr(DUEU^*)=\frac{1}{N}\sum_{ij}D_{ii}E_{jj}|U_{ij}|^2$, so the orthogonality $tr(DUEU^*)=tr(D)tr(E)$ amounts to $|U_{ij}|=1/\sqrt N$. Moreover, the diagram
$$\begin{matrix}\Delta&\subset&M_N(\mathbb C)\\ \cup&&\cup\\ \mathbb C&\subset&H\Delta H^*\end{matrix}$$
is a commuting square in the sense of subfactor theory: the conditional expectations onto $\Delta$ and $H\Delta H^*$ commute, and their product is the expectation onto $\mathbb C$. By iterating the basic construction, any such commuting square produces a subfactor $A_0\subset A_1$ of the hyperfinite ${\rm II}_1$ factor $R$, of index $N$ in the sense of \cite{jo1}, whose planar algebra $P=(P_k)$, with $P_k=A_0'\cap A_k$, can be computed from the commuting square, by Ocneanu's compactness theorem \cite{jsu,ocn}. For the Hadamard commuting square, this computation was done by Jones. In order to formulate his result, we view the elements $T\in P_k$ as vectors $T\in(\mathbb C^N)^{\otimes k}$, that is, as linear maps $\mathbb C\to(\mathbb C^N)^{\otimes k}$, in their fixed tensor realization, explained in section \ref{sec:qg} below, and we set $T^\circ=id\otimes T\otimes id$. Also, we use the normalized Gram matrix of the row quotients, and its products along the rows of indices, given by
$$G_{ia}^{jb}=\frac{1}{N}\sum_kH_{ik}\bar{H}_{jk}\bar{H}_{ak}H_{bk}\quad,\quad G^k_{i_1\ldots i_k,j_1\ldots j_k}=G_{i_ki_{k-1}}^{j_kj_{k-1}}\ldots G_{i_2i_1}^{j_2j_1}$$
With these conventions, the result of Jones is as follows \cite{jo2,jon}:

\begin{theorem}\label{thm:jones}
The planar algebra $P=(P_k)$ of the subfactor associated to a complex Hadamard matrix $H\in M_N(\mathbb C)$ is given by
$$T\in P_k\iff T^\circ G^2=G^{k+2}T^\circ$$
with $T^\circ$ and $G$ being as above.
\end{theorem}

\begin{proof}
This follows from Ocneanu compactness, by a direct computation. See \cite{jo2,jon}, and also \cite[Theorem 14.23]{banicaHM}.
\end{proof}

For the readers coming from combinatorics, the only thing to be retained from the above is that the subfactor associated to $H$ produces a sequence of finite dimensional algebras $P_k\subset M_{N^k}(\mathbb C)$, given by explicit linear conditions in terms of $H$, whose dimensions $c_k=\dim P_k$ are the invariants studied in this paper. These algebras are subject to strong structural constraints, coming from the planar operations, and the theory of quantum permutation groups, explained in the next section, is the most convenient way of organizing them. The normalization $1/N$ in $G$ matters: for $H=F_N$ we have $G_{ia}^{jb}=\delta_{i-j,a-b}$, while without it, the 2 sides of the equation would differ by a factor $N^k$. For the readers coming from subfactors, with $N=A_0$ and $M=A_1$ we have $N'\cap M_{k-1}=A_0'\cap A_k=P_k$, and $c_k$ is its dimension, as explained in section \ref{sec:qg} below.

\section{Quantum permutation groups}

\subsection{Quantum permutations}\label{sec:qperm}

\noindent\textbf{Magic matrices.} Consider first the symmetric group $S_N$, viewed as a group of permutation matrices. The coordinate functions $u_{ij}:S_N\to\{0,1\}$, given by $u_{ij}(\sigma)=\delta_{i\sigma(j)}$, are commuting projections in the algebra $C(S_N)$, and since each $\sigma$ is a bijection, they sum up to 1 on each row and each column of the matrix $u=(u_{ij})$. A square matrix with entries in a $C^*$-algebra which are projections, summing up to 1 on each row and each column, is called magic, and by Gelfand duality $C(S_N)$ is the universal commutative $C^*$-algebra generated by the entries of an $N\times N$ magic matrix. The group structure of $S_N$ is encoded by the maps $\Delta(u_{ij})=\sum_ku_{ik}\otimes u_{kj}$, $\varepsilon(u_{ij})=\delta_{ij}$, $S(u_{ij})=u_{ji}$, which are the transposes of the multiplication, of the unit, and of the inversion. Following Wang \cite{wa1}, $C(S_N^+)$ is obtained by removing the commutativity: it is the universal $C^*$-algebra generated by the entries of an $N\times N$ magic matrix, with comultiplication, counit and antipode given by the same formulae. This is a Woronowicz algebra in the sense of \cite{wo1}, and $S_N^+$ is called the quantum permutation group.

\medskip

\noindent\textbf{Small values of $N$.} The entries of a magic matrix of size $N\leq3$ commute, so that $S_N^+=S_N$ for $N\leq3$ \cite{wa1}. At $N\geq4$ this is no longer true, because for any 2 projections $p,q$ in a $C^*$-algebra, the following matrix is magic:
$$\begin{pmatrix}p&1-p&0&0\\1-p&p&0&0\\0&0&q&1-q\\0&0&1-q&q\end{pmatrix}$$
Since the universal $C^*$-algebra generated by 2 projections is $C^*(\mathbb Z_2*\mathbb Z_2)$, which is noncommutative, and infinite dimensional, so is $C(S_N^+)$, for any $N\geq4$. In fact, $S_4^+$ is a twist of $SO_3$, in the sense explained below \cite{bb2}.

\medskip

\noindent\textbf{Quantum subgroups.} A closed quantum subgroup $G\subset S_N^+$ is given by a quotient $C(S_N^+)\to C(G)$ by a Hopf ideal, and the image of $u$, still denoted $u$, is called the fundamental corepresentation of $G$. The basic examples are the closed subgroups of $S_N$, the finite quantum groups, those for which $C(G)$ is finite dimensional, and the group duals: if $\Gamma$ is a discrete group generated by elements $g_1,\ldots,g_r$ of orders $N_1,\ldots,N_r$, with $N_1+\ldots+N_r=N$, then $C^*(\Gamma)$ is a quotient of $C(S_N^+)$, with $u$ block diagonal, the $s$-th block being obtained from $C^*(\langle g_s\rangle)\simeq C(\mathbb Z_{N_s})$ by Fourier transform \cite{bi1}. Given 2 quantum subgroups $G,K\subset S_N^+$, we can form their intersection $G\cap K$, and the quantum group $\langle G,K\rangle$ that they generate, which are the analogues of the classical constructions.

\medskip

\noindent\textbf{Fixed vectors.} The tensor powers of $u$ are the matrices $u^{\otimes k}=(u_{i_1j_1}\ldots u_{i_kj_k})$, and their fixed spaces are given by
$$Fix(u^{\otimes k})=\left\{\xi\in(\mathbb C^N)^{\otimes k}\,\Big|\,u^{\otimes k}(\xi\otimes1)=\xi\otimes1\right\}$$
More generally, we have the intertwiner spaces $Hom(u^{\otimes k},u^{\otimes l})$. By Woronowicz's results \cite{wo1}, $C(G)$ has a unique Haar state $\int_G$, and the Peter--Weyl theory gives
$$c_k=\dim Fix(u^{\otimes k})=\int_G\chi^k\quad,\quad\chi=\sum_iu_{ii}$$
where $\chi$ is the main character, satisfying $0\leq\chi\leq N$. For a classical group $G\subset S_N$, the function $\chi$ counts the fixed points, and $c_k$ is the number of orbits of $G$ on $\{1,\ldots,N\}^k$.

\medskip

\noindent\textbf{Tannakian duality.} The quantum group $G\subset S_N^+$ is uniquely determined by its intertwiner spaces $Hom(u^{\otimes k},u^{\otimes l})$, which form a tensor $C^*$-category \cite{wo2}, and the correspondence $G\to Hom$ is decreasing, the smaller quantum groups having more intertwiners. In particular, the intertwiner spaces of $\langle G,K\rangle$ are the intersections of those of $G$ and $K$. By Frobenius duality, everything is determined by the fixed spaces $Fix(u^{\otimes k})$, so these spaces, or rather their dimensions $c_k$, are the basic invariants of $G$.

\subsection{Partitions, Temperley--Lieb}\label{sec:tl}

\noindent\textbf{Partitions.} For a partition $\pi$ of $\{1,\ldots,k\}$, consider the vector
$$\xi_\pi=\sum_{i_1\ldots i_k}\delta_\pi(i)\,e_{i_1}\otimes\ldots\otimes e_{i_k}$$
where $\delta_\pi(i)=1$ if the indices are equal inside each block of $\pi$, and $\delta_\pi(i)=0$ otherwise. For instance at $k=2$ the partitions $\{1,2\}$ and $\{1\}\{2\}$ give the vectors $\sum_ae_a\otimes e_a$ and $\sum_{ab}e_a\otimes e_b$. The fixed spaces of $S_N$ are spanned by all the vectors $\xi_\pi$, which are linearly independent for $N\geq k$, and this gives the Bell numbers $c_k=1,2,5,15,52,\ldots$ For $S_N^+$, the fixed spaces are spanned by the vectors $\xi_\pi$ with $\pi$ noncrossing, meaning that there are no $a<b<c<d$ with $a,c$ in one block and $b,d$ in another block, and these vectors are linearly independent for $N\geq4$ \cite{ba1x,bsp}. Since the number of noncrossing partitions of $\{1,\ldots,k\}$ is the Catalan number $Cat_k$, this gives $c_k=Cat_k$. At $k=4$, the only crossing partition $\{1,3\}\{2,4\}$ gives the vector $\sum_{ab}e_a\otimes e_b\otimes e_a\otimes e_b$, which is fixed by $S_N$, and not by $S_N^+$.

\medskip

\noindent\textbf{Consequences.} For any quantum group $G\subset S_N^+$ with $N\geq4$ the $Cat_k$ vectors $\xi_\pi$ with $\pi\in NC(k)$ are fixed, so $c_k\geq Cat_k$, and by Tannakian duality we have equality for all $k$ iff $G=S_N^+$. In the language of free probability \cite{nsp}, $c_k=Cat_k$ for all $k$ means that $\chi$ follows the free Poisson law $\pi_1$, of density $\frac{1}{2\pi}\sqrt{4x^{-1}-1}$ on $[0,4]$. More generally, the free Poisson law $\pi_t$ has moments $\sum_{\pi\in NC(k)}t^{|\pi|}$, and the free cumulants $\kappa_n$ of a law are defined by the moment-cumulant formula $c_k=\sum_{\pi\in NC(k)}\prod_{\beta\in\pi}\kappa_{|\beta|}$, so that $\pi_t$ has all its free cumulants equal to $t$. The free convolution $\boxplus$ is additive on the free cumulants, and corresponds to the sum of free variables. These are the free analogues of the Poisson laws and of the cumulants.

\medskip

\noindent\textbf{Temperley--Lieb.} The intertwiners of $S_N^+$ are best described via the Temperley--Lieb algebra $TL_n(\delta)$, spanned by the diagrams of noncrossing pairings between 2 rows of $n$ points, with multiplication given by stacking the diagrams, the closed loops being replaced by $\delta$ factors \cite{jo1,tle}. By splitting each point into 2 points, the noncrossing partitions of $k$ points correspond to the noncrossing pairings of $2k$ points, and this fattening procedure gives an isomorphism $End(u^{\otimes k})\simeq TL_{2k}(\sqrt N)$ \cite{ba1x}. The Jones--Wenzl projection $f_n\in TL_n(\delta)$, for $\delta\geq2$, is the unique nonzero idempotent which is killed by all the diagrams containing a cap. It is given by $f_1=1$ and by Wenzl's recurrence \cite{wen}
$$f_{n+1}=f_n\otimes1-\frac{P_{n-1}(\delta)}{P_n(\delta)}(f_n\otimes1)E_n(f_n\otimes1)$$
where $E_n$ is the diagram with a cap and a cup at the positions $n,n+1$, and where $P_0=1$, $P_1=\delta$, $P_{n+1}=\delta P_n-P_{n-1}$ are Chebyshev polynomials, $P_n(\delta)$ being the trace of $f_n$.

\medskip

\noindent\textbf{Representations of $S_N^+$.} For $N\geq4$, the projection $p_s\in End(u^{\otimes s})$ corresponding to $f_{2s}\in TL_{2s}(\sqrt N)$ is minimal, its range is an irreducible corepresentation $r_s$, of dimension $d_s=P_{2s}(\sqrt N)$, and the corepresentations $r_0=1,r_1,r_2,\ldots$ obtained in this way are all the irreducible corepresentations of $S_N^+$, with fusion rules $r_s\otimes r_t=r_{|s-t|}+\ldots+r_{s+t}$, as for $SO_3$ \cite{ba1x}. Concretely, the range of $p_s$ consists of the vectors which are orthogonal to the images of all the intertwiners coming from lower levels, and so is the intersection of the kernels of the lowering maps, namely the merging of 2 adjacent legs, $e_a\otimes e_b\to\delta_{ab}e_a$, and the summation over a leg, $e_a\to1$. This description will be our computational definition of the Jones--Wenzl pieces, in section \ref{sec:jw}.

\medskip

\noindent\textbf{Planar algebras.} The vectors of $(\mathbb C^N)^{\otimes k}$ can be viewed as $k$-boxes, which can be composed along planar tangles, by contracting the indices along the strings, and this is the spin planar algebra of Jones \cite{jon}. For any $G\subset S_N^+$, the spaces $Fix(u^{\otimes k})$ form a planar subalgebra, and conversely, any planar subalgebra of the spin planar algebra comes from a unique quantum group $G\subset S_N^+$ \cite{ba2}. The tangles drawn in an annulus, generated by the rotation of the legs and by the Temperley--Lieb operations, form the annular category, and any planar algebra decomposes as a sum of irreducible annular modules, determined by their lowest weight vectors, those killed by all the annular operations decreasing the number of legs \cite{jo2}. This decomposition, extended to arbitrary tracial states, will be the key to the spectral computations of section \ref{sec:jw}.

\subsection{Hopf algebra tools}\label{sec:hopf}

\noindent\textbf{Hopf images.} Given a representation $\pi:C(S_N^+)\to B$, there is a largest Hopf ideal contained in $\ker\pi$, and the corresponding quantum group $G\subset S_N^+$ is called the Hopf image of $\pi$ \cite{bb1}. This is the smallest quantum subgroup $G\subset S_N^+$ such that $\pi$ factorizes through $C(G)$. Classically, if the elements $\pi(u_{ij})$ commute, they generate an algebra $C(Z)$, with $Z$ a compact space, each point of $Z$ produces a permutation, and the Hopf image is the closed subgroup of $S_N$ generated by these permutations. In general, the fixed spaces of the Hopf image are computed by Tannakian duality, in terms of the matrices $\pi(u_{ij})$.

\medskip

\noindent\textbf{Crossed coproducts.} Let $X$ be a finite abelian group, written additively, acting by automorphisms $r\to x\cdot r$ on a discrete group $\Gamma$. The crossed coproduct $C^*(\Gamma)\rtimes C(X)$ is the algebra $C^*(\Gamma)\otimes C(X)$, with comultiplication
$$\Delta(r\otimes\delta_x)=\sum_{y\in X}(r\otimes\delta_y)\otimes((-y)\cdot r\otimes\delta_{x-y})$$
for $r\in\Gamma$ and $x\in X$. This is a Woronowicz algebra, fitting into an exact sequence $C(X)\to C^*(\Gamma)\rtimes C(X)\to C^*(\Gamma)$, so that the corresponding quantum group is an extension of the group dual $\widehat{\Gamma}$ by $X$ \cite{ade,schn}. Its Haar state is given by $r\otimes\delta_x\to\delta_{r1}/|X|$. These quantum groups appear for the deformed Fourier matrices, with $\Gamma$ being a semidirect product $\mathbb Z^n\rtimes Y$, see section \ref{sec:cumul}.

\medskip

\noindent\textbf{Quotients.} The Hopf algebra quotients of $C^*(\Gamma)\rtimes C(X)$ which are injective on $C(X)$ were described by Bichon \cite{bic}. Such a quotient is given by a quotient datum $(N,\Phi)$, consisting of a normal subgroup $N\subset\Gamma$, which is stable under $X$, and of a morphism $\Phi:N\to C(X,\T)$, subject to compatibility conditions, and the quotient is obtained by identifying each $r\in N$ with the function $\Phi(r)$, namely
$$C^*(\Gamma/N)\rtimes_\Phi C(X)=C^*(\Gamma)\rtimes C(X)\big/\left<r\otimes1=1\otimes\Phi(r)\,\big|\,r\in N\right>$$
When $\Gamma/N$ is finite, the quotient is finite dimensional, of dimension $|\Gamma/N|\cdot|X|$. Intuitively, in a given model, $N$ consists of the elements of $\Gamma$ which act as functions on $X$, and the main result of section \ref{sec:cumul} is the computation of this pair, for the deformed Fourier matrices.

\medskip

\noindent\textbf{Free wreath products.} Given quantum subgroups $G\subset S_N^+$ and $G'\subset S_M^+$, the free wreath product $C(G\wr_*G')$ is generated by $M$ free copies $u^{(1)},\ldots,u^{(M)}$ of the fundamental corepresentation of $C(G)$, and by the fundamental corepresentation $v$ of $C(G')$, with the relations $[u^{(i)}_{ab},v_{ij}]=0$, the fundamental corepresentation being $w_{ia,jb}=u^{(i)}_{ab}v_{ij}$ \cite{bi0}. For instance $S_n^+\wr_*S_m^+$ is the quantum automorphism group of the graph $mK_n$ formed by $m$ disjoint copies of the complete graph $K_n$ \cite{bb4}, the hyperoctahedral quantum group is $H_N^+=\mathbb Z_2\wr_*S_N^+$ \cite{bbc}, and $H_2^+\wr_*\mathbb Z_2$ is the quantum automorphism group of the binary tree with 8 leaves \cite{dkr,dfk}. When the factors are group duals or abelian groups, the free wreath products are crossed coproducts as above, see section \ref{sec:bb}.

\medskip

\noindent\textbf{Free products, twists.} Given $G\subset S_N^+$ and $G'\subset S_M^+$, the dual free product $G\,\hat{*}\,G'\subset S_{N+M}^+$ corresponds to the free product $C(G)*C(G')$, with fundamental corepresentation $u\oplus u'$, and its Haar state is the free product of the Haar states \cite{wa0}. Thus its main character is a sum of 2 free variables, and its law is $\mu\boxplus\mu'$. As for the twists, $S_4^+=SO_3^{-1}$ means that $C(S_4^+)$ is obtained from $C(SO_3)$ by twisting its multiplication by a 2-cocycle on the Klein group $V\subset SO_3$ of diagonal matrices \cite{bb2}, and similarly, $H_2^+=O_2^{-1}$ \cite{bbc}. A twisted quantum group has the same representation theory, at the level of fusion rules, as the original group.

\medskip

\noindent\textbf{Coamenability.} A quantum group $G\subset S_N^+$ is coamenable when its Haar state is faithful on the full $C^*$-algebra, and by the analogue of Kesten's criterion, this happens iff $N$ belongs to the support of the law $\mu$ of $\chi$, that is, iff $\lim_kc_k^{1/k}=N$ \cite{ba3x}. For instance the finite quantum groups, and the crossed coproducts with $\Gamma$ amenable, are coamenable, while $S_N^+$ with $N\geq5$ is not, because $\lim_kCat_k^{1/k}=4$.

\subsection{The quantum group of a matrix}\label{sec:qg}

\noindent\textbf{Definition.} Given a complex Hadamard matrix $H\in M_N(\T)$ with rows $H_1,\ldots,H_N$, the rank 1 projections
$$P_{ij}=Proj(H_i/H_j)$$
form a magic matrix, because $\langle H_i/H_j,H_i/H_k\rangle=N\delta_{jk}$ and $\langle H_i/H_j,H_k/H_j\rangle=N\delta_{ik}$. Thus we have a representation $\pi:C(S_N^+)\to M_N(\mathbb C)$, $u_{ij}\to P_{ij}$, and the quantum group $G=G_H\subset S_N^+$ associated to $H$ is its Hopf image, the minimal one producing a factorization $\pi:C(S_N^+)\to C(G)\to M_N(\mathbb C)$ \cite{bb1,bbs}. For the Fourier matrix $F_G$ we obtain the group $G$ itself, acting on itself, and for a tensor product $H'\otimes H''$ we obtain the product $G'\times G''$ \cite[Theorem 14.5]{banicaHM}.

\medskip

\noindent\textbf{An example.} For $H=F_2$, with rows $H_1=(1,1)$ and $H_2=(1,-1)$, we have $H_1/H_1=H_2/H_2=(1,1)$ and $H_1/H_2=H_2/H_1=(1,-1)$, so the magic matrix is
$$P=\begin{pmatrix}Q_+&Q_-\\ Q_-&Q_+\end{pmatrix}\quad,\quad Q_\pm=Proj\left(\begin{matrix}1\\ \pm1\end{matrix}\right)$$
Its entries commute, they generate $C(\mathbb Z_2)\simeq\mathbb C^2$, and $G=\mathbb Z_2$. The fixed vectors can be computed directly, via the matrix $T_k$ of section \ref{sec:conv}, with entries $tr(P_{i_1j_1}\ldots P_{i_kj_k})$. At $k=2$, in the basis $11,12,21,22$, we obtain
$$T_2=\frac{1}{2}\begin{pmatrix}1&0&0&1\\0&1&1&0\\0&1&1&0\\1&0&0&1\end{pmatrix}$$
which has eigenvalues $1,1,0,0$, and so $c_2=2$, the fixed vectors being $e_1\otimes e_1+e_2\otimes e_2$ and $e_1\otimes e_2+e_2\otimes e_1$. More generally, for $F_N$ we have $G=\mathbb Z_N$, and $c_k=N^{k-1}$. At the other extreme, the free answer is $c_k=Cat_k$, meaning $G=S_N^+$, and the general question is to locate $G_H$ between these 2 extremes.

\medskip

\noindent\textbf{Recovering the planar algebra.} By Tannakian duality \cite{wo2}, the Hopf image is described by the formula $Hom(u^{\otimes k},u^{\otimes l})=Hom(P^{\otimes k},P^{\otimes l})$, the spaces on the right being taken in a formal sense, and a computation with the vectors $H_i/H_j$ shows that $T\in Hom(u^{\otimes k},u^{\otimes l})$ iff $T^\circ G^{k+2}=G^{l+2}T^\circ$, with the notations of Theorem \ref{thm:jones} \cite[Theorem 14.9]{banicaHM}. By comparing with Theorem \ref{thm:jones}, we obtain that the planar algebra of the Hadamard subfactor is given by $P_k=Fix(u^{\otimes k})$ \cite{bbs}. Moreover, the subfactor itself is of the form $A^G\subset(\mathbb C^N\otimes A)^G$, with $A=R\rtimes\widehat{G}$ \cite{ba0}. The main character $\chi=\sum_iu_{ii}$ being a sum of $N$ projections, its law $\mu=\mu_H$ with respect to the Haar functional is a probability measure on $[0,N]$, with moments
$$c_k=\int_G\chi^k=\dim Fix(u^{\otimes k})=\dim(P_k)$$
by the Peter--Weyl theory of \cite{wo1}. Thus the Poincar\'e series of the subfactor $f(z)=\sum_kc_kz^k$ is the Stieltjes transform of $\mu$, namely $f(z)=\int_G\frac{1}{1-z\chi}$. The subfactor has finite depth precisely when $G$ is finite, and the main problem is that of computing $\mu$, and so the numbers $c_k$, and if possible $G$ itself, for the various complex Hadamard matrices.

\medskip

\noindent\textbf{Dictionary.} The literature uses several languages, as follows:

\medskip

\begin{center}
\begin{tabular}{ll}
\toprule
subfactors, commuting squares & quantum groups\\
\midrule
Hadamard matrix, biunitary, spin model, type II matrix & $H\in M_N(\T)$\\
relative commutant $N'\cap M_{k-1}$ & $Fix(P^{\otimes k})$, of dimension $c_k$\\
``only Temperley--Lieb'', principal graph $A_\infty$ & $c_k=Cat_k$, that is, $G=S_N^+$\\
finite depth & $G$ finite\\
connected components of the profile graph & orbitals, $c_2$\\
intermediate subfactor & $G$ preserves a partition of $\{1,\ldots,N\}$\\
defect, isolation, span condition & geometry of the moduli space\\
\bottomrule
\end{tabular}
\end{center}

\medskip

\noindent Here we use $c_k=\dim(N'\cap M_{k-1})$, normalized so that $F_N$ gives $c_2=N$. Beware of the index shifts in the literature: \cite{cn} call $N'\cap M_1$ the second relative commutant, while \cite{ben} call it the first. See \cite{jmn,nom,pop} for the various equivalent formulations, and \cite{ba1,mwe,szo,tz2} for the defect and the isolated matrices.

\medskip

\noindent\textbf{Known facts.} We will use the following results:
\begin{enumerate}
\item As explained in section \ref{sec:tl}, we have $c_k\geq Cat_k$ for $N\geq4$, with equality for all $k$ iff $G=S_N^+$. See \cite{bcm} for the related maximality questions for $S_N\subset S_N^+$. A single extra fixed vector, at some level, means that $G$ is smaller.

\item At level 2 \cite{cn,jmn,jon}, $c_2$ is the number of connected components of the graph on pairs $(a,c)$, with $(a,c)\sim(b,d)$ when $\langle H_a/H_b,H_c/H_d\rangle\neq0$, and is the dimension of the Nomura algebra of $H$. We have $c_2=N$ iff $H$ is a tensor product of Fourier matrices, and $c_2=2$ for real Hadamard matrices of order $\equiv4$ mod $8$ \cite{jmn}. The matrices of Di\c t\u a type have an intermediate subfactor, in the sense of \cite{bha,gj}, and so $c_2\geq3$ \cite{cn}. This holds as well in the general form $H\otimes_Q(K_1,\ldots,K_M)$ of section 1, the indicator of the partition of the rows according to the outer index being a fixed vector of $G_H$, as one checks directly.

\item For the Di\c t\u a deformations $F_M\otimes_QF_N$ with generic $Q$ \cite{bb3}, we have $C(G)=C^*(\Gamma_{X,Y})\rtimes C(X)$ with $\Gamma_{X,Y}=\mathbb Z^{(|X|-1)(|Y|-1)}\rtimes Y$, and an explicit formula for $c_k$, see section \ref{sec:bb}. See also \cite{bic,bur}. The case of non-generic $Q$ was open.

\item The angle operator $\Theta$ of \cite{jon,mo1} has as 1-eigenspace the standard invariant, and $\tau(\Theta^n)$ gives $c_n$. It is not clear whether this operator is the matrix $T_k$ below in disguise, but its eigenvalues found in \cite{mo2,mo3} appear as eigenvalues of $NT_2$, see section \ref{sec:infdepth}.
\end{enumerate}

\medskip

\noindent\textbf{Previous work.} The only computations of quantum invariants beyond level 2 that we are aware of are those of Camp and Nicoara \cite{cn,ni}:

\medskip

\begin{center}
\begin{tabular}{lccc}
\toprule
matrix & $c_2$ & $c_3$ & $c_4$\\
\midrule
$F_4^s$, $s=\pm1,\pm i$ (Fourier) & 4 & & \\
$F_4^s$, other $s\in\T$ & 3 & 10 & 35, or 36 if $s^8=1$ primitive\\
$F_6(a,b)$, generic Di\c t\u a & 4 & & \\
$F_8(a,b,c,d,z)$, generic Di\c t\u a & 4 & & \\
Beauchamp--Nicoara $BN_6^q$, Tao $T_6$, Haagerup & 2 & & \\
Petrescu $P_7^q$, all $q$ & 2 & 5 (at $q=1$) & \\
\bottomrule
\end{tabular}
\end{center}

\medskip

\noindent They conjecture that the principal graph of $F_4^s$ is $D_{2k}^{(1)}$ when $s$ is a primitive root of unity of order $2k$, and $D_\infty^{(1)}$ otherwise, and that $BN_6^q$ and $P_7^q$ are supertransitive. They also report that every non-Di\c t\u a, non-Fourier matrix they tested, for $N\leq11$, has $c_2=2$, and $c_3=5$ where computed. These are the first observations of the type summarized in section \ref{sec:nondita}. Montgomery \cite{mo2,mo3} proved that the Paley II matrices and $P_7$ give infinite depth subfactors, and conjectured that Paley II is $A_\infty$. For Paley I, see section \ref{sec:infdepth}. The rest of this literature \cite{fnw,hn1,hn2,ho,hw,nw1,nw2,nw4,nwo,sch,whi} deals with the moduli space and with the deformation theory, and the deformations of \cite{nwh,wor} will be discussed in section \ref{sec:dita}.

\subsection{Transposition}\label{sec:transp}

\noindent\textbf{Conjugation and transposition.} The 4 matrices $H,\bar{H},H^t,H^*$ are related as follows \cite{ba0,ba5}:

\begin{proposition}\label{prop:conj}
We have $G_{\bar{H}}=G_H^\sigma$, where $C(G^\sigma)$ is $C(G)$ with comultiplication $\Sigma\Delta$, with $\Sigma$ the flip. Thus $\mu_{\bar{H}}=\mu_H$ and $\mu_{H^*}=\mu_{H^t}$. Also, if $G_H$ is finite, then $G_{H^t}=\widehat{G}_H$.
\end{proposition}

\begin{proof}
The first assertion and the last one are from \cite[Proposition 2.2, Theorem 2.3]{ba5}. As for the laws, the main character of $G^\sigma$ is the same element $\chi=\sum_iu_{ii}$ of the same algebra, with the same Haar functional, and $H^*=\overline{H^t}$.
\end{proof}

Thus at most 2 computations are needed per matrix, for $H$ and $H^t$, and in fact 1 suffices, by Theorem \ref{thm:transp} below. The truncations of $\mu$, from \cite{ba5}, are defined as follows. With $\int_G^r=(tr\circ\rho)^{*r}$, where $\rho:C(G)\to M_N(\mathbb C)$ is the factorizing map, we have $\int_G=\lim_k\frac{1}{k}\sum_{r\leq k}\int_G^r$, and the law $\mu^r$ of $\chi$ with respect to $\int_G^r$ has moments $c_k^r=Tr(T_k^r)$, where
$$(T_k)_{i_1\ldots i_k,j_1\ldots j_k}=tr(P_{i_1j_1}\ldots P_{i_kj_k})$$
Since $T_k$ is positive, with spectrum in $[0,1]$, see section \ref{sec:conv}, we have in fact $c_k=\lim_{r\to\infty}Tr(T_k^r)$, decreasingly. We have the following results, from \cite{ba5}:

\begin{theorem}\label{thm:trunc}
The truncations $\mu^r=\mu^r_H$ have the following properties:
\begin{enumerate}
\item $\mu^r$ is the law of the Gram matrix $X^{(r)}=(\langle\xi_a,\xi_b\rangle)_{a,b\in\{1,\ldots,N\}^r}$ of the unit vectors $\xi_{a_1\ldots a_r}=N^{-r/2}\,\frac{H_{a_1}}{H_{a_2}}\otimes\ldots\otimes\frac{H_{a_r}}{H_{a_1}}$, with respect to the normalized trace.

\item We have the duality formula $\int_{G_H}^r(\chi/N)^p=\int_{G_{H^t}}^p(\chi/N)^r$, for any $p,r$.

\item For $H=F_M\otimes_QF_N$ we have $\mu_H=\mu_{H^t}$, for any $Q\in M_{M\times N}(\T)$.

\item We have $\mu_H(\{N\})=\mu_{H^t}(\{N\})$, and $\mu_H(\{N\})=1/|G_H|$ when $G_H$ is finite.
\end{enumerate}
\end{theorem}

\begin{proof}
These are \cite[Theorem 4.4, Theorem 5.2, Propositions 4.1, 4.5, 4.6]{ba5}.
\end{proof}

In all our computations, without any exception, we found $c_k(H)=c_k(H^t)$. In fact, this holds in general, and the Di\c t\u a assumption in Theorem \ref{thm:trunc} (3) can be removed:

\begin{theorem}\label{thm:transp}
For any complex Hadamard matrix $H\in M_N(\T)$ and any $k$ we have
$$T_k(H)=N^{-k-1}AA^*\quad,\quad T_k(H^t)=N^{-k-1}A^*A$$
where $A\in M_{N^k}(\mathbb C)$ is given by $A_{i,x}=\prod_s\bar{H}_{i_sx_s}H_{i_{s+1}x_s}$, with indices mod $k$. Thus $T_k(H)\simeq T_k(H^t)$, so $\mu^r_H=\mu^r_{H^t}$ for any $r$, and $\mu_H=\mu_{H^t}$.
\end{theorem}

\begin{proof}
With $\xi_{ij}=N^{-1/2}H_i/H_j$, cyclicity gives $(T_k)_{ij}=\frac{1}{N}\prod_s\langle\xi_{i_{s+1}j_{s+1}},\xi_{i_sj_s}\rangle$, and expanding each scalar product as
$$\langle\xi_{i_{s+1}j_{s+1}},\xi_{i_sj_s}\rangle=\frac{1}{N}\sum_{x_s}\bar{H}_{i_sx_s}H_{j_sx_s}H_{i_{s+1}x_s}\bar{H}_{j_{s+1}x_s}$$
gives the first formula. In $(A^*A)_{xy}=\sum_i\prod_sH_{i_sx_s}\bar{H}_{i_{s+1}x_s}\bar{H}_{i_sy_s}H_{i_{s+1}y_s}$ each $i_s$ appears in the factors $s-1$ and $s$, and summing over it gives $N$ times a scalar product of the vectors of $H^t$, which gives the second formula. Finally $AA^*\simeq A^*A$, and the truncated moments are the numbers $Tr(T_k^r)$.
\end{proof}

In particular $H$ is wild iff $H^t$ is, and the Jones--Wenzl pieces of section \ref{sec:jw} satisfy $\Phi_s(H)\simeq\Phi_s(H^t)$. On the fixed vectors, $v\to N^{-(k+1)/2}Av$ is an explicit isometry from $Fix(G_{H^t}^{\otimes k})$ to $Fix(G_H^{\otimes k})$, which might help with the question of the relation between $G_H$ and $G_{H^t}$ for $G_H$ infinite, where the dual is no longer compact, and where Proposition \ref{prop:conj} should be replaced by something weaker, such as a monoidal equivalence. Also, $T_k(H)=\frac{1}{N}\sum_xProj(\zeta_x)$ is $\frac{1}{N}$ times the frame operator of the unit vectors $\zeta_x=N^{-k/2}\,\eta_{x_kx_1}\otimes\eta_{x_1x_2}\otimes\ldots\otimes\eta_{x_{k-1}x_k}$, with $\eta_{ab}=H^a/H^b$ the column ratios. Combined with Theorem \ref{thm:trunc} (2), the array $N^{-p}Tr(T_p^r)$ is symmetric in $p,r$: in vertex model terms, the partition function on the $p\times r$ torus is invariant under the rotation of the torus. The duality formula gives a useful reformulation of the problem:

\begin{theorem}\label{thm:gram}
We have $c_p(H)=N^p\mu^p_{H^t}(\{N\})$, that is, $c_p(H)$ is the multiplicity of the eigenvalue $N$ of the Gram matrix $X^{(p)}$ associated to $H^t$, and also, by Theorem \ref{thm:transp}, of the one associated to $H$.
\end{theorem}

\begin{proof}
In the duality formula, the left side is $c_p^r(H)/N^p$, which converges to $c_p(H)/N^p$ with $r\to\infty$. The right side is $\int_0^N(x/N)^rd\mu^p_{H^t}(x)$, which converges to $\mu^p_{H^t}(\{N\})$, because $\mu^p_{H^t}$ is supported on $[0,N]$. Finally, by Theorem \ref{thm:trunc} (1) the measure $\mu^p_{H^t}$ is the law of an $N^p\times N^p$ matrix with respect to the normalized trace, so $N^p\mu^p_{H^t}(\{N\})$ is the multiplicity of the eigenvalue $N$.
\end{proof}

We checked this formula numerically at $p=1,2,3$, for $\mathrm{had}_{16.3}$, $F_{8,5a}$, $F_{8,5b}$, $T_6$ and $F_4^s$, with the multiplicities $1,3,11$, then $1,4,19$, then $1,3,11$, then $1,2,5$, and then $1,3,10$, matching the values of $c_1,c_2,c_3$ computed below.

\section{Deformed Fourier matrices, numerics}\label{sec:numerics}

\subsection{Conventions and methods}\label{sec:conv}

\noindent\textbf{Notations.} We write $e(t)=e^{2\pi it}$, $w=e(1/3)$ and $i=e(1/4)$, and the indices run over $0,\ldots,N-1$, with $(F_N)_{ij}=e(ij/N)$. For the Di\c t\u a deformations we use the convention $(H\otimes_QK)_{ia,jb}=Q_{ib}H_{ij}K_{ab}$ of section 1, with $Q\in M_{M\times N}(\T)$ dephased, and with the double index $ia$ standing for $iN+a$, as in \cite[Theorem 5.12]{banicaHM}. A matrix in logarithmic form is given by $H'$, with $H_{ij}=\omega^{H'_{ij}}$ for a stated root of unity $\omega$, and the real matrices are displayed as sign patterns. We denote by $Cat_k=1,2,5,14,42,132,429,1430,4862$ the Catalan numbers, at $k=1,\ldots,9$.

\medskip

\noindent\textbf{Equivalences.} By Proposition \ref{prop:conj}, only $H$ and $H^t$ need to be computed. For each matrix, we first determine the equivalences between $H,H^t,\bar{H},H^*$, by placing them in the known parametrizations, and by an exact equivalence test, so that each equivalence class is computed once. The matrices themselves are built by the module \texttt{matrices.py}, which contains all the conventions of the present paper, and the equivalences are tested by \texttt{equiv.py}, by using, in the real case, McKay's reduction to a graph isomorphism problem.

\medskip

\noindent\textbf{Exact computation.} With $\xi_{ij}=(H_i/H_j)/\sqrt N$ the unit vector spanning $P_{ij}$, consider the matrix $T_k\in M_{N^k}(\mathbb C)$ given by
$$(T_k)_{i_1\ldots i_k,j_1\ldots j_k}=tr(P_{i_1j_1}\ldots P_{i_kj_k})=\frac{1}{N}\prod_{s=1}^k\langle\xi_{a_{s+1}},\xi_{a_s}\rangle$$
where $a_s=(i_s,j_s)$, with $a_{k+1}=a_1$, and with $tr$ the normalized trace. We have $T_k=N^{-k-1}VV^*$, with $V_{i,x}=\prod_sH_{i_sx_{s-1}}\bar{H}_{i_sx_s}$ for $i,x\in\{1,\ldots,N\}^k$, with $x_0=x_k$, so $T_k$ is positive. Also, $T_k$ is a contraction, being the image of the unitary $P^{\otimes k}\in M_{N^k}(M_N(\mathbb C))$ under the conditional expectation $id\otimes tr$. Thus the spectrum of $T_k$ is contained in $[0,1]$, and $c_k=\dim\ker(T_k-1)$ is the multiplicity of the eigenvalue 1. For $N^k\leq5000$ we compute this multiplicity directly, from the singular values of $T_k-1$, and we record the ratio between the smallest nonzero singular value and the largest zero one. This, and the methods below, are implemented in the module \texttt{ck.py}, on which all the other programs are built.

\medskip

\noindent\textbf{Seeded computation.} For larger $N^k$, we first generate candidate fixed vectors at level $k$ from those at lower levels, namely leg doublings, tensor products and cyclic rotations, and we verify each of them, $\|T_kv-v\|<10^{-8}$, after orthogonalization. This gives an orthonormal family $S$ of $s$ numerically fixed vectors, and so $c_k\geq s$, up to the numerical accuracy. The noncrossing seeds below are exactly fixed, but the other candidates are built from numerically computed vectors, and a small residual alone does not certify a fixed vector. The candidates are processed in blocks of 64, the independent directions being extracted from the Gram matrix of each block. In order to decide whether $S$ spans the whole fixed space, we use either a deflated Lanczos computation on $S^\perp$, which is the method used for the values at level 5 written without star, or the probabilistic count below. Note that $T_k$ is never formed: we only use matrix-vector products, which factor through $V$, at a cost of order $N^{k+2}$.

\medskip

\noindent\textbf{Probabilistic count.} On $S^\perp$ we take a random Gaussian block $X\in M_{n\times b}$, with $n=N^k$, and apply the Chebyshev filter $p(T_k)$, where $p(x)=T_{50}(2x/a-1)$ with $a=0.95$, which is bounded by 1 on $[0,a]$, and equals $\cosh(50\,\mathrm{arccosh}(2/a-1))\simeq4\cdot10^{9}$ at $x=1$. After Rayleigh--Ritz on the span of $p(T_k)X$, we count the Ritz values $>1-10^{-4}$, and check the corresponding vectors directly, with residuals $\|T_kv-v\|$ of order $10^{-8}$. We denote by $X^*$ the value $s+\#$ obtained in this way, which is $c_k$ with high confidence, but not with a deterministic error bound: a residual of order $10^{-8}$ does not distinguish the eigenvalue 1 from an eigenvalue $1-10^{-10}$, for instance. The failure modes are a random block nearly orthogonal to a fixed vector, which is excluded by taking $b$ well above the number of vectors found, restarting with $2b$ if this number exceeds $b-10$, and an eigenvalue in $(a,1)$, which would show up as an intermediate Ritz value. As a test, against known exact values, we recover $c_5=197$ for $F_{8,5b}$, $c_4=142$ for $\mathrm{had}_{16.2}$, and $c_4=46$ for $\mathrm{had}_{16.3}$.

\medskip

\noindent\textbf{Noncrossing seeds.} For any $G\subset S_N^+$, the $Cat_k$ vectors $\xi_\pi$ with $\pi\in NC(k)$ are fixed. For the matrices expected to be Temperley--Lieb, we use these vectors as seeds, and the random block then only checks whether something beyond Temperley--Lieb is fixed. This makes the computation of $c_7$ at $N=6$, in dimension $6^7=279936$, a matter of minutes.

\medskip

\noindent\textbf{Rotation sectors.} Since the trace is cyclic, $T_k$ commutes with the simultaneous cyclic rotation of the $k$ legs, $(T_k)_{\rho(i),\rho(j)}=(T_k)_{ij}$. Thus $T_k$ is block diagonal with respect to the $k$ eigenspaces of the rotation, of dimension about $N^k/k$ each, and $c_k$ is the sum of the multiplicities of the eigenvalue 1 in the blocks. Working in one block at a time, with one coordinate per rotation orbit, divides the memory by about $k$. This is how the values $c_6$ of section \ref{sec:f85} were obtained, and the decomposition is implemented in \texttt{sectors.py} and \texttt{lwsector.py}, the latter restricting moreover to the lowest weight vectors of section \ref{sec:jw}.

\medskip

\noindent\textbf{Numerical status, stars.} The values of $c_k$ in this paper are obtained in floating point arithmetic, and are numerical results, with the exception of the bounds $c_k\geq Cat_k$, which are exact, of the values given by closed formulae, as those of section \ref{sec:cumul}, and of the exact integer computations of section \ref{sec:infdepth}. Among the numerical values, a star marks those coming from the probabilistic count above, for the integers $c_k^*$ and for the Jones--Wenzl multiplicities $e_k^*$ of section \ref{sec:jw}, which are high confidence values. The values without star come from the full singular value decomposition of $T_k-1$, with the gap ratios indicated, or from a deflated Lanczos computation of the top eigenvalues of $T_k$ on the orthogonal complement of the fixed vectors found, or from a converged norm computation, via Theorem \ref{thm:jw}. These are more reliable, but are not proofs either: a converged Ritz value is an eigenvalue up to its residual, and a lower bound for the top of the spectrum, but a Krylov iteration gives no global upper bound, an eigenvector nearly orthogonal to the starting block being possible in principle. For the norms $\|\Phi_k\|$ of section \ref{sec:jw} the values are Ritz values, and so lower bounds, and a star means that the iteration was stopped with a residual above $10^{-4}$. At an algebraic point, a proof of $c_k=Cat_k$ could be obtained from an exact rank computation, such as a nonzero minor of $T_k-1$ of the right size modulo a suitable prime, but we have not done this.

\medskip

\noindent\textbf{Exact checks.} Some of the numerical results are complemented by exact computations, in integer or cyclotomic arithmetic: the integrality obstructions for the Paley I matrices, in section \ref{sec:paley1}, the value $c_2=2$ at the Haagerup point of the Karlsson family, in section \ref{sec:wild}, which gives $c_2^{gen}=2$ for this family, and the value $c_2=3$, together with the Di\c t\u a block structures, at a point of $F_{8,5b}$, of its transpose, and of $F_{8,5c}$, with entries roots of unity, in section \ref{sec:nondita}. These exact checks are done by the programs \texttt{trace\_exact.py}, using the GMP library for big integer products, \texttt{c2exact.py}, working with Gaussian integers, and \texttt{f85bpoint.py} and \texttt{f85cpoint.py}, working with exponents of roots of unity, and reduction modulo cyclotomic polynomials. On the other hand, all the other values of $c_k$, and all the norms, are numerical, as explained above, and turning them into proofs would require exact rank computations, or rigorous enclosures for the relevant spaces.

\medskip

\noindent\textbf{Positivity.} Since $\chi=\sum_iu_{ii}$ satisfies $0\leq\chi\leq N$, the sequence $(c_k)$, with $c_0=1$, is the moment sequence of a probability measure on $[0,N]$. Thus the Hankel matrices $(c_{i+j})$, $(c_{i+j+1})$ and $(Nc_{i+j}-c_{i+j+1})$ must be positive, which gives free lower bounds on $c_k$, in terms of $c_1,\ldots,c_{k-1}$.

\medskip

\noindent\textbf{Cost and checks.} On a laptop with 8GB of memory, typical times are: 3 minutes for $F_{8,5a}$ at $k=5$, 40 minutes for $\mathrm{had}_{16.4}$ at $k=5$, from 1 minute ($N=6$) to 2.5 hours ($N=10$) for the matrices of sections \ref{sec:wild} and \ref{sec:wild8} at $k=6$, 75 minutes for $P_{20}^1$ at $k=5$, and from 77 minutes ($F_{8,5b}$) to 4 hours ($F_{8,5a}$) at $k=6$, by rotation sectors. Every matrix below was checked numerically to be Hadamard, $|H_{ij}|=1$ and $HH^*=N1_N$ to precision $10^{-9}$, at 5 random points for the families, and the formulas used in the computations are the ones displayed here. The exact checks mentioned above use only integer arithmetic, with Python integers, or with the GMP library, and are independent of the floating point computations, the corresponding programs, and their output files, being included in the archive.
\medskip

\noindent\textbf{Programs.} All the programs, with their output files, are in the code archive accompanying the paper, described by a README file. The basic layer consists of the modules \texttt{matrices.py}, \texttt{ck.py}, \texttt{equiv.py} mentioned above, and of \texttt{dita\_test.py}, which searches for the Di\c t\u a block structures, as explained in section \ref{sec:f85}. The first sweeps through the basis of matrices were done by the programs \texttt{stage0.py} to \texttt{stage46.py}, \texttt{level5.py} and \texttt{sun.py}, and the more specialized programs are mentioned in the text, at the place where they are used. Most of these programs run in seconds or minutes, and all of them use only \texttt{numpy} and \texttt{scipy}, besides the GMP library mentioned above. The archive contains, for each of them, the version actually used, together with its output files, as described in its README file. The laptop programs below, on the other hand, run from several minutes to a full night each, the longest single jobs, at the highest levels, taking several hours, as indicated by the timings recorded in their log files, which are included in the archive, next to the programs themselves. The longest computations were done on 2 laptops with 8GB of memory, overnight or during the day, by programs named after felines and fighters, which write their results in log and json files, stop by themselves after a fixed time, and resume where they stopped when launched again. These programs are as follows:

\medskip

\begin{center}
\footnotesize
\begin{tabular}{lp{11.2cm}c}
\toprule
program & computation & section\\
\midrule
\texttt{ocelot.py} & $c_5$ for $F_{8,5a}$ and $\mathrm{had}_{16.4}$, $c_4$ for $\mathrm{had}_{16.1}$, $c_6$ for $F_{8,5b}$ & \ref{sec:f85}, \ref{sec:had16}\\
\texttt{hoost.py} & $c_6$ for $F_{8,5a},F_{8,5b},F_{8,5c}$, by rotation sectors & \ref{sec:f85}\\
\texttt{sandcat.py} & multiplicities $e_5,e_6$ for the families $F_{8,5}$, then converged norms $\|\Phi_k\|$ at orders 7 to 24 & \ref{sec:f85}, \ref{sec:jwpieces}\\
\texttt{margay.py} & independent check of $c_5=212$ for $\mathrm{had}_{16.3}$, with a new seed and Lanczos & \ref{sec:had16}\\
\texttt{lynx.py} & $c_1,\ldots,c_4$ for the 45 classes of order 24 with $c_2=2$, for $H$ and $H^t$ & \ref{sec:orders}\\
\texttt{fedor.py} & $c_1,c_2,c_3$ for the 487 classes of order 28, for $H$ and $H^t$, then $c_4$ & \ref{sec:orders}\\
\texttt{bobcat.py} & Jones--Wenzl counts for the Paley matrices of orders 20 to 36 & \ref{sec:realres}\\
\texttt{cormier.py} & Jones--Wenzl counts for the Paley matrices of orders 44 to 108 & \ref{sec:realres}\\
\texttt{bones.py} & glued matrices of order 10, with $c_2=2$, up to $e_6$, and $e_7$ for 1 of them & \ref{sec:nine}\\
\texttt{caracal.py} & $c_7$ for the 7 matrices of order 6, with noncrossing seeds & \ref{sec:level5}\\
\texttt{puma.py} & $c_6$ for the matrices of orders 6 to 10, and $c_5$ for $P_{20}^1$ & \ref{sec:level5}\\
\texttt{gsp.py} & $c_6$ for $I_{8,12},K_{8,30},L_{8,30},L_{8,30}^t,P_7^1,BN_9$, and $c_5$ at order 20 & \ref{sec:level5}\\
\texttt{crocop.py} & norms $\|\Phi_k\|$ at the low levels, for the wild matrices of orders 6 to 12 & \ref{sec:jwpieces}\\
\texttt{serval.py} & norms $\|\Phi_k\|$ at the next levels, up to level 8 at order 6 & \ref{sec:jwpieces}\\
\texttt{jaguar.py} & Jones--Wenzl counts $e_6,e_7,e_8$, for the wild matrices of orders 6 to 12 & \ref{sec:jwpieces}\\
\texttt{aerts.py} & $e_9$ for the 7 matrices of order 6, and $e_7$ for $H_9,BN_9$ & \ref{sec:jwpieces}\\
\texttt{lennox.py} & $e_8$ for the 3 matrices of order 7, $e_6$ for $P_{12}^1$, $e_7$ for $BN_{10}$ & \ref{sec:jwpieces}\\
\texttt{sefo.py} & converged norms $\|\Phi_k\|$ for the 7 matrices of order 6, up to level 9 & \ref{sec:jwpieces}\\
\bottomrule
\end{tabular}
\end{center}

\medskip

\noindent Each program comes with the modules that it uses, and the times indicated below are those of these laptops.

\subsection{Fourier matrices and $F_4^s$}

For the Fourier matrices $F_N$ with $N=2,\ldots,8$, and for $F_G$ with $G=\mathbb Z_2^2,\mathbb Z_2^3,\mathbb Z_2\times\mathbb Z_4$, the first 2 being the Walsh matrices $W_4,W_8$, we find $c_k=N^{k-1}$ for $k\leq4$, as it should be. This was the calibration run of the methods, done by the program \texttt{stage0.py}. At $N=4$, every complex Hadamard matrix is equivalent to a matrix of the following form, which is $F_2\otimes_QF_2$ with $Q=(^1_1{\ }^1_s)$, see \cite[Theorem 5.13]{banicaHM}:
$$F_4^s=\begin{pmatrix}1&1&1&1\\1&-1&1&-1\\1&s&-1&-s\\1&-s&-1&s\end{pmatrix}$$

\noindent We have $F_4^s\sim F_4^{-s}\sim F_4^{\bar{s}}\sim F_4^{-\bar{s}}$, so the relevant parameter is $s^2$, up to conjugation. With $s=e(a/n)$ for all $n\leq12$ and $(a,n)=1$, plus 3 random points, we find, for $H$ and $H^t$ alike:

\medskip

\begin{center}
\begin{tabular}{lcccc}
\toprule
$s$ & $c_1$ & $c_2$ & $c_3$ & $c_4$\\
\midrule
$s^4=1$ (Fourier, $\mathbb Z_2^2$ or $\mathbb Z_4$) & 1 & 4 & 16 & 64\\
$s^8=1$, primitive & 1 & 3 & 10 & 36\\
all other tested $s$, including the random ones & 1 & 3 & 10 & 35\\
\bottomrule
\end{tabular}
\end{center}

\medskip

\noindent This agrees with \cite{cn}, and the singular value gaps are all above $10^{13}$. Up to level 4, only $s^2$ of order 4 produces a jump. For $s^2$ of order 3 or 6, any jump must happen at $k\geq5$, and testing the $D_{2k}^{(1)}$ conjecture of \cite{cn} in these cases requires the levels $k\geq5$.

\subsection{Deformations of tensor products}\label{sec:dita}

The Di\c t\u a deformations are the basic source of families of complex Hadamard matrices. A conceptual approach to the deformations of the Fourier matrices, and more generally of the group commuting squares, is that of Nicoara and White \cite{nwh}: for a finite abelian group, explicit partial isometries $a$ produce one-parameter families $U_t=\exp(it(a+a^*))$, with $U_tF_N$ Hadamard for any $t\in\mathbb R$, which attain the defect bound, and which were combined into multi-parameter families by Worley \cite{wor}, see also \cite{nw4}. We use here the Di\c t\u a parametrization, the comparison with the families of \cite{nwh,wor} being an interesting question. At order 6 we have the following 2 families, with $F_6^{(rs)}=F_2\otimes_QF_3$ for $Q=(^1_1{\ }^1_r{\ }^1_s)$, and $F_6^{(^r_s)}=F_3\otimes_QF_2$ for $Q^t=(^1_1{\ }^1_r{\ }^1_s)$, see \cite[Proposition 6.19]{banicaHM} and \cite{bbs}:
$$F_6^{(rs)}=\begin{pmatrix}
1&1&1&1&1&1\\
1&w&w^2&1&w&w^2\\
1&w^2&w&1&w^2&w\\
1&r&s&-1&-r&-s\\
1&wr&w^2s&-1&-wr&-w^2s\\
1&w^2r&ws&-1&-w^2r&-ws
\end{pmatrix}\qquad
F_6^{(^r_s)}=\begin{pmatrix}
1&1&1&1&1&1\\
1&-1&1&-1&1&-1\\
1&r&w&wr&w^2&w^2r\\
1&-r&w&-wr&w^2&-w^2r\\
1&s&w^2&w^2s&w&ws\\
1&-s&w^2&-w^2s&w&-ws
\end{pmatrix}$$

\noindent At order 8 we have the 4 right deformations of the tensor products $F_2\otimes F_4$, $F_4\otimes F_2$, $F_2\otimes W_4$, $W_4\otimes F_2$, with $Q$ dephased, and the parameters $q_1,q_2,q_3$ filled into $Q$ row by row:

{\small\setlength{\arraycolsep}{3pt}$$F_2\otimes_QF_4=\begin{pmatrix}1 & 1 & 1 & 1 & 1 & 1 & 1 & 1\\ 1 & i & -1 & -i & 1 & i & -1 & -i\\ 1 & -1 & 1 & -1 & 1 & -1 & 1 & -1\\ 1 & -i & -1 & i & 1 & -i & -1 & i\\ 1 & q_{1} & q_{2} & q_{3} & -1 & -q_{1} & -q_{2} & -q_{3}\\ 1 & iq_{1} & -q_{2} & -iq_{3} & -1 & -iq_{1} & q_{2} & iq_{3}\\ 1 & -q_{1} & q_{2} & -q_{3} & -1 & q_{1} & -q_{2} & q_{3}\\ 1 & -iq_{1} & -q_{2} & iq_{3} & -1 & iq_{1} & q_{2} & -iq_{3}\end{pmatrix}$$}
{\small\setlength{\arraycolsep}{3pt}$$F_4\otimes_QF_2=\begin{pmatrix}1 & 1 & 1 & 1 & 1 & 1 & 1 & 1\\ 1 & -1 & 1 & -1 & 1 & -1 & 1 & -1\\ 1 & q_{1} & i & iq_{1} & -1 & -q_{1} & -i & -iq_{1}\\ 1 & -q_{1} & i & -iq_{1} & -1 & q_{1} & -i & iq_{1}\\ 1 & q_{2} & -1 & -q_{2} & 1 & q_{2} & -1 & -q_{2}\\ 1 & -q_{2} & -1 & q_{2} & 1 & -q_{2} & -1 & q_{2}\\ 1 & q_{3} & -i & -iq_{3} & -1 & -q_{3} & i & iq_{3}\\ 1 & -q_{3} & -i & iq_{3} & -1 & q_{3} & i & -iq_{3}\end{pmatrix}$$}
{\small\setlength{\arraycolsep}{3pt}$$F_2\otimes_QW_4=\begin{pmatrix}1 & 1 & 1 & 1 & 1 & 1 & 1 & 1\\ 1 & -1 & 1 & -1 & 1 & -1 & 1 & -1\\ 1 & 1 & -1 & -1 & 1 & 1 & -1 & -1\\ 1 & -1 & -1 & 1 & 1 & -1 & -1 & 1\\ 1 & q_{1} & q_{2} & q_{3} & -1 & -q_{1} & -q_{2} & -q_{3}\\ 1 & -q_{1} & q_{2} & -q_{3} & -1 & q_{1} & -q_{2} & q_{3}\\ 1 & q_{1} & -q_{2} & -q_{3} & -1 & -q_{1} & q_{2} & q_{3}\\ 1 & -q_{1} & -q_{2} & q_{3} & -1 & q_{1} & q_{2} & -q_{3}\end{pmatrix}$$}
{\small\setlength{\arraycolsep}{3pt}$$W_4\otimes_QF_2=\begin{pmatrix}1 & 1 & 1 & 1 & 1 & 1 & 1 & 1\\ 1 & -1 & 1 & -1 & 1 & -1 & 1 & -1\\ 1 & q_{1} & -1 & -q_{1} & 1 & q_{1} & -1 & -q_{1}\\ 1 & -q_{1} & -1 & q_{1} & 1 & -q_{1} & -1 & q_{1}\\ 1 & q_{2} & 1 & q_{2} & -1 & -q_{2} & -1 & -q_{2}\\ 1 & -q_{2} & 1 & -q_{2} & -1 & q_{2} & -1 & q_{2}\\ 1 & q_{3} & -1 & -q_{3} & -1 & -q_{3} & 1 & q_{3}\\ 1 & -q_{3} & -1 & q_{3} & -1 & q_{3} & 1 & -q_{3}\end{pmatrix}$$}

\noindent At order 9 we have the deformation of $F_3\otimes F_3$, with 4 parameters, which gives for instance the Fourier matrix $F_9$, up to equivalence, at $Q_{jd}=e^{2\pi ijd/9}$:

{\small\setlength{\arraycolsep}{3pt}$$F_3\otimes_QF_3=\begin{pmatrix}1 & 1 & 1 & 1 & 1 & 1 & 1 & 1 & 1\\ 1 & w & w^2 & 1 & w & w^2 & 1 & w & w^2\\ 1 & w^2 & w & 1 & w^2 & w & 1 & w^2 & w\\ 1 & q_{1} & q_{2} & w & wq_{1} & wq_{2} & w^2 & w^2q_{1} & w^2q_{2}\\ 1 & wq_{1} & w^2q_{2} & w & w^2q_{1} & q_{2} & w^2 & q_{1} & wq_{2}\\ 1 & w^2q_{1} & wq_{2} & w & q_{1} & w^2q_{2} & w^2 & wq_{1} & q_{2}\\ 1 & q_{3} & q_{4} & w^2 & w^2q_{3} & w^2q_{4} & w & wq_{3} & wq_{4}\\ 1 & wq_{3} & w^2q_{4} & w^2 & q_{3} & wq_{4} & w & w^2q_{3} & q_{4}\\ 1 & w^2q_{3} & wq_{4} & w^2 & wq_{3} & q_{4} & w & q_{3} & w^2q_{4}\end{pmatrix}$$}

\noindent At 3 random points of each family, for $H$ and $H^t$ alike, we find the following values, which are stable on all the points, the program being \texttt{stage2.py}:

\medskip

\begin{center}
\begin{tabular}{lccccc}
\toprule
family & parameters & $c_1$ & $c_2$ & $c_3$ & $c_4$\\
\midrule
$F_2\otimes_QF_3$, $F_3\otimes_QF_2$ & 2 & 1 & 4 & 18 & 86\\
$F_2\otimes_QF_4$, $F_4\otimes_QF_2$, $F_2\otimes_QW_4$, $W_4\otimes_QF_2$ & 3 & 1 & 5 & 28 & 167\\
$F_3\otimes_QF_3$ & 4 & 1 & 5 & 29 & 181\\
\bottomrule
\end{tabular}
\end{center}

\medskip

\noindent Thus the generic values depend only on the sizes $\{M,N\}$, and not on the order of the factors, nor on $F_4$ versus $W_4$, as predicted by \cite{bb3}, see section \ref{sec:bb}. The special values, found by setting each parameter in turn to a root of unity of order $\leq12$, the others being random, and all the parameters equal to such a root, at $k\leq4$ for $N=6$ and $k\leq3$ for $N=8,9$, are as follows. At $N=6$, all parameters $\pm1$ give $F_6$, with $c_2,c_3,c_4=6,36,216$, all parameters of order 3 or 6 give $c_4=90$, and all parameters $\pm i$ give $c_4=88$. At $N=8,9$, the only special point is the Fourier one.

\subsection{The families $F_{8,5a}$, $F_{8,5b}$, $F_{8,5c}$}\label{sec:f85}

Among the regular complex Hadamard matrices of order 8, classified in \cite{gt}, there are 7 families in which only the cycle structures $A,B$ appear. A direct search for the Di\c t\u a block structure, namely row groups whose quotients are constant on a common column partition, done by the program \texttt{dita\_test.py}, shows that 3 of them have Di\c t\u a block structures, and these are the following 5-parameter families, $F_{8,5a}$ and $F_{8,5b}$ being of Di\c t\u a type in the row convention of section 1, and $F_{8,5c}$ in the transposed one, that is, $F_{8,5c}^t$ is of Di\c t\u a type, as explained in section \ref{sec:nondita}, with $a,b,c,d,e\in\T$ independent, displayed as in \cite{gt}. The other 4 families have no such structure, and are discussed in section \ref{sec:wild8}.

\begin{center}\normalsize\setlength{\arraycolsep}{3.8pt}
\begin{adjustbox}{max width=\textwidth}$\displaystyle F_{8,5a}(a,b,c,d,e)=\begin{pmatrix}
1 & 1 & 1 & 1 & 1 & 1 & 1 & 1 \\
1 & -1 & -bd & bd & -bd & bd & 1 & -1 \\
1 & -\frac{1}{c} & -\frac{bd}{a} & d & \frac{bd}{a} & -d & -1 & \frac{1}{c} \\
1 & 1 & -1 & -1 & -1 & -1 & 1 & 1 \\
1 & -\frac{1}{c} & \frac{bd}{a} & -d & -\frac{bd}{a} & d & -1 & \frac{1}{c} \\
1 & -1 & bd & -bd & bd & -bd & 1 & -1 \\
1 & \frac{1}{c} & e & \frac{ae}{b} & -e & -\frac{ae}{b} & -1 & -\frac{1}{c} \\
1 & \frac{1}{c} & -e & -\frac{ae}{b} & e & \frac{ae}{b} & -1 & -\frac{1}{c}
\end{pmatrix}$\end{adjustbox}
\end{center}

\begin{center}\normalsize\setlength{\arraycolsep}{3.8pt}
\begin{adjustbox}{max width=\textwidth}$\displaystyle F_{8,5b}(a,b,c,d,e)=\begin{pmatrix}
1 & 1 & 1 & 1 & 1 & 1 & 1 & 1 \\
1 & -1 & -bd & bd & -bd & bd & 1 & -1 \\
1 & -\frac{1}{c} & -\frac{bd}{a} & d & \frac{bd}{a} & -d & -1 & \frac{1}{c} \\
1 & \frac{1}{c} & -\frac{bd}{a} & -d & \frac{bd}{a} & d & -1 & -\frac{1}{c} \\
1 & 1 & -1 & -1 & -1 & -1 & 1 & 1 \\
1 & -1 & bd & -bd & bd & -bd & 1 & -1 \\
1 & e & \frac{bd}{a} & cde & -\frac{bd}{a} & -cde & -1 & -e \\
1 & -e & \frac{bd}{a} & -cde & -\frac{bd}{a} & cde & -1 & e
\end{pmatrix}$\end{adjustbox}
\end{center}

\begin{center}\normalsize\setlength{\arraycolsep}{3.8pt}
\begin{adjustbox}{max width=\textwidth}$\displaystyle F_{8,5c}(a,b,c,d,e)=\begin{pmatrix}
1 & 1 & 1 & 1 & 1 & 1 & 1 & 1 \\
1 & -1 & ad & -ad & -bd & bd & 1 & -1 \\
1 & -\frac{1}{c} & d & d & -d & -d & -1 & \frac{1}{c} \\
1 & 1 & -1 & -1 & -1 & -1 & 1 & 1 \\
1 & -\frac{1}{c} & -d & -d & d & d & -1 & \frac{1}{c} \\
1 & -1 & -ad & ad & bd & -bd & 1 & -1 \\
1 & \frac{1}{c} & e & -e & \frac{be}{a} & -\frac{be}{a} & -1 & -\frac{1}{c} \\
1 & \frac{1}{c} & -e & e & -\frac{be}{a} & \frac{be}{a} & -1 & -\frac{1}{c}
\end{pmatrix}$\end{adjustbox}
\end{center}

\noindent At 3 random points, for $H$ and $H^t$ alike, we find:

\medskip

\begin{center}
\begin{tabular}{lcc}
\toprule
family & $c_1,c_2,\ldots$ & next eigenvalue at the last level\\
\midrule
$F_{8,5a}$ & $1,4,19,98,531,2974^*$ & 0.993\\
$F_{8,5b}$ & $1,3,11,45,197,903$ & 0.986\\
$F_{8,5c}$ & $1,3,11,45,197,903$ & $<0.95$\\
\bottomrule
\end{tabular}
\end{center}

\medskip

\noindent These values were obtained by the programs \texttt{sun.py}, up to level 5, and then \texttt{ocelot.py} and \texttt{hoost.py}, the latter working at level 6 by rotation sectors. The value $c_5=531^*$ comes from $526$ seeds and $5$ new fixed vectors. The next eigenvalue at level 5, namely $0.966$, is above the filter threshold $a=0.95$, but it is well separated from 1, and is correctly not counted. At level 6, the count by rotation sectors gives $c_6=903^*$ for $F_{8,5b}$ and $F_{8,5c}$, with the same sector contributions $165,139,159,142,159,139$, and $c_6=2974^*$ for $F_{8,5a}$, with sector contributions $520,474,516,474,516,474$, the sectors $s$ and $-s$ contributing equally in all cases. The same method gives again $c_5=531^*$ for $F_{8,5a}$, with sector contributions $107,106,106,106,106$, from independent random blocks. As for the next Ritz values at level 6, these are $0.986$ in all sectors for $F_{8,5b}$, and $0.993,0.967,0.981,0.967$ in the sectors $0,\pm1,\pm2,3$ for $F_{8,5a}$, while for $F_{8,5c}$ the filter finds no eigenvalue of $T_6$ in $(a,1)$, at the same parameters as for $F_{8,5b}$. Thus $F_{8,5b}$ and $F_{8,5c}$, which have the same $c_k$ as far as computed, are distinguished by the spectrum of $T_6$. The program \texttt{sandcat.py} computes the multiplicities $e_5,e_6$ of the eigenvalue 1 in the Jones--Wenzl pieces of section \ref{sec:jw}, and finds $e_5=9,9,55$ and $e_6=27,27$ for $F_{8,5b},F_{8,5c},F_{8,5a}$, which confirms $c_5=531$ and $c_6=903$. As for the special points, for $F_{8,5a}$ and $F_{8,5b}$, all parameters $\pm1$ give $c_2,c_3=8,64$, and all parameters $\pm i$ give $c_2,c_3=5,28$, the generic values of $F_2\otimes_QF_4$. For $F_{8,5c}$ we also get $c_2,c_3=4,19$, the generic values of $F_{8,5a}$, at $c=\pm1$, and along the diagonal at all the other tested roots of unity.

\section{Deformed Fourier, algebraic results}\label{sec:cumul}

\subsection{Deformed Fourier models}\label{sec:bb}

In this section $X,Y$ are finite abelian groups, written additively, with $M=|X|$ and $N=|Y|$, the double indices being $ia$ with $i\in X$, $a\in Y$. The parameter matrix $Q\in M_{X\times Y}(\T)$ is dephased, $Q_{0c}=Q_{i0}=1$, and we consider the deformed Fourier matrix $F_X\otimes_QF_Y$. The quantum groups of these matrices with $Q$ generic were computed in \cite{bb3}, and Bichon developed in \cite{bic} the tools for the non-generic case. We first recall what is needed from these papers, then state our main result, which covers all the parameters, and finally discuss the examples.

\medskip

\noindent\textbf{The free wreath product.} For a Di\c t\u a deformation, the magic matrix is a $Q$-twisted tensor product of those of the factors,
$$(U^{H\otimes_QK}_{ia,jb})_{kc,ld}=\frac{Q_{ic}Q_{jd}}{Q_{id}Q_{jc}}\,(U^H_{ij})_{kl}(U^K_{ab})_{cd}$$
and the deformation can only decrease the truncated moments $c_p^r=Tr(T_p^r)$, as shown in \cite[Theorem 2.8]{bb3}. The free wreath products, recalled in section \ref{sec:hopf}, enter as follows \cite{bi0,bb3}:

\begin{proposition}\label{prop:wreath}
For $H\in M_M(\T)$, $K\in M_N(\T)$ and any $Q\in M_{M\times N}(\T)$, the representation $u_{ia,jb}\to U^{H\otimes_QK}_{ia,jb}$ of $C(S_{MN}^+)$ factorizes through $C(S_N^+\wr_*S_M^+)$, via $v_{ij}\to U^H_{ij}\otimes1$ and $u^{(i)}_{ab}\to\sum_jU^{H\otimes_QK}_{ia,jb}$. In particular $G_{H\otimes_QK}\subset S_N^+\wr_*G_H$.
\end{proposition}

\begin{proof}
Let $W=U^{H\otimes_QK}$. Since $U^K$ is magic, $\sum_a(U^K_{ab})_{cd}=\delta_{cd}$, and the $Q$-ratio equals 1 at $c=d$, so $V_{ij}=\sum_aW_{ia,jb}=U^H_{ij}\otimes1$, independently of $b$. Thus $V$ is magic, and $U^{(i)}_{ab}=\sum_jW_{ia,jb}$ is magic too: its entries are sums of orthogonal projections from the same row of $W$, its row sums are $\sum_{jb}W_{ia,jb}=1$, and its column sums are $\sum_jV_{ij}=1$. Now $W_{ia,j'b}\leq V_{ij'}$ and $W_{ia',jb}\leq V_{ij}$, and $V_{ij'}V_{ij}=0$ for $j'\neq j$, so $U^{(i)}_{ab}V_{ij}=\sum_{a'}W_{ia,jb}W_{ia',jb}=W_{ia,jb}$, and similarly $V_{ij}U^{(i)}_{ab}=W_{ia,jb}$. This gives the factorization. Finally, the elements $V_{ij}=U^H_{ij}\otimes1$ satisfy the relations of $C(G_H)$, so the factorization descends to $C(S_N^+\wr_*G_H)$, and the last assertion follows from the minimality of the Hopf image.
\end{proof}

The same proof applies to the general form $H\otimes_Q(K_1,\ldots,K_M)$ of section 1, because the computation of $V_{ij}$ only uses the orthogonality of the rows of $K_i$, and so $G_H\subset S_N^+\wr_*S_M^+$ for all the matrices of Di\c t\u a type, the $M$ copies of $S_N^+$ being free.

\noindent\textbf{The Fourier case.} For $H=F_X$ the entries $U_{ij}$ depend only on $i-j$, and the factorization descends to $C(Y\wr_*X)$, where $v$ commutes with all the $u^{(i)}$. The Fourier transform $C(Y)\simeq C^*(Y)$ on each copy gives then an isomorphism $C(Y\wr_*X)\simeq C^*(Y^{*X})\rtimes C(X)$, where $c^{(i)}$ denotes the element $c\in Y$ of the $i$-th copy of $Y$, and where $C^*(\Gamma)\rtimes C(X)$ is the crossed coproduct, equal to $C^*(\Gamma)\otimes C(X)$ as an algebra, with the comultiplication twisted by the action of $X$ on $\Gamma$ by translation of the copies \cite[Section 3]{bb3}. With $\varepsilon_{ke}=\sum_iF_{ik}e_{ie}$, for $k\in X$ and $e\in Y$, the model is given by
$$\pi(c^{(i)})\,\varepsilon_{ke}=\theta_{ic}^{ke}\,\varepsilon_{k,e-c},\qquad\theta_{ic}^{ke}=\frac{Q_{i,e-c}\,Q_{i-k,e}}{Q_{ie}\,Q_{i-k,e-c}}$$
and $C(X)$ acts by scalars on the blocks $\mathrm{span}(\varepsilon_{ke})_e$, which are preserved by the elements $c^{(i)}$. Thus the words of total degree $0\in Y$ act diagonally, and so commute, and we obtain, for any $Q$, a factorization \cite[Theorem 3.10, Proposition 4.2]{bb3}
$$C(S_{MN}^+)\to C^*(\Gamma_{X,Y})\rtimes C(X)\to M_{MN}(\mathbb C),\qquad\Gamma_{X,Y}\simeq T\rtimes Y,\quad T\simeq\mathbb Z^{(M-1)(N-1)}$$
where $\Gamma_{X,Y}$ is the quotient of $Y^{*X}$ by the commutators of words of degree 0, and where $T$ is the subgroup of the words of degree 0, which is free abelian on the elements $t_{ic}=(-c)^{(0)}c^{(i)}$, with $i,c\neq0$. For $X=Y=\mathbb Z_2$ this is the infinite dihedral group.

\medskip

\noindent\textbf{Quotient data.} According to Bichon \cite{bic}, a quotient datum for $C^*(\Gamma)\rtimes C(X)$ is a pair $(N,\Phi)$, with $N\subset\Gamma$ a normal, $X$-stable subgroup, and $\Phi:N\to C(X)^\times$ a group morphism satisfying some compatibility conditions, and to such a pair corresponds a Hopf algebra quotient $C^*(\Gamma/N)\rtimes_\Phi C(X)$. Given a representation $\theta$ which is injective on $C(X)$, its Hopf image is of this form, with $(N,\Phi)$ being the maximal pair such that $\theta(r\otimes1)=\theta(1\otimes\Phi(r))$ for any $r\in N$ \cite[Proposition 4.1]{bic}. In our situation, this gives:

\begin{proposition}\label{prop:nq}
For any $Q$, the elements of $\Gamma_{X,Y}$ acting by scalars on all the blocks belong to $T$, and $t_R=\prod_{i,c\neq0}t_{ic}^{R_{ic}}$ acts on $\varepsilon_{ke}$ by a monomial $Q^{D_{ke}(R)}$, with $D_{ke}$ linear in $R$, given by
$$D_{ke}(t_{ic})=[i,e-c]+[i-k,e]-[i,e]-[i-k,e-c]+[-k,e-c]-[-k,e]$$
with $[j,d]$ denoting the exponent vector of $Q_{jd}$. If the group $N_Q$ of these elements is $X$-stable, then the quantum group of $F_X\otimes_QF_Y$ is given by $C^*(\Gamma_{X,Y}/N_Q)\rtimes_\Phi C(X)$, with $\Phi(t_R)(k)=Q^{D_{k0}(R)}$, and the law of its main character is given by
$$c_k=\frac{1}{M}\,\#\left\{(i_1,c_1),\ldots,(i_k,c_k)\in X\times Y\,\Big|\,c_1^{(i_1)}\ldots c_k^{(i_k)}\in N_Q\right\}$$
\end{proposition}

\begin{proof}
A word of degree $d\in Y$ maps $\varepsilon_{ke}$ to a multiple of $\varepsilon_{k,e-d}$, so it can act by scalars only if $d=0$. The formula of $D_{ke}$ follows from the formula of $\theta_{ic}^{ke}$, by applying $c^{(i)}$, then $(-c)^{(0)}$. The elements acting by scalars on all blocks form a group, which is normal, because the elements $c^{(j)}$ preserve the blocks, and $\Phi$ is a morphism. Since any pair $(N,\Phi)$ as above has $N\subset N_Q$, the first assertion follows from \cite[Proposition 4.1]{bic}. As for the law, we have $\chi=(\sum_{i,c}c^{(i)})\otimes\delta_0$, and the Haar functional of $C^*(\Gamma/N_Q)\rtimes_\Phi C(X)$ maps $r\otimes\delta_0$ to $\frac{1}{M}\Phi(r)(0)$ for $r\in N_Q$, and to $0$ otherwise. Since $D_{0e}=0$, all the elements of $N_Q$ act trivially on the block $k=0$, so $\Phi(r)(0)=1$, which gives the formula.
\end{proof}

\noindent We have checked the formula for $c_k$ against the direct computation with the model, at levels $k\leq4$, at many points, generic or not, of orders 6, 8 and 9, with the program \texttt{nongen.py}. For $Q$ generic, $N_Q=\{1\}$, and this is the theorem of \cite{bb3}, where the law was computed as follows, via an explicit embedding $\Gamma_{X,Y}\subset\mathbb Z^{(M-1)N}\rtimes Y$:

\begin{theorem}\label{thm:bb}
For $Q\in M_{X\times Y}(\T)$ generic, the quantum group associated to $F_X\otimes_QF_Y$ is given by $C(G)=C^*(\Gamma_{X,Y})\rtimes C(X)$, and the law of its main character is
$$\mu=\left(1-\frac{1}{N}\right)\delta_0+\frac{1}{N}\,law(A)$$
where $A=qq^*$ is the Gram matrix of the rows of a Haar distributed matrix $q\in M_{M\times N}(\T)$, with respect to $E\otimes tr_M$. Equivalently, $c_p$ is $\frac{1}{MN}$ times the number of pairs $i\in X^p$, $d\in Y^p$ such that the multisets $[(i_r,d_r)]_r$ and $[(i_r,d_{r-1})]_r$ coincide, with $d_0=d_p$.
\end{theorem}

\begin{proof}
This is \cite[Theorem 4.6, Propositions 5.1 and 5.2]{bb3}. With the partial sums $d_r=c_1+\ldots+c_r$, the word $c_1^{(i_1)}\ldots c_p^{(i_p)}$ is sent by the embedding to $(\sum_rb_{i_rd_{r-1}}-b_{i_rd_r},d_p)$, which gives the multiset formula, and the indicator of the multiset equality is $\int\prod_rq_{i_rd_r}\bar{q}_{i_{r+1}d_r}\,dq$, which gives the Gram matrix formula.
\end{proof}

\noindent These formulae are implemented in \texttt{bbformula.py} and \texttt{bb\_moments.py}. The Gram matrix formula gives $c_2=M+N-1$ and $c_3=M^2+N^2+3MN-3M-3N+2$, and the multiset formula gives the generic values below. In the regime $M=\alpha K$, $N=\beta K$ with $K\to\infty$, only the noncrossing partitions contribute at leading order, by Biane's inequality \cite{bia}, and $c_p\simeq K^{p-1}\sum_{\pi\in NC(p)}\alpha^{|\pi|-1}\beta^{p-|\pi|}$, in agreement with the Marchenko--Pastur behavior of $qq^*$ \cite[Lemma 5.3]{bb3}. Note also that $\Gamma_{X,Y}$ is amenable, so that $G$ is coamenable.

\medskip

\begin{center}
\begin{tabular}{lccccc}
\toprule
$|X|\times|Y|$ & $c_1$ & $c_2$ & $c_3$ & $c_4$ & $c_5$\\
\midrule
$2\times2$ & 1 & 3 & 10 & 35 & 126\\
$2\times3$, $3\times2$ & 1 & 4 & 18 & 86 & 426\\
$2\times4$, $4\times2$ & 1 & 5 & 28 & 167 & 1036\\
$3\times3$ & 1 & 5 & 29 & 181 & 1181\\
\bottomrule
\end{tabular}
\end{center}

\subsection{The main theorem}\label{sec:nongen}

Our main result describes $N_Q$ and $\Phi$ for any $Q$. Let $\mathbb Z[X\times Y]$ be the group ring, with the group element $(j,d)$ written $x^jy^d$, let $I_X\subset\mathbb Z[X]$ and $I_Y\subset\mathbb Z[Y]$ be the augmentation ideals, and identify $T$ with $I_X\otimes I_Y\subset\mathbb Z[X\times Y]$, via
$$t_{ic}\to(x^i-1)(y^c-1)$$
We also use the following element, encoding the parameters:
$$\theta=\sum_{j,d}\theta_{jd}\,x^jy^d\in(\mathbb R/\mathbb Z)[X\times Y],\qquad e^{2\pi i\theta_{jd}}=Q_{-j,d}$$
the space $(\mathbb R/\mathbb Z)[X\times Y]$ being viewed as a $\mathbb Z[X\times Y]$-module, so that $r\theta$ makes sense for $r\in\mathbb Z[X\times Y]$.

\begin{theorem}\label{thm:cyclic}
For any finite abelian groups $X,Y$, and any dephased $Q$, the quantum group of $F_X\otimes_QF_Y$ is $C^*(\Gamma_{X,Y}/N_Q)\rtimes_\Phi C(X)$, where
$$N_Q=\left\{r\in I_X\otimes I_Y\,\Big|\,r\theta\in(\mathbb R/\mathbb Z)^X+(\mathbb R/\mathbb Z)^Y\right\}$$
that is, $N_Q$ consists of the elements $r$ such that $r\theta$ is of the form $a(j)+b(d)$, a condition which reads $(x-1)(y-1)r\theta=0$ when $X,Y$ are cyclic. Moreover:
\begin{enumerate}
\item Write $Q_{jd}=\zeta_{jd}\prod_s\gamma_s^{\beta_s(j,d)}$, with $\zeta_{jd}$ roots of unity, and with $\gamma_s$ multiplicatively independent, and not roots of unity. Then $N_Q=\{1\}$, that is, the quantum group is the generic one, if and only if, for any pair $(\chi,\psi)$ of nontrivial characters of $X,Y$, there is an $s$ such that
$$\sum_{j,d}\beta_s(j,d)\,\bar{\chi}(j)\psi(d)\neq0$$

\item If $Q_{jd}=e^{2\pi ia_{jd}/L}$ for all $j,d$, the quantum group is finite, of dimension $MN[T:N_Q]$, where $[T:N_Q]$ is the order of the image of $I_X\otimes I_Y$ in $(\mathbb Z/L)[X\times Y]/((\mathbb Z/L)^X+(\mathbb Z/L)^Y)$, under $r\to r\hat{\theta}$, with $\hat{\theta}=\sum_{j,d}a_{-j,d}\,x^jy^d$.

\item In all cases, the law of the main character is given by the formula in Proposition \ref{prop:nq}, and only depends on $N_Q$.
\end{enumerate}
\end{theorem}

\begin{proof}
With $\varphi_{jd}=\theta_{-j,d}$, so that $e^{2\pi i\varphi_{jd}}=Q_{jd}$, Proposition \ref{prop:nq} shows that $t_{ic}$ acts on $\varepsilon_{ke}$ by $e^{2\pi iV(k,e)}$, with
$$V(k,e)=\varphi_{i,e-c}+\varphi_{i-k,e}-\varphi_{ie}-\varphi_{i-k,e-c}+\varphi_{-k,e-c}-\varphi_{-k,e}$$
On the other hand, with $F=(x^i-1)(y^c-1)\theta$, a direct expansion gives $F(k,e)-F(0,e)=-V(k,e)$. Thus the element of $T$ corresponding to $r$ acts by scalars on the block $k$ precisely when $F(k,e)-F(k,0)-F(0,e)+F(0,0)=0$ for any $e$, with $F=r\theta$, and the condition for all $k$ is that $r\theta$ is of the form $a(j)+b(d)$. This condition is invariant under the multiplication by the elements of $X$ and of $Y$, which implement the action of $X$ on $T$, and the conjugation by $Y$, so $N_Q$ is normal and $X$-stable. As for the compatibility conditions in the definition of a quotient datum \cite[Definition 2.1]{bic}, the action of $X$ on $\Gamma_{X,Y}$ is the translation of the copies, $c^{(i)}\to c^{(i+j)}$, which corresponds on $T$ to the multiplication by $x^j$. With $F_r=r\theta$, the scalar by which $r\in N_Q$ acts on the block $k$ is $\Phi_r(k)=e^{2\pi i(F_r(0,0)-F_r(k,0))}$, and since $F_{x^jr}(k,e)=F_r(k-j,e)$, we obtain
$$\Phi_r(l+h)=\Phi_r(l)\,\Phi_{x^{-l}r}(h)$$
which is the required compatibility between the fibres. Also, $\Phi$ is invariant under the conjugation by $\Gamma_{X,Y}$, because its elements preserve the blocks, on which $r$ acts by scalars. Thus $(N_Q,\Phi)$ is a quotient datum, and Proposition \ref{prop:nq} applies. For (1), we have $N_Q=\{1\}$ if and only if $N_Q\otimes\mathbb Q=0$. Multiplying by an integer kills the torsion part of $\theta$, and since the numbers $\frac{1}{2\pi}\log\gamma_s$, together with 1, are linearly independent over $\mathbb Q$, we see that $r\in N_Q\otimes\mathbb Q$ if and only if $r\hat{\beta}_s\in\mathbb Q^X+\mathbb Q^Y$ for any $s$, where $\hat{\beta}_s=\sum_{j,d}\beta_s(-j,d)x^jy^d$. Now by decomposing $\mathbb Q[X\times Y]$ into its simple components, indexed by the Galois orbits of the pairs of characters $(\chi,\psi)$, the functions of the form $a(j)+b(d)$ are those supported by the components with $\chi=1$ or $\psi=1$, while $I_X\otimes I_Y\otimes\mathbb Q$ is the sum of the other components. Thus $N_Q\otimes\mathbb Q$ is the sum of the components with $\chi,\psi\neq1$ on which all the $\hat{\beta}_s$ vanish, and the value of $\hat{\beta}_s$ at $(\chi,\psi)$ being the above sum, this gives (1). Finally, (2) is clear from the description of $N_Q$, the dimension of $C^*(\Gamma_{X,Y}/N_Q)\rtimes_\Phi C(X)$ being $|Y|[T:N_Q]\cdot|X|$, and (3) is Proposition \ref{prop:nq}.
\end{proof}

\noindent In other words, the generic locus is described by a finite number of components, one for each Galois orbit of pairs of nontrivial characters, each of them producing its own type of degeneration, and this explains all the phenomena below. All the formulas in this subsection and the next one were checked against a brute force computation of the elements of $\Gamma_{X,Y}$ acting by scalars, for $M\times N$ up to $5\times2$, $3\times5$ and $4\times4$, including the non-cyclic cases $X$ or $Y$ equal to $\mathbb Z_2^2$, or $\mathbb Z_2^3$, with more than $50,000$ elements tested. The programs, in the folder \texttt{tsukuba} of the archive, are \texttt{general.py} and \texttt{gen\_final.py} for the cyclic case, \texttt{abelian.py} and \texttt{klein.py} for the general abelian case, and \texttt{unified.py} for the case $M=2$ or $N=2$ below.

\subsection{Examples}

\noindent\textbf{The case $M=2$ or $N=2$.} For $X=\mathbb Z_2$ we have $(x-1)^2=-2(x-1)$, and the condition in Theorem \ref{thm:cyclic} becomes a condition in $\mathbb Z[Y]$. For $Y=\mathbb Z_n$, with augmentation ideal $I\subset\mathbb Z[\mathbb Z_n]$, and similarly for $Y=\mathbb Z_2$ and $X=\mathbb Z_n$, we obtain:

\begin{theorem}\label{thm:2n}
Consider the matrices $F_2\otimes_QF_n$ with $Q_1=(1,q_1,\ldots,q_{n-1})$, and $F_n\otimes_QF_2$ with $Q^t_1=(1,q_1,\ldots,q_{n-1})$, and let $\Theta=\sum_d\theta_dx^d$ in $(\mathbb R/\mathbb Z)[\mathbb Z_n]$, with $e^{2\pi i\theta_d}=q_d^2$, respectively with $e^{2\pi i\theta_d}=q_{-d}^2$, where $q_0=1$.
\begin{enumerate}
\item $N_Q=\{r\in I\,|\,(x-1)r\Theta=0\}$, that is, $N_Q$ consists of the elements $r$ such that $r\Theta$ is constant, and in the case $F_2\otimes_QF_n$, $\Phi(r)=(1,e^{2\pi i\lambda})$, where $\lambda\in\frac{1}{n}\mathbb Z/\mathbb Z$ is this constant.

\item The laws of the main characters of $F_2\otimes_QF_n$ with $Q_1=(1,q_1,\ldots,q_{n-1})$, and of $F_n\otimes_QF_2$ with $Q^t_1=(1,q_{n-1},\ldots,q_1)$, coincide.

\item If $q_d=e^{2\pi ia_d/L}$ for all $d$, the quantum group is finite, of dimension $2n[\mathbb Z^{n-1}:N_Q]$, with
$$[\mathbb Z^{n-1}:N_Q]=\left|\left(\psi I+L\mathbb Z[\mathbb Z_n]\right)/L\mathbb Z[\mathbb Z_n]\right|,\qquad\psi=2(x-1)\sum_da_dx^d$$
for $F_2\otimes_QF_n$, and with the reversed vector $(a_{-d})_d$ for $F_n\otimes_QF_2$.
\end{enumerate}
\end{theorem}

\begin{proof}
For $F_2\otimes_QF_n$, write $r=(x-1)\otimes s$ with $s\in I$. Since $\theta$ is supported by $x\otimes\mathbb Z[\mathbb Z_n]$, with $(x-1)x=1-x$, the condition $(x-1)(y-1)r\theta=0$ becomes $2(y-1)s\theta_1=0$, where $\theta_1$ is the row $j=1$ of $\theta$, which gives (1) with $\Theta=2\theta_1$, and the constant value $\lambda$ of $s\Theta$ satisfies $n\lambda=\varepsilon(s)\sum_d\theta_d=0$. The case of $F_n\otimes_QF_2$ is similar, with the roles of $x,y$ exchanged, the reflection $j\to-j$ in the definition of $\theta$ producing $q_{-d}$. As for (2), the transpose of $F_2\otimes_QF_n$ is equivalent to $F_n\otimes_{Q^t}F_2$, so by Theorem \ref{thm:transp} the law for $F_2\otimes_QF_n$ is that for $F_n\otimes_QF_2$ with $Q^t_1=(1,q_1,\ldots,q_{n-1})$, and the relabeling $i\to-i$ of the rows and the columns of $F_n$, which leaves $F_n$ invariant, shows that this matrix is equivalent to the one with $Q^t_1=(1,q_{n-1},\ldots,q_1)$. Note that the equality of the groups $N_Q$ alone would not suffice here, the ambient groups $\Gamma_{X,Y}$ being different. Finally, (3) is Theorem \ref{thm:cyclic} (2).
\end{proof}

\noindent For $n=2$, this gives back the dimension $4\cdot o(q^4)$ found by Bichon for $F_4^q$ \cite{bic}, the jump $c_4=36$ at $q^8=1$ primitive coming from a finite quotient of the infinite dihedral group. For $n$ prime, there is a single component in Theorem \ref{thm:cyclic} (1), and we obtain \cite[Theorem 5.1]{bic}: the quantum group is the generic one, as soon as one of the parameters is not a root of unity.

\medskip

\noindent\textbf{Order 6.} For $n=3$ the module $I$ is the ring $\mathbb Z[\omega]$ of Eisenstein integers, with $\omega=e^{2\pi i/3}$, and everything can be made explicit:

\begin{theorem}\label{thm:order6}
Consider $F_2\otimes_QF_3$ with $Q_1=(1,p,q)$, and $F_3\otimes_QF_2$ with $Q^t_1=(1,p,q)$.
\begin{enumerate}
\item For $F_2\otimes_QF_3$, with $\alpha=p^2$, $\beta=q^2$, the group $N_Q$ consists of the elements $t_1^mt_2^n$ with $\beta^m\alpha^n=\alpha^{-m-n}\beta^n=\alpha^m\beta^{-m-n}$, and $\Phi(t_1^mt_2^n)=(1,\beta^m\alpha^n)$, with values in $\{1\}\times\mathbb Z_3$.

\item For $F_3\otimes_QF_2$, with $B=p^2/q$, $C=q^2/p$, the group $N_Q$ consists of the elements $t_1^mt_2^n$ with $B^{2m}C^{2m+2n}=B^{2m+2n}C^{2n}=1$, and $\Phi(t_1^mt_2^n)=(1,(BC)^mC^n,B^m(BC)^n)$, with values in $\{1\}\times\mathbb Z_2^2$.

\item If $p=e^{2\pi ia/L}$, $q=e^{2\pi ib/L}$, the quantum group is finite, of dimension $6[\mathbb Z^2:N_Q]$, with
$$[\mathbb Z^2:N_Q]=\frac{L^2}{\left|\mathbb Z[\omega]/(L,w)\right|},\qquad w=2(1-\omega)(a+b\omega)=2(a+b)+2(2b-a)\omega$$
the same for the 2 matrices, as is the law of the main character.
\end{enumerate}
\end{theorem}

\begin{proof}
Here $t_1,t_2$ are $t_{11},t_{12}$ for $F_2\otimes_QF_3$, and $t_{11},t_{21}$ for $F_3\otimes_QF_2$. On the block $k=1$ of $F_2\otimes_QF_3$, the element $t_1^mt_2^n$ acts by $diag(\beta^m\alpha^n,\alpha^{-m-n}\beta^n,\alpha^m\beta^{-m-n})$, which gives (1), the product of these 3 numbers being 1. For $F_3\otimes_QF_2$, it acts on the blocks $k=1,2$ by $diag(x,x^{-1})$ and $diag(y,y^{-1})$, with $x=(BC)^mC^n$ and $y=B^m(BC)^n$, which gives (2). For (3), the equations in (1) are given by the integer matrix $2\begin{pmatrix}a+b&2a-b\\ 2b-a&a+b\end{pmatrix}$, acting modulo $L$, and those in (2) by its transpose. This matrix has determinant $12(a^2-ab+b^2)$, and the gcd of its entries is $2\gcd(a+b,2b-a)$, and the same holds for the matrix of the multiplication by $w$ on $\mathbb Z[\omega]$, whose determinant is the norm of $w$. Thus the 2 matrices have the same Smith normal form, and the order of the image modulo $L$ is $L^2/|\mathbb Z[\omega]/(L,w)|$.
\end{proof}

\noindent For instance, the Fourier point $p=q=1$ gives $\mathbb Z_6$, while $p=1$ and $q$ of order $3,4,12,8$ give dimensions $18,24,72,96$, with $c_4=90,88,86,86$ and $c_5=486,456,426,426$, against the generic $c_4,c_5=86,426$, and with $\Phi$ nontrivial in some of these cases, for instance for $F_2\otimes_QF_3$ with $q$ of order 3 or 12, and for $F_3\otimes_QF_2$ with $q$ of order 4, 8 or 12. This completes the partial results of \cite[Section 5]{bic} in the case $3\times2$. These values come from the programs \texttt{formula6.py}, \texttt{eisen.py} and \texttt{table6.py}, the last one comparing with the direct computation.

\medskip

\noindent\textbf{Order 8.} For $n=4$ the module $I=\mathbb Z[x]/(x^3+x^2+x+1)$ embeds with finite index into $\mathbb Z\times\mathbb Z[i]$, via $x\to(-1,i)$, and there are 2 components in Theorem \ref{thm:cyclic} (1), corresponding to $x\to-1$ and $x\to\pm i$. Thus, contrary to the case $n$ prime, the generic locus has a nontrivial structure even for parameters which are not roots of unity:

\begin{theorem}\label{thm:nongen}
Consider $F_2\otimes_QF_4$, with $Q_1=(1,q_1,q_2,q_3)$, with not all $q_c$ roots of unity, and set $r=q_1q_3/q_2$. Then exactly one of the following happens:
\begin{enumerate}
\item $r$ is not a root of unity, and $q_2,q_1/q_3$ are not both roots of unity. Then the quantum group is the generic one, $C^*(\mathbb Z^3\rtimes\mathbb Z_4)\rtimes C(\mathbb Z_2)$.

\item $r$ is a root of unity. Then $N_Q=\langle t^m\rangle$, with $t=t_{11}t_{12}^{-1}t_{13}$ and $m=o(r^4)$, and $\Phi(t^m)=(1,r^{2m})$, with $r^{2m}=\pm1$.

\item $q_2$ and $q_1/q_3$ are roots of unity. Then, with $\alpha=(q_3/q_1)^2$ and $\beta=q_2^2$, the group $N_Q$ consists of the elements $t_{11}^at_{12}^bt_{13}^{-a}$ with $\alpha^a\beta^b=\alpha^b\beta^{-a}=\pm1$, and $\Phi$ is given by this sign, on the block $k=1$.
\end{enumerate}
For $F_4\otimes_QF_2$, with $Q^t_1=(1,q_3,q_2,q_1)$, the same trichotomy holds, with the same subgroups $N_Q$ of $I$, and the same laws of the main character.
\end{theorem}

\begin{proof}
The component $x\to-1$ of $\Theta=2\sum_d\theta_dx^d$ is $2(\theta_0-\theta_1+\theta_2-\theta_3)$, which corresponds to $r^{-2}$, and the component $x\to i$ is $2(\theta_0-\theta_2)+2i(\theta_1-\theta_3)$, which corresponds to $q_2$ and $q_1/q_3$. Thus Theorem \ref{thm:cyclic} (1) gives the trichotomy, the cases (2) and (3) being exclusive, because together they would force all the parameters to be roots of unity. In case (2), $N_Q$ is contained in the line of $(1,-1,1)$, whose integer points are the powers $t^j$, and these act on the block $k=1$ by $r^{2j},r^{-2j},r^{2j},r^{-2j}$, so by scalars precisely when $r^{4j}=1$. Case (3) is similar, the element $t_{11}^at_{12}^bt_{13}^{-a}$ acting on the block $k=1$ by $x,y,x^{-1},y^{-1}$, with $x=\alpha^a\beta^b$, $y=\alpha^b\beta^{-a}$. The last assertion follows from Theorem \ref{thm:2n} (1) and (2).
\end{proof}

\noindent Observe that the quantum groups themselves are different for $F_2\otimes_QF_4$ and $F_4\otimes_QF_2$. In the generic case, the characters of $C^*(T\rtimes Y)\otimes C(X)$ correspond to those of $T_Y\times Y$, times the points of $X$, where $T_Y$ is the group of coinvariants. For $Y=\mathbb Z_4$ we have $T_Y\simeq\mathbb Z_4$, while for $Y=\mathbb Z_2$, acting by $-1$ on $T\simeq\mathbb Z^3$, we have $T_Y\simeq\mathbb Z_2^3$, and so there are $2\cdot4\cdot4=32$ characters in the first case, and $4\cdot2\cdot8=64$ in the second one.

\noindent In particular, the relation $q_2=q_1q_3$ makes the quantum group smaller than the generic one, even with $q_1,q_3$ not roots of unity, so the primality assumption in \cite[Theorem 5.1]{bic} cannot be removed. The law is then given by $c_1,\ldots,c_6=1,5,28,169,1076,7124$, instead of the generic $1,5,28,167,1036,6608$, the value $c_4=169$ being confirmed by the direct computation with the $8\times8$ model. At roots of unity of order dividing 8, we find the following values, the same for $F_2\otimes_QF_4$ with $Q_1=(1,q_1,q_2,q_3)$ and for $F_4\otimes_QF_2$ with $Q^t_1=(1,q_3,q_2,q_1)$, with $\zeta=e^{2\pi i/8}$, the values for $(1,i,i)$ and $(1,\zeta,\zeta)$ being confirmed by the direct computation, with the programs \texttt{index8.py} and \texttt{table8.py}:

\medskip

\begin{center}
\begin{tabular}{lccc}
\toprule
$(q_1,q_2,q_3)$ & $[\mathbb Z^3:N_Q]$ & dimension & $c_1,\ldots,c_4$\\
\midrule
$(1,1,1)$ & 1 & 8 & $1,8,64,512$\\
$(1,i,i)$ & 2 & 16 & $1,6,40,288$\\
$(1,1,i)$ & 4 & 32 & $1,5,28,176$\\
$(1,\zeta,\zeta)$ & 8 & 64 & $1,5,28,170$\\
$(1,1,\zeta)$ & 16 & 128 & $1,5,28,167$\\
\bottomrule
\end{tabular}
\end{center}

\medskip

\noindent The other deformed Fourier matrix of order 8 is $F_2\otimes_QF_{\mathbb Z_2^2}$, with $Q_1=(1,q_a,q_b,q_c)$, the indices being the elements $0,a,b,c$ of $\mathbb Z_2^2$. Here the 3 nontrivial characters of $\mathbb Z_2^2$ are rational, so there are 3 components, and Theorem \ref{thm:cyclic} (1) shows that, when not all the parameters are roots of unity, the quantum group is the generic one, except when one of the 3 numbers $q_a/q_bq_c$, $q_b/q_aq_c$, $q_c/q_aq_b$ is a root of unity, in which case $N_Q$ has rank 1, or 2 of them, in which case it has rank 2.

\medskip

\noindent\textbf{Order 9.} For $F_3\otimes_QF_3$ there are 2 components, corresponding to $(\chi,\psi)=(\omega,\omega)$ and $(\omega,\bar{\omega})$, and Theorem \ref{thm:cyclic} (1) shows that, when not all the parameters are roots of unity, the quantum group is the generic one, except in the following 2 cases, which are mutually exclusive:
\begin{enumerate}
\item $Q_{12}/Q_{21}$ and $Q_{11}Q_{22}/Q_{12}$ are roots of unity.

\item $Q_{11}/Q_{22}$ and $Q_{12}Q_{21}/Q_{11}$ are roots of unity.
\end{enumerate}
In both cases $N_Q$ has rank 2, and for instance with $Q_{12}=Q_{21}=\gamma$ and $Q_{11}Q_{22}=\gamma$, we find $c_1,\ldots,c_4=1,5,31,213$, instead of the generic $1,5,29,181$, the values $c_1,c_2,c_3$ being confirmed by the direct computation with the model, by the programs \texttt{three.py}, \texttt{three2.py} and \texttt{three3.py}. At roots of unity, we find for instance, the other parameters being 1:

\medskip

\begin{center}
\begin{tabular}{lccc}
\toprule
$F_3\otimes_QF_3$ & $[T:N_Q]$ & dimension & $c_1,\ldots,c_4$\\
\midrule
$Q=1$, or $Q_{jd}=\omega^{jd}$ & 1 & 9 & $1,9,81,729$\\
$Q_{11}=\omega$, or $Q_{11}=Q_{22}=\omega$, or $Q_{12}=Q_{21}=\omega$ & 3 & 27 & $1,5,33,261$\\
$Q_{11}=e^{2\pi i/6}$ & 48 & 432 & $1,5,29,181$\\
$Q_{11}=i$ & 256 & 2304 & $1,5,29,181$\\
\bottomrule
\end{tabular}
\end{center}

\medskip

\noindent\textbf{Other Di\c t\u a matrices.} For a Di\c t\u a matrix $H\otimes_QK$ which is not of deformed Fourier type, Proposition \ref{prop:wreath} still gives $G\subset S_N^+\wr_*G_H$, but the above method requires the inner pieces to be group algebras, as is the case when $K$ is a Fourier matrix. The families $F_{8,5a},F_{8,5b},F_{8,5c}$ of section \ref{sec:f85}, whose sequences $1,4,19,98,531$ and $1,3,11,45,197$ are strictly below the generic one $1,5,28,167,1036$ of the $8\times8$ deformations, are of this other type, with inner blocks of type $F_4^s$, which for $F_{8,5b}$ depend on the outer row, and are discussed at the end of the present section.

\subsection{Cumulants}

\noindent\textbf{Definitions.} The free cumulants $\kappa_n=\kappa_n(\chi)$ are given by the moment-cumulant formula over the noncrossing partitions, $c_k=\sum_{\pi\in NC(k)}\prod_{\beta\in\pi}\kappa_{|\beta|}$, see \cite{nsp}. The Boolean cumulants $b_n$ are given by the same formula over the interval partitions, or equivalently by $M(z)=1/(1-B(z))$, with $M(z)=1+\sum_kc_kz^k$ and $B(z)=\sum_nb_nz^n$, or by $1/G(z)=z-\sum_kb_{k+1}z^{-k}$, with $G$ the Cauchy transform of $\mu$. Combinatorially, $c_k$ counts the closed paths of length $k$ based at $1$ on the fusion graph of $u$, and $b_k$ counts those which return to $1$ for the first time at step $k$, that is, the fixed vectors of $u^{\otimes k}$ which do not factorize as $\xi\otimes\eta$ over a prefix of the legs.

\medskip

\noindent\textbf{Values.} For the deformed Fourier matrices, we have:

\medskip

\begin{center}
\begin{tabular}{lll}
\toprule
matrix & $\kappa_1,\kappa_2,\ldots$ & $b_1,b_2,\ldots$\\
\midrule
$F_2\otimes_QF_2$, generic & $1,2,3,2,-5,-20$ & $1,2,5,14,42,132$\\
$F_2\otimes_QF_3$, generic & $1,3,8,17,20,-40$ & $1,3,11,45,197,905$\\
$F_2\otimes_QF_4$, generic & $1,4,15,50,135$ & $1,4,19,100,562$\\
$F_3\otimes_QF_3$, generic & $1,4,16,60,200$ & $1,4,20,112,672$\\
$F_{8,5b}$, $F_{8,5c}$ & $1,2,4,8,16,32$ & $1,2,6,22,90,394$\\
$F_{8,5a}$ & $1,3,9,25,60,105^*$ & $1,3,12,55,273,1428^*$\\
\bottomrule
\end{tabular}
\end{center}

\medskip

\noindent Several observations are in order. For $F_2\otimes_QF_2$ the Boolean cumulants are the Catalan numbers, shifted by 1, while $\kappa_2\kappa_4-\kappa_3^2=2\cdot2-3^2=-5<0$. Since the free cumulants $(\kappa_{n+2})_{n\geq0}$ of a freely infinitely divisible law form the moment sequence of a positive measure, this shows that the law of $\chi$ is not freely infinitely divisible. For $F_2\otimes_QF_3$ the Boolean cumulants coincide with the little Schr\"oder numbers $1,3,11,45,197,903$ up to $n=5$. For $F_{8,5b}$ and $F_{8,5c}$ the free cumulants are $\kappa_n=2^{n-1}$, at the computed orders $n\leq6$, which suggests $\chi\sim2\pi_{1/2}$, with $\pi_t$ the free Poisson law of parameter $t$, and the Boolean cumulants are, at the same orders, the large Schr\"oder numbers, counting the separable permutations, those avoiding the patterns $2413$ and $3142$. Finally, for $F_{8,5a}$ the free cumulants are close to $3^{n-1}$ without being equal, while the Boolean cumulants are the ternary numbers $\frac{1}{2n+1}\binom{3n}{n}$, up to $n=6$. Note also that the sequence shifting left under 3 iterations of the INVERT transform, the 3-fold analogue of the little Schr\"oder numbers, is $1,4,19,100,562$, which is that of $F_2\otimes_QF_4$, and not that of $F_{8,5a}$.

\section{The families \texorpdfstring{$F_{8,5}$}{F(8,5)}, wreath products}\label{sec:f85q}

In this section we study the 3 families $F_{8,5a},F_{8,5b},F_{8,5c}$ of section \ref{sec:f85}, which are of Di\c t\u a type, $F_{8,5c}$ up to transposition, but not of deformed Fourier type, their inner blocks being matrices of type $F_4^s$, so that Theorem \ref{thm:cyclic} does not apply. We identify the candidates for their quantum groups, which are free wreath products, prove the inclusions, and reduce the equalities to freeness questions, for 2 copies of $S_4^+$, respectively of $H_2^+$, the main result being Theorem \ref{thm:f85b}.

\subsection{Free wreath products}

The moments $c_k=\sum_{\pi\in NC(k)}2^{k-|\pi|}$ of $2\pi_{1/2}$ are the little Schr\"oder numbers $1,3,11,45,197,903,\ldots$ (OEIS A001003), which are, up to a factor 2, the moments of the free Poisson law of parameter 2. Since $\lim_kc_k^{1/k}=(1+\sqrt2)^2=3+2\sqrt2$, which lies strictly between the free value 4 and the value $N=8$ of the coamenable case, the corresponding quantum group would not be coamenable, unlike for the generic deformations \cite{bb3}. At generic parameters, the values of section \ref{sec:f85} for $F_{8,5b}$ and $F_{8,5c}$ are the little Schr\"oder numbers, up to $c_6=903$, and for $F_{8,5a}$ the Boolean cumulants of $\chi$ are the ternary numbers $\frac{1}{2n+1}\binom{3n}{n}$, up to the value $c_6=2974^*$. For $F_{8,5b}$ and $F_{8,5c}$, there is a natural candidate for the quantum group. Consider the free wreath product $S_4^+\wr_*\mathbb Z_2\subset S_8^+$, consisting of 2 free copies of $S_4^+$, exchanged by $\mathbb Z_2$. With the notations of section \ref{sec:bb}, here $v=\begin{pmatrix}p&1-p\\1-p&p\end{pmatrix}$, the projection $p$ commutes with both copies $u^{(0)},u^{(1)}$, and so $C(S_4^+\wr_*\mathbb Z_2)=(C(S_4^+)*C(S_4^+))\otimes C(\mathbb Z_2)$ as an algebra, with the Haar state being the tensor product of the Haar states.

\begin{proposition}\label{prop:f85}
The main character of $S_4^+\wr_*\mathbb Z_2$ has law
$$\mu=\frac{1}{2}\,\delta_0+\frac{1}{2}\,\pi_2=D_2(\pi_{1/2})$$
so its moments are the little Schr\"oder numbers, and this quantum group is not coamenable. Moreover, the magic matrix $W$ of $F_{8,5b}$, and that of $F_{8,5c}^t$, admit a partition of the indices into 2 blocks of size 4, $ia$ with $i\in\{0,1\}$, such that $\sum_aW_{ia,jb}$ does not depend on $b$, and so the quantum groups of $F_{8,5b}$ and $F_{8,5c}^t$ are contained in $S_4^+\wr_*\mathbb Z_2$.
\end{proposition}

\begin{proof}
We have $\chi=\sum_{i,a}u^{(i)}_{aa}v_{ii}=p(\chi^{(0)}+\chi^{(1)})$, with $p$ central, of trace $\frac{1}{2}$, and with $\chi^{(0)},\chi^{(1)}$ free, each having law $\pi_1$, the free Poisson law, since $S_4^+$ has $c_k=Cat_k$. Thus $\mu=\frac{1}{2}\delta_0+\frac{1}{2}(\pi_1\boxplus\pi_1)=\frac{1}{2}\delta_0+\frac{1}{2}\pi_2$, whose moments are $c_k=\frac{1}{2}\sum_{\pi\in NC(k)}2^{|\pi|}$, the little Schr\"oder numbers, and the equality with $D_2(\pi_{1/2})$ follows by comparing the moments $\sum_{\pi\in NC(k)}2^{k-|\pi|}$. The support of $\pi_2$ is $[(\sqrt2-1)^2,(\sqrt2+1)^2]$, which does not contain 8, so $G$ is not coamenable. As for the second assertion, the condition is a polynomial identity in the parameters, which is checked directly, with the blocks $\{0,1,4,5\},\{2,3,6,7\}$ for $F_{8,5b}$, and $\{0,1,6,7\},\{2,3,4,5\}$ for $F_{8,5c}^t$. The proof of Proposition \ref{prop:wreath} then applies verbatim, and gives a factorization through $S_4^+\wr_*S_2^+=S_4^+\wr_*\mathbb Z_2$.
\end{proof}

\noindent The numerics, done by the programs \texttt{wstruct.py} and \texttt{even.py}, are consistent with these inclusions being equalities. At random parameters, each of the 2 inner copies $U^{(i)}_{ab}=\sum_jW_{ia,jb}$ is a $4\times4$ magic matrix whose Hopf image has $c_1,\ldots,c_5=1,2,5,14,42$, the values of $S_4^+$. By the classification of the quantum subgroups of $S_4^+$ \cite{bb2}, this excludes all of them except the twisted icosahedral group $A_5^\tau$, which has the same moments up to $c_5$, and differs only at $c_6$, and the fact that each copy generates $S_4^+$ is proved in Proposition \ref{prop:rkl} below. Also, the block diagonal $8\times8$ magic matrix formed by the 2 copies has $c_1,\ldots,c_4=2,6,22,90$, which are the moments of $\pi_2$, the values of the dual free product $S_4^+\,\hat{*}\,S_4^+$, so the 2 copies are free, up to level 4. Together with the agreement of the $c_k$ up to $k=6$, this raises the question whether the quantum groups of $F_{8,5b}$ and $F_{8,5c}^t$, at generic parameters, are both equal to $S_4^+\wr_*\mathbb Z_2$, which would make them non-coamenable, unlike the quantum groups of the generic Di\c t\u a deformations of \cite{bb3}. Since $\mu_H=\mu_{H^t}$ by Theorem \ref{thm:transp}, this would give $\chi\sim D_2(\pi_{1/2})$ for all the 4 matrices $F_{8,5b}$, $F_{8,5b}^t$, $F_{8,5c}$, $F_{8,5c}^t$. The case of $F_{8,5a}$ is different. Its magic matrix has both block structures, with 2 blocks of size 4 and with 4 blocks of size 2, and these are nested, the pairs $\{0,3\},\{1,5\}$ and $\{2,4\},\{6,7\}$ being contained in the blocks $\{0,1,3,5\}$ and $\{2,4,6,7\}$. When the 2 blocks of size 4 are labeled compatibly with the pairs, the 2 inner copies have $c_1,\ldots,c_4=1,3,10,35$, the values of the hyperoctahedral quantum group $H_2^+=\mathbb Z_2\wr_*\mathbb Z_2=O_2^{-1}$ \cite{bbc}, which is the quantum group of the generic $F_4^s$, and so the natural candidate is $H_2^+\wr_*\mathbb Z_2$, as discussed at the end of the present section.

\medskip

\subsection{Structure of $F_{8,5b}$}

We now explain the inclusion in Proposition \ref{prop:f85}, and what we can prove towards the equality with $S_4^+\wr_*\mathbb Z_2$, in the case of $F_{8,5b}$. Let $P_\pm\in M_2(\mathbb C)$ be the projections onto $\frac{1}{\sqrt2}(1,\pm1)$, and for $K,L\in M_4(\T)$ Hadamard, consider the matrix $R^{KL}_{ab}=Proj(K_a\bar{L}_b)$, where $K_a\bar{L}_b$ is the componentwise product of the row $K_a$ with the conjugate of the row $L_b$. This matrix is magic, because $\langle K_a\bar{L}_b,K_a\bar{L}_{b'}\rangle=\langle L_{b'},L_b\rangle=4\delta_{bb'}$, and similarly for the columns, and for $K=L$ we recover the magic matrix $U^K$.

\begin{proposition}\label{prop:f85b}
Up to permutations of rows and columns, we have
$$F_{8,5b}=\begin{pmatrix}K&K\\ L&-L\end{pmatrix},\quad K=\begin{pmatrix}1&1&1&1\\ 1&-1&-s&s\\ 1&1&-1&-1\\ 1&-1&s&-s\end{pmatrix},\quad L=\begin{pmatrix}1&-\bar{c}&-\frac{bd}{a}&d\\ 1&\bar{c}&-\frac{bd}{a}&-d\\ 1&e&\frac{bd}{a}&cde\\ 1&-e&\frac{bd}{a}&-cde\end{pmatrix}$$
with $s=bd$, and with respect to $\mathbb C^8=\mathbb C^2\otimes\mathbb C^4$, the magic matrix of $F_{8,5b}$ is
$$W=\begin{pmatrix}P_+\otimes U^K&P_-\otimes R^{KL}\\ P_-\otimes R^{LK}&P_+\otimes U^L\end{pmatrix}$$
Thus the outer magic matrix is $U^{F_2}\otimes1$, and the 2 inner copies are $u^{(0)}\to U^K\oplus R^{KL}$ and $u^{(1)}\to U^L\oplus R^{LK}$, acting on $(P_+\otimes\mathbb C^4)\oplus(P_-\otimes\mathbb C^4)$.
\end{proposition}

\begin{proof}
The first assertion is a reordering of the rows and columns of $F_{8,5b}$, with the rows $0,1,4,5$ and $2,3,6,7$, and the columns $0,1,2,3$ and $6,7,4,5$. The rows of $F_{8,5b}$ are then $(K_a,K_a)$ and $(L_a,-L_a)$, so the quotients are $(K_a/K_b,K_a/K_b)$, $(L_a/L_b,L_a/L_b)$, $(K_a\bar{L}_b,-K_a\bar{L}_b)$ and $(L_a\bar{K}_b,-L_a\bar{K}_b)$, whose projections are $P_+\otimes U^K_{ab}$, $P_+\otimes U^L_{ab}$, $P_-\otimes R^{KL}_{ab}$ and $P_-\otimes R^{LK}_{ab}$. Summing over the inner indices, and using the magic property of $U^K,U^L,R^{KL},R^{LK}$, gives the outer and inner magic matrices.
\end{proof}

\noindent We will also need the following matrix. With $D=diag(L_0)$, we have $\frac{1}{4}KDL^*=1\oplus V_3$, where, with $w=\bar{c}\bar{e}$,
$$V_3=\frac{1}{4}\begin{pmatrix}2(1-s)&(1+s)(1+w)&(1+s)(1-w)\\ 0&2(1-w)&2(1+w)\\ 2(1+s)&(1-s)(1+w)&(1-s)(1-w)\end{pmatrix}$$
Indeed, the entries of $\frac{1}{4}KDL^*$ are $\frac{1}{4}\sum_mK_{im}L_{0m}\bar{L}_{jm}$, and since $L_0\bar{L}_1=(1,-1,1,-1)$, $L_0\bar{L}_2=(1,-w,-1,w)$, $L_0\bar{L}_3=(1,w,-1,-w)$, a direct computation gives this formula.

\noindent Recall from \cite{bb2} that $S_4^+=SO_3^{-1}$ is a twist of $SO_3$. It has the Pauli model $\pi:C(S_4^+)\to C(SU_2,M_4(\mathbb C))$, given by $u_{ab}\to[x\to Proj(c_axc_b)]$, with $c_0,\ldots,c_3$ being $1,i\sigma_1,i\sigma_2,i\sigma_3$, acting on $M_2(\mathbb C)\simeq\mathbb C^4$, with scalar product $\langle a,b\rangle=tr(ab^*)$. This model is faithful, as shown in \cite{bc}, and its fibers $\pi_x$ were studied by Enders \cite{en}, who showed that they are generically irreducible, and locally pairwise inequivalent, and deduced from this that $C(S_4^+)$ is not semiprojective. We denote by $\rho:SU_2\to SO_3$ the covering map, given by $xc_kx^*=\sum_l\rho(x)_{lk}c_l$, and the fiber $\pi_x$ becomes, in the coordinates where $u\simeq1\oplus a$, the representation $\theta_{\rho(x)}$ below. The Hopf images of these fibers are as follows.

\begin{lemma}\label{lem:pauli}
Let $a$ be the 3-dimensional fundamental corepresentation of $S_4^+=SO_3^{-1}$, so that $u\simeq1\oplus a$, and for $y\in SO_3$ let $\theta_y:C(S_4^+)\to M_4(\mathbb C)$ be the Pauli representation $a_{ij}\to y_{ij}\tau_i\otimes\tau_j$, with $\tau_k=i\sigma_k$, which is the fiber at $y$ of the embedding of \cite[Theorem 5.1]{bb2}. Then the Hopf image of $\theta_g$ is $\Gamma_g^\tau$, the twist of the closed subgroup $\Gamma_g=\overline{\langle V,g\rangle}\subset SO_3$, where $V$ is the Klein group of diagonal matrices.
\end{lemma}

\begin{proof}
For a closed subgroup $V\subset\Gamma\subset SO_3$ the cocycle of \cite{bb2}, which is induced by the restriction to $V$, descends to $C(\Gamma)$, and the formula $a_{ij}\to\tau_i\otimes\tau_j\otimes x_{ij}$ gives, as in \cite[Theorem 5.1]{bb2}, an embedding $\theta_\Gamma:C(\Gamma^\tau)\to M_4(\mathbb C)\otimes C(\Gamma)$, whose fibers are the representations $\theta_y$, with $y\in\Gamma$. Since $\theta_g$ is a fiber of $\theta_{\Gamma_g}$, it factorizes through $C(\Gamma_g^\tau)$, and so its Hopf image is contained in $\Gamma_g^\tau$. Conversely, let $I$ be the Hopf ideal defining the Hopf image of $\theta_g$, and $S=\{y\in SO_3\,|\,\theta_y(I)=0\}$. This set is closed, and contains $g$. Let $\Omega=\frac{1}{\sqrt2}(e_1\otimes e_2-e_2\otimes e_1)$, which satisfies $(\tau_k\otimes\tau_k)\Omega=\Omega$ for $k=1,2,3$, and consider the isometry $J(\xi\otimes\eta)=\xi\otimes\Omega\otimes\eta$. We have then
$$(\theta_y\otimes\theta_z)\Delta(a_{ij})J=\sum_ky_{ik}z_{kj}(\tau_i\otimes\tau_k\otimes\tau_k\otimes\tau_j)J=(yz)_{ij}J(\tau_i\otimes\tau_j)=J\theta_{yz}(a_{ij})$$
so $\theta_{yz}$ is a subrepresentation of $(\theta_y\otimes\theta_z)\Delta$. Since $I$ is a Hopf ideal, it follows that $S$ is a closed subsemigroup of $SO_3$, hence a closed subgroup, and in particular $1\in S$. Now $\theta_1$ maps $a$ to the diagonal matrix with entries $\tau_k\otimes\tau_k$, which are commuting symmetries, so $\theta_1$ factorizes through $C(V)$, where $V\subset SO_3^{-1}$ is the classical diagonal subgroup, the twist being trivial on it, and the Hopf image of $\theta_1$ is $V$. Thus $I\subset\ker(C(S_4^+)\to C(V))$, and since each $\theta_v$ with $v\in V$ factorizes through $C(V)$, we obtain $V\subset S$. Thus $\Gamma_g\subset S$, so $\theta_{\Gamma_g}(I)=0$, and by injectivity of $\theta_{\Gamma_g}$ the ideal $I$ is contained in the kernel of $C(S_4^+)\to C(\Gamma_g^\tau)$. Thus the Hopf image contains $\Gamma_g^\tau$.
\end{proof}

\noindent This agrees with the values found numerically for the Pauli models, by the programs \texttt{pauli.py}, \texttt{fibers.py} and \texttt{finite.py}: $O_2$ around a coordinate axis gives $c_1,\ldots,c_4=1,3,10,35$, $D_4$ gives $1,3,10,36$, $A_4$ gives $1,2,6,22$, and $SO_3$ gives $1,2,5,14$. In order to apply this to $F_{8,5b}$, we need to know which fibers appear. Here we have the following general fact, which comes from the faithfulness of $\pi$:

\begin{lemma}\label{lem:fiber}
Every irreducible representation $\sigma:C(S_4^+)\to M_4(\mathbb C)$ is unitarily equivalent to a fiber $\pi_x$, with $x\in SU_2$, and we have then, for any $a,b\in\{1,2,3\}$,
$$Tr(\sigma(u_{ab})\sigma(u_{00}))=\rho(x)_{ab}^2$$
where $Tr$ is the usual, non-normalized trace of $M_4(\mathbb C)$.
\end{lemma}

\begin{proof}
Let $\xi\in\mathbb C^4$ be a unit vector. Since $\sigma$ is irreducible, $\varphi=\langle\sigma(\cdot)\xi,\xi\rangle$ is a pure state of $C(S_4^+)$, which by \cite{bc} is a $C^*$-subalgebra of $C(SU_2,M_4(\mathbb C))$, via $\pi$. The extensions of $\varphi$ to $C(SU_2,M_4(\mathbb C))$ form a weak$^*$-closed face of the state space, whose extreme points are therefore pure states, and these are of the form $f\to\langle f(x)\eta,\eta\rangle$. Thus $\varphi=\langle\pi_x(\cdot)\eta,\eta\rangle$, and by uniqueness of the GNS construction, $\sigma$ is equivalent to the cyclic subrepresentation of $\pi_x$ generated by $\eta$, which is 4-dimensional, and so is $\pi_x$ itself. As for the formula, $c_axc_b$ and $x$ are unit vectors of $M_2(\mathbb C)$, and since $tr(c_ac_l)=-\delta_{al}$ for $a,l\geq1$, we have
$$\langle c_axc_b,x\rangle=tr(c_axc_bx^*)=\sum_l\rho(x)_{lb}\,tr(c_ac_l)=-\rho(x)_{ab}$$
so that $Tr(\pi_x(u_{ab})\pi_x(u_{00}))=|\langle c_axc_b,x\rangle|^2=\rho(x)_{ab}^2$, as claimed.
\end{proof}

\noindent With this, the 2 inner copies can be identified by hand, with the above matrix $V_3$ as the only input:

\begin{proposition}\label{prop:rkl}
Assume $s\neq\pm1$ and $w\notin\{\pm1,\pm i\}$. Then $R^{KL}$ is an irreducible representation of $C(S_4^+)$, equivalent to a Pauli fiber $\pi_x$ with $|\rho(x)_{ab}|=|(V_3)_{ab}|$ for any $a,b$, and its Hopf image is $S_4^+$. The same holds for $R^{LK}$.
\end{proposition}

\begin{proof}
Let $\xi_{ab}=K_a\bar{L}_b$, so that $R^{KL}_{ab}=Proj(\xi_{ab})$, with $||\xi_{ab}||^2=4$. Since $K_0=1$, the above formula of $V_3$ gives $\langle\xi_{ab},\xi_{00}\rangle=\sum_mK_{am}L_{0m}\bar{L}_{bm}=4(1\oplus V_3)_{ab}$. Also, $\langle\xi_{a0},\xi_{0b}\rangle=\sum_mK_{am}\bar{L}_{0m}L_{bm}$ is the conjugate of the same quantity with $K$ replaced by $\bar{K}$, that is, with $s$ replaced by $\bar{s}$, and since the entries of $V_3$ are products of constants and of the numbers $1\pm s$, $1\pm w$, whose moduli don't change under $s\to\bar{s}$, we obtain $|\langle\xi_{a0},\xi_{0b}\rangle|=4|(1\oplus V_3)_{ab}|$ as well.

Now let $T\in M_4(\mathbb C)$ commute with all the $R^{KL}_{ab}$. Both $(R^{KL}_{0b})_b$ and $(R^{KL}_{a0})_a$ are families of 4 orthogonal rank 1 projections summing up to 1, so we have $T\xi_{0b}=\mu_b\xi_{0b}$ and $T\xi_{a0}=\lambda_a\xi_{a0}$, for certain scalars $\lambda_a,\mu_b$, and the same holds for $T^*$, with the conjugate scalars. Thus
$$\lambda_a\langle\xi_{a0},\xi_{0b}\rangle=\langle T\xi_{a0},\xi_{0b}\rangle=\langle\xi_{a0},T^*\xi_{0b}\rangle=\mu_b\langle\xi_{a0},\xi_{0b}\rangle$$
so that $\lambda_a=\mu_b$ when $(V_3)_{ab}\neq0$. For $s,w\neq\pm1$ all the entries of $V_3$ are nonzero, except for the $(2,1)$ one, so the corresponding bipartite graph is connected, and we obtain $\mu_1=\mu_2=\mu_3=\mu$. Thus $T=\mu_0R^{KL}_{00}+\mu(1-R^{KL}_{00})$. But $|(V_3)_{22}|=\frac{1}{2}|1-w|\in(0,1)$, so $\xi_{22}$ is neither proportional nor orthogonal to $\xi_{00}$, and $R^{KL}_{00}$ does not commute with $R^{KL}_{22}$, which gives $\mu_0=\mu$. Thus $R^{KL}$ is irreducible.

By Lemma \ref{lem:fiber}, $R^{KL}$ is equivalent to a fiber $\pi_x$, with $\rho(x)_{ab}^2=Tr(R^{KL}_{ab}R^{KL}_{00})=\frac{1}{16}|\langle\xi_{ab},\xi_{00}\rangle|^2=|(V_3)_{ab}|^2$. Let $g=\rho(x)$. By Lemma \ref{lem:pauli}, the Hopf image is $\Gamma_g^\tau$, and we must prove that $\Gamma_g=SO_3$. The closed subgroups of $SO_3$ containing $V$ are $SO_3$ itself, the subgroups of the 3 copies of $O_2$ around the coordinate axes, whose elements have a diagonal entry $\pm1$, and the tetrahedral, octahedral and icosahedral groups containing $V$. The pairs formed by such a group and a Klein subgroup are conjugate in $SO_3$, except for the octahedral group, where $V$ can be normal or not, and since the normalizer of $V$ consists of signed permutation matrices, a direct check on representatives shows that the entries of the elements of these groups have moduli in the set $F=\{0,\frac{1}{2},\frac{1}{\sqrt2},\frac{\varphi}{2},\frac{1}{2\varphi},1\}$, with $\varphi$ the golden ratio. Now for $s,w\neq\pm1$ the entries of $V_3$ have moduli $<1$, so $\Gamma_g$ is not contained in a copy of $O_2$. Also, the second row of $g$ has moduli $0,\frac{1}{2}|1-w|,\frac{1}{2}|1+w|$, and the only way of writing 1 as a sum of 2 squares of nonzero elements of $F$ is $1=\frac{1}{2}+\frac{1}{2}$, which would require $|1-w|^2=2$, that is $w=\pm i$. Thus $g$ belongs to no finite subgroup containing $V$, so $\Gamma_g$ is infinite, and $\Gamma_g=SO_3$. Finally, $R^{LK}_{ab}=(R^{KL}_{ba})^t$, so $R^{LK}$ is obtained from $R^{KL}$ by composing with the antipode and transposing, which does not change the Hopf image.
\end{proof}

\noindent Numerically, as checked by the program \texttt{rklcheck.py}, $x$ is a lift of $r_2(\delta)r_1(\gamma)$, where $r_k$ denotes the rotation around the $k$-th axis, and $\delta=\frac{\pi}{2}-\frac{1}{2}\arg(bd)$, $\gamma=\frac{\pi}{2}+\frac{1}{2}\arg(ce)$, but the above proof only uses the moduli of the entries of $\rho(x)$, which can be read on $V_3$. Note also that the numerics alone could hardly have done this step: the twist of the icosahedral group $A_5\supset V$ has $c_1,\ldots,c_5=1,2,5,14,42$, like $S_4^+$, and only $c_6$ separates them, with $c_6=133$ for $A_5^\tau$, and $c_6=132$ for $R^{KL}$.

\medskip

\subsection{Freeness, reduction}

What remains is the freeness of the 2 copies, which we have not been able to establish. We can however reduce it, in the spirit of \cite{bb3,bic}, to a question about a single model of the dual free product $S_4^+\,\hat{*}\,S_4^+$. Let $A=C(S_4^+)*C(S_4^+)$, with the 2 copies $u^{(0)},u^{(1)}$, and $\nu$ the automorphism exchanging them. Since the outer model is the Fourier model $U^{F_2}$, the factorization of Proposition \ref{prop:f85} descends, as in section \ref{sec:bb}, to the quotient where $v_{ij}$ depends only on $j-i\in\mathbb Z_2$, and then $v$ commutes with both copies. Thus we have a factorization through $A\otimes C(\mathbb Z_2)$, with $v_{ij}\to1\otimes\delta_{j-i}$, which is a Hopf algebra with the comultiplication coming from $w_{ia,jb}=u^{(i)}_{ab}v_{ij}$, namely
$$\Delta(u^{(i)}_{ab})=\sum_{k,c}u^{(i)}_{ac}v_{ik}\otimes u^{(k)}_{cb}$$
which is the crossed coproduct $A\rtimes C(\mathbb Z_2)$, with $\mathbb Z_2$ acting via $\nu$ \cite{bb3,bic}. In this picture, Proposition \ref{prop:f85b} tells us that the model $\theta$ of $F_{8,5b}$ is block diagonal with respect to $\mathbb C^8=(P_+\otimes\mathbb C^4)\oplus(P_-\otimes\mathbb C^4)$, with $1\otimes\delta_g$ mapped to the projection $E_g$ onto the $g$-th block, and with $A$ acting on the 2 blocks via the models
$$\sigma^0=(U^K,U^L),\qquad\sigma^1=(R^{KL},R^{LK})$$
the first matrix in each pair being the image of $u^{(0)}$, and the second one, that of $u^{(1)}$.

\begin{proposition}\label{prop:reduction}
Consider the model $\tau=(\sigma^1\otimes\sigma^1\nu)\Delta$ of $A$, acting on $\mathbb C^4\otimes\mathbb C^4$. If the Hopf image of $\sigma^0\oplus\tau$ is $S_4^+\,\hat{*}\,S_4^+$, and in particular if the Hopf image of $\tau$ is $S_4^+\,\hat{*}\,S_4^+$, then the quantum group of $F_{8,5b}$ is $S_4^+\wr_*\mathbb Z_2$.
\end{proposition}

\begin{proof}
Let $I\subset A\rtimes C(\mathbb Z_2)$ be the Hopf ideal defining the Hopf image of $\theta$. Since $\delta_0,\delta_1$ are central idempotents, we have $I=(J_0\otimes\delta_0)\oplus(J_1\otimes\delta_1)$, with $J_g=\{a\in A\,|\,a\otimes\delta_g\in I\}$, and since $\theta(1\otimes\delta_g)=E_g\neq0$, we have $1\notin J_g$. The above formula of $\Delta$ gives, for $a\in A$ and $h\in\mathbb Z_2$,
$$\Delta(a\otimes\delta_h)=\sum_{g\in\mathbb Z_2}(a_{(1)}\otimes\delta_g)\otimes(\nu^g(a_{(2)})\otimes\delta_{h-g})$$
and by multiplying by $(1\otimes\delta_g)\otimes(1\otimes\delta_{h-g})$, we see that $a\in J_h$ implies $\Delta(a)\in J_g\otimes A+A\otimes\nu^g(J_{h-g})$, for any $g\in\mathbb Z_2$.

We first prove that $J_0=0$ implies $I=0$. Let $a\in J_1$. With $g=1$ and $g=0$ we obtain $\Delta(a)\in J_1\otimes A$ and $\Delta(a)\in A\otimes J_1$, so that $\Delta(a)\in J_1\otimes J_1$. By applying $m(S\otimes id)$ we obtain $\varepsilon(a)1\in J_1$, so $\varepsilon(a)=0$, and then $a=(\varepsilon\otimes id)\Delta(a)\in\varepsilon(J_1)J_1=0$.

Now $J_0$ is the image of $I$ by the Hopf algebra map $A\rtimes C(\mathbb Z_2)\to A$ given by $a\otimes\delta_h\to\delta_{h0}a$, so it is a Hopf ideal of $A$. Since $\theta(a\otimes\delta_0)=\sigma^0(a)E_0$, we have $J_0\subset\ker\sigma^0$. Also, for $a\in J_0$ we have $(\theta\otimes\theta)\Delta(a\otimes\delta_0)=0$, and the restriction of this operator to the block $E_1\otimes E_1$ is $(\sigma^1\otimes\sigma^1\nu)\Delta(a)=\tau(a)$, so $J_0\subset\ker\tau$. Thus $J_0$ is contained in the Hopf ideal defining the Hopf image of $\sigma^0\oplus\tau$, which is $0$ by assumption, and so $I=0$.
\end{proof}

\noindent\textbf{The anti structure.} The pair $\sigma^1$ has a special form. With $\pi=R^{KL}$ we have $R^{LK}_{ab}=(R^{KL}_{ba})^t$, so that $\sigma^1=(\pi,\pi^t\circ S)$, where $\pi^t(a)=\pi(a)^t$, the composition of the 2 antimultiplicative maps $S$ and $x\to x^t$ being a representation. Since $u^t$ is magic, there is also a transposition automorphism $\iota$ of the algebra $C(S_4^+)$, given by $\iota(u_{ij})=u_{ji}$. For a Pauli fiber, since $\bar{c}_k=\pm c_k$, we have $\pi_x(u_{ba})^t=Proj(\bar{c}_b\bar{x}\bar{c}_a)=\pi_{\bar{x}}(u_{ba})$, so $\pi_x^t\circ S=\pi_{\bar{x}}\circ\iota$, these 2 representations agreeing on the generators. Thus, if $\pi=Ad(F)\circ\pi_x$, as in Proposition \ref{prop:rkl}, we obtain
$$\sigma^1\simeq(\pi_x,Ad(M)\circ\pi_{\bar{x}}\circ\iota),\qquad M=F^*\bar{F}$$
with $M$ a symmetric unitary. In terms of the representations $\theta_y$ of Lemma \ref{lem:pauli}, we have $\theta_y(\iota(a_{ij}))=y_{ji}\tau_j\otimes\tau_i$, so the second copy is a Pauli fiber with the 2 tensor legs exchanged, that is, twisted the opposite way. Numerically, this appears to be what produces freeness. With $x,y\in SU_2$ generic, and with the programs \texttt{anti.py}, which transposes the 2 magic indices, that is, uses $\iota$, \texttt{even4.py}, \texttt{wtest.py} and \texttt{wtest4.py}, the free values being the moments $2,6,22,90$ of $\pi_2$, and the values of the twist of $SO_3\times SO_3$, with independent copies, being $2,6,22,92$:

\medskip

\begin{center}
\begin{tabular}{lcccc}
\toprule
model of $A$ & $c_1$ & $c_2$ & $c_3$ & $c_4$\\
\midrule
$(\pi_x,\pi_y)$ & 2 & 6 & 22 & 92\\
$(\pi_x,\pi_y\circ\iota)$ & 2 & 6 & 22 & 90\\
$\sigma^1$, for $F_{8,5b}$ & 2 & 6 & 22 & 90\\
$\tau$, for $F_{8,5b}$ & 2 & 6 & 22 & 90\\
\bottomrule
\end{tabular}
\end{center}

\medskip

\noindent Indeed, for $(\pi_x,\pi_y)$ both copies are fibers of the same embedding of \cite[Theorem 5.1]{bb2}, at the point $(\rho(x),\rho(y))$ of the abelianization of $C(SO_3)*C(SO_3)$, so the Hopf image is contained in the twist of $SO_3\times SO_3$, and the value $c_4=92$ is the one of this twist. Note also that the value $c_4=90$ for $(\pi_x,Ad(W)\circ\pi_y)$ appears as well for the various unitaries $W$ tested, outside the Pauli normalizer, the 2 copies being then not fibers of the same embedding.

\medskip

\noindent\textbf{Placement models.} In order to understand the tensor powers of such pairs, consider a space $(\mathbb C^2)^{\otimes Q}\otimes K$, and for $c=0,1$ an ordered pair $(p_c,q_c)$ of distinct qubits, and a representation $X^c$ of $C(SO_3)$ on the tensor product of the other factors. The formula
$$a^{(c)}_{ij}\to\tau_i^{(p_c)}\tau_j^{(q_c)}\otimes X^c_{ij}$$
where $\tau^{(p)}$ denotes $\tau$ acting on the qubit $p$, defines then, via the embedding of \cite[Theorem 5.1]{bb2} followed by $X^c$, a model of $A$, that we call a placement model. As basic examples, $(\theta_y,\theta_z)$ is the placement model with $Q=2$, $K=\mathbb C$, both copies placed on $(1,2)$, and $X^0,X^1$ the evaluations at $y,z$, and $(\theta_y,\theta_z\circ\iota)$ is the one with copy 0 placed on $(1,2)$, copy 1 placed on $(2,1)$, and $X^1$ the evaluation at $z^t$.

\begin{lemma}\label{lem:placement}
The tensor product $(\alpha\otimes\beta)\Delta$ of 2 placement models $\alpha,\beta$, with data $(p_c,q_c,X^c)$ and $(p_c',q_c',Y^c)$, is a placement model, with copy $c$ placed on $(p_c,q_c')$, and given by
$$Z^c_{ij}=\sum_k(\tau_k\otimes\tau_k)^{(q_c,p_c')}\otimes X^c_{ik}\otimes Y^c_{kj}$$
In particular, $(\theta_y\otimes\theta_z)\Delta\simeq\oplus_{v\in V}\theta_{yvz}$, for any $y,z\in SO_3$.
\end{lemma}

\begin{proof}
Since $a$ is a corepresentation, $\Delta(a^{(c)}_{ij})=\sum_ka^{(c)}_{ik}\otimes a^{(c)}_{kj}$ is mapped to $\sum_k\tau_i^{(p_c)}\tau_k^{(q_c)}X^c_{ik}\otimes\tau_k^{(p_c')}\tau_j^{(q_c')}Y^c_{kj}=\tau_i^{(p_c)}\tau_j^{(q_c')}\otimes Z^c_{ij}$, because $X^c$ does not act on the qubit $q_c$, and $Y^c$ does not act on $p_c'$. The operators $\tau_k\otimes\tau_k$ commute, and in the Bell basis of $\mathbb C^2\otimes\mathbb C^2$ they are $diag(v_{kk})$, with $v$ running over $V$, the singlet $\Omega$ corresponding to $v=1$. Thus $Z^c=\oplus_{v\in V}X^cvY^c$, which is a commuting family, with values in $SO_3$, and so a representation of $C(SO_3)$. The last assertion is the case $Q=2$, $K=\mathbb C$, of points.
\end{proof}

\noindent For the pair $(\theta_y,\theta_z\circ\iota)$, the tensor square has copy 0 placed on the qubits $(1,4)$, contracted over $(2,3)$, and copy 1 placed on $(2,3)$, contracted over $(1,4)$. Thus each copy is placed on the qubits over which the other one is contracted, and there the operators $\tau_i\otimes\tau_j$ with $i\neq j$ of one copy do not commute with the operators $\tau_k\otimes\tau_k$ entering in the representation $Z^c$ of the other copy, which they permute, up to signs. For $(\theta_y,\theta_z)$, on the other hand, both copies remain on the same qubits at all levels, and the representations $Z^0,Z^1$ commute. This seems to be the mechanism behind the above table, as checked with \texttt{placement.py}, and leads to the question whether, for $x,y\in SU_2$ generic, the Hopf image of $(\pi_x,\pi_y\circ\iota)$ is $S_4^+\,\hat{*}\,S_4^+$, that is, whether a twisted point and an anti-twisted point of $SO_3^{-1}$ are free. Together with Proposition \ref{prop:reduction}, a positive answer would give the equality $G=S_4^+\wr_*\mathbb Z_2$ for $F_{8,5b}$, modulo the unitary $M$. Since the condition $c_k=c_k(S_4^+\,\hat{*}\,S_4^+)$ defines, at each level $k$, the complement of a real algebraic subset of $SU_2\times SU_2$, it would be enough to establish freeness at a single point $(x,y)$, as in the proofs of the freeness of 2 generic rotations. We have no such point, however, and the numerical evidence for this freeness stops at $k=4$.

\medskip

\subsection{The case of $F_{8,5a}$}

Here the same reordering as in Proposition \ref{prop:f85b}, with the rows $0,1,3,5$ and $2,4,6,7$, and the columns $0,1,2,3$ and $6,7,4,5$, gives $F_{8,5a}=\begin{pmatrix}K&K\\ L&-L\end{pmatrix}$, with the same $K$ as for $F_{8,5b}$, and with
$$L=\begin{pmatrix}1&-\bar{c}&-t&d\\ 1&-\bar{c}&t&-d\\ 1&\bar{c}&e&\frac{ae}{b}\\ 1&\bar{c}&-e&-\frac{ae}{b}\end{pmatrix},\qquad t=\frac{bd}{a}$$
so the magic matrix has the same form as in Proposition \ref{prop:f85b}, with $R^{LK}_{ab}=(R^{KL}_{ba})^t$. The block structures being nested, the proof of Proposition \ref{prop:wreath}, applied twice, shows that the quantum group of $F_{8,5a}$ is contained in $H_2^+\wr_*\mathbb Z_2$, which is the quantum automorphism group of the binary tree with 8 leaves \cite{dkr,dfk}. The main character of this latter quantum group has law $\frac{1}{2}\delta_0+\frac{1}{2}(\mu\boxplus\mu)$, with $\mu=\frac{1}{2}\delta_0+\frac{1}{2}\alpha$ being the law for $H_2^+$, where $\alpha$ is the arcsine law on $[0,4]$, by the proof of Proposition \ref{prop:f85}, and its moments are
$$c_1,\ldots,c_6=1,4,19,98,531,2974$$
which are exactly the values found for $F_{8,5a}$, at all the levels computed. Also, the block diagonal $8\times8$ magic matrix formed by the 2 inner copies has $c_1,\ldots,c_4=2,8,38,196$, the moments of $\mu\boxplus\mu$, while 2 independent copies would give $c_4=204$, as computed by the programs \texttt{blocks.py} and \texttt{freeconv.py}. As a first algebraic step, we have:

\begin{proposition}\label{prop:f85a}
Assume that $c$ or $b/a$ is not a root of unity. Then the $4\times4$ magic matrix of $F_{8,5a}$ obtained by summing over the pairs generates $H_2^+$. Also, the 2 middle level projections $P_0,P_1$ of the inner copies, acting on $P_-\otimes\mathbb C^4$, produce a representation of $C^*(D_\infty)=C(\mathbb Z_2)*C(\mathbb Z_2)$ which is faithful on the group $D_\infty$, and so inner faithful.
\end{proposition}

\begin{proof}
The pair level magic matrix preserves the partition into 2 blocks of 2 pairs, so it factorizes through $\mathbb Z_2\wr_*\mathbb Z_2=H_2^+$, which is the crossed coproduct $C^*(D_\infty)\rtimes C(\mathbb Z_2)$ of section \ref{sec:bb}, with $X=Y=\mathbb Z_2$, the 2 inner involutions $g_0,g_1$ generating $D_\infty$, and the block projection $E$ representing $C(\mathbb Z_2)$. By Proposition \ref{prop:nq}, or rather by its proof, which applies verbatim here, the Hopf image is $C^*(D_\infty/N)\rtimes_\Phi C(\mathbb Z_2)$, with $N$ consisting of the elements acting by scalars on the 2 blocks. A reflection $g_0(g_0g_1)^n$ acts on the block $E$ as $g_0$, which is not scalar there, because $E\cdot\frac{1+g_0}{2}$ has rank 2, inside a block of rank 4. As for the translation $g_0g_1$, a direct computation from the formula of $L$ shows that it acts trivially on the block $E$, and with eigenvalues $c^{\pm2},(b/a)^{\pm2}$ on the block $1-E$. Thus $(g_0g_1)^m$ acts by scalars only if $c^{4m}=(b/a)^{4m}=1$, so $N=\{1\}$, which gives the first assertion. As for the second one, with $K_0=1$ we have $P_0=Proj(\bar{L}_0)+Proj(\bar{L}_1)$ and $P_1=Proj(L_0)+Proj(L_0\bar{K}_2)$, and since $L_0\pm L_1$ and $L_0\pm L_0\bar{K}_2$ are supported on the first 2 and on the last 2 coordinates, we obtain
$$P_0=Proj(1,-c)\oplus Proj(-\bar{t},\bar{d}),\qquad P_1=Proj(1,-\bar{c})\oplus Proj(-t,d)$$
Thus, with $S_i=2P_i-1$, the unitary $S_0S_1$ has eigenvalues $c^{\pm2},(b/a)^{\pm2}$. Since $P_0,P_1$ are group-like, in the sense that $2P_i-1$ are group-like elements, the Hopf image of this representation of $C^*(D_\infty)$ is the group algebra of the image of $D_\infty$, and this image is $D_\infty$ itself as soon as $S_0S_1$ has infinite order. Note that the representation of the algebra $C^*(D_\infty)$ itself is not faithful, being finite dimensional: with $u=g_0g_1$ and $E=\{c^{\pm2},(b/a)^{\pm2}\}$, the element $\prod_{\lambda\in E}(u-\lambda)\neq0$ is in its kernel.
\end{proof}

\noindent Proposition \ref{prop:reduction} applies verbatim to $F_{8,5a}$, with $H_2^+$ instead of $S_4^+$. Here the model $\theta$ splits as $\varepsilon\oplus\sigma\oplus\pi_2$ on the block $E_0$, with $\varepsilon$ the counit, $\sigma$ the classical permutation $(03)(15)(24)(67)$, which flips all the pairs, and $\pi_2$ irreducible of dimension 2, and as an irreducible representation of dimension 4 on the block $E_1$. Numerically, at random parameters, each of the 2 copies of $C(H_2^+)$ has $c_1,\ldots,c_5=1,3,10,35,126$ in $\sigma^0\oplus\tau$, the values of $H_2^+$, and the model $\sigma^0\oplus\tau$ itself has $c_1,\ldots,c_4=2,8,38,196$, the moments of $\mu\boxplus\mu$, which are the values of $H_2^+\,\hat{*}\,H_2^+$, while $\sigma^0$ alone has $c_1,\ldots,c_4=2,10,56,338$. Thus, here too, the question is the freeness of the 2 copies.

As an algebra, $C(H_2^+)*C(H_2^+)=C^*(\Lambda)$, with $\Lambda=(D_\infty\times\mathbb Z_2)*(D_\infty\times\mathbb Z_2)$, its Haar state being the canonical trace of $\Lambda$, and there are 2 natural families of intermediate quantum groups. The first one consists of the dihedral quotients at the middle level, which are the quantum group analogues of the quotients of $A_2*A_2$ studied by Izumi, Morrison and Penneys \cite{imp}, and these are excluded by Proposition \ref{prop:f85a}. The second one consists of the free wreath products with amalgamation of Freslon \cite{fre}, obtained by identifying the 2 copies of a Hopf subalgebra $C(K)\subset C(H_2^+)$, and these are excluded too:

\begin{proposition}\label{prop:amalg}
Assume that $c$ or $b/a$ is not a root of unity. Then the quantum group of $F_{8,5a}$ is not contained in the free wreath product with amalgamation $H_2^+\wr_{*,K}\mathbb Z_2$, for any Hopf subalgebra $C(K)\neq\mathbb C$ of $C(H_2^+)$.
\end{proposition}

\begin{proof}
We have $H_2^+=O_2^{-1}$ \cite{bbc}, a twist of $O_2$, so the irreducible corepresentations of $H_2^+$ are, as for $O_2$, the trivial one, a 1-dimensional one $\delta$, and 2-dimensional ones $\rho_k$ with $k\geq1$, with $\rho_k\otimes\rho_l=\rho_{k+l}\oplus\rho_{|k-l|}$, where $\rho_0=1\oplus\delta$. A Hopf subalgebra $C(K)\neq\mathbb C$ corresponds to a full subcategory, stable by tensor products and conjugation, and containing a nontrivial irreducible corepresentation, so it contains $\delta$, because $\rho_k\otimes\rho_k\supset\rho_0$. Now $\delta$ is the unique nontrivial group-like element of $C(H_2^+)$, namely $m=2(u_{00}+u_{01})-1$, with $\{0,1\},\{2,3\}$ being the 2 pairs, and this element is invariant under the automorphisms of $H_2^+$, so the amalgamation forces the relation $\nu^0(m)=\nu^1(m)$, for any identification of the 2 copies. But in the model, $\nu^0(m)$ and $\nu^1(m)$ are mapped, on the block $P_-\otimes\mathbb C^4$, to the symmetries $S_0,S_1$ from the proof of Proposition \ref{prop:f85a}, and since $S_0S_1$ has eigenvalues $c^{\pm2},(b/a)^{\pm2}$, not all equal to 1, we have $S_0\neq S_1$.
\end{proof}

\noindent In subfactor terms, by Tarrago and Wahl \cite{twa}, the equality $G=H_2^+\wr_*\mathbb Z_2$ means that the planar algebra of $F_{8,5a}$ is the Bisch--Jones free composition \cite{bjo} of 3 subfactors of index 2, and the above results exclude the known intermediate objects between the tensor and the free compositions. We have not been able, however, to exclude the other intermediate quantum groups, such as the generalized free wreath products of Fima and Troupel \cite{ftr}, and the method of Proposition \ref{prop:reduction} does not go further. Indeed, below the top level, the Hopf subalgebras $C(\mathbb Z_2)$ of the 2 copies of $C(H_2^+)$ generate $C^*(D_\infty)$, which is not central, and the quantum subgroups which can be read on the model are small: for instance, the largest invariant subspaces of $(\theta\otimes\theta)\Delta$ and of its triple analogue on which the pair level magic matrix acts as the identity have dimensions 10 and 54, and on them the 4 symmetries inside the pairs generate only a virtually abelian group, as found by the programs \texttt{pairfixed.py}, \texttt{flipgroup.py} and \texttt{imagealg.py}. Thus, what remains is a genuine freeness question, which might well require analytic methods. Summarizing, we have proved the following result:

\begin{theorem}\label{thm:f85b}
For $s=bd\neq\pm1$ and $w=\bar{c}\bar{e}\notin\{\pm1,\pm i\}$, the quantum group $G$ of $F_{8,5b}$ satisfies:
\begin{enumerate}
\item $G\subset S_4^+\wr_*\mathbb Z_2$, and each of the 2 inner copies generates $S_4^+$.

\item The model splits as $\sigma^0\oplus\sigma^1$, with $\sigma^1=(\pi,\pi^t\circ S)$, and $\pi$ being a Pauli fiber, up to conjugation.

\item If the model $\sigma^0\oplus(\sigma^1\otimes\sigma^1\nu)\Delta$ generates $S_4^+\,\hat{*}\,S_4^+$, then $G=S_4^+\wr_*\mathbb Z_2$.
\end{enumerate}
\end{theorem}

\begin{proof}
This follows from Propositions \ref{prop:f85}, \ref{prop:rkl} and \ref{prop:reduction}, and from the above discussion of the anti structure.
\end{proof}

Similarly, by Propositions \ref{prop:f85a} and \ref{prop:amalg}, and by the above discussion of the case of $F_{8,5a}$, the quantum group $G'$ of $F_{8,5a}$ is contained in the quantum automorphism group $H_2^+\wr_*\mathbb Z_2$ of the binary tree, and when $c$ or $b/a$ is not a root of unity, $G'$ projects onto $H_2^+$ at the pair level, is contained in no free wreath product with amalgamation $H_2^+\wr_{*,K}\mathbb Z_2$ with $K\neq\{1\}$, and is equal to $H_2^+\wr_*\mathbb Z_2$ as soon as the corresponding model $\sigma^0\oplus(\sigma^1\otimes\sigma^1\nu)\Delta$ generates $H_2^+\,\hat{*}\,H_2^+$.

\section{Real Hadamard matrices}\label{sec:real}

We discuss here the real Hadamard matrices. The numerics show that the matrices with $c_2=2$, including the Paley matrices up to order 108, are Temperley--Lieb as far as computed, while those of order 16, which are all of Di\c t\u a type, have sequences $c_k$ which decrease from one class to the next. Then we turn to infinite depth, which is the first step towards wildness, and which we establish for the Paley II matrices, and more generally for the matrices coming from symmetric conference matrices, and for the Paley I matrices, with $11\leq q<10^5$.

\subsection{Constructions}

Let $\chi$ be the quadratic character of $\mathbb F_q$, and $Q_{ab}=\chi(b-a)$. The Paley I matrix, for $q\equiv3$ mod 4, is $P^1_{q+1}=1+\begin{pmatrix}0&\mathbb I\\-\mathbb I&Q\end{pmatrix}$, and the Paley II matrix, for $q\equiv1$ mod 4, is obtained from $\begin{pmatrix}0&\mathbb I\\ \mathbb I&Q\end{pmatrix}$ by replacing $0\to\begin{pmatrix}1&-1\\-1&-1\end{pmatrix}$ and $\pm1\to\pm\begin{pmatrix}1&1\\1&-1\end{pmatrix}$, as in \cite[Theorem 1.19]{banicaHM}, see also \cite{pal}. Note that $P^1_{q+1}$ is skew, in the sense that $P^1_{q+1}+(P^1_{q+1})^t=2$, because $\chi(-1)=-1$ for $q\equiv3$ mod 4, while $P^2_{2(q+1)}$ is symmetric. The elements of $\mathbb F_q$ are ordered as $0,1,\ldots,q-1$ for $q$ prime, and as polynomials $c_0+c_1x+c_2x^2$ in lexicographic order of $(c_0,c_1,c_2)$ for $q=9,25,27,49$, these matrices being built by the module \texttt{pmatrices.py}. The Williamson matrices are built from symmetric circulant matrices $A,B,C,D\in M_K(\pm1)$ with first entry $+1$ and $A^2+B^2+C^2+D^2=4K$, placed as in \cite[Theorem 1.24]{banicaHM}, the quadruple used being the first one found by a lexicographic search. Finally, $\mathrm{had}_{N.k}$ denotes the $k$-th matrix of order $N$ in McKay's library \cite{mck}, in his order, starting from $k=0$.

\medskip

\noindent\textbf{Classes.} The number of equivalence classes is 1, 5, 3, 60, 487 at the orders 12, 16, 20, 24, 28. At order 12, the Paley matrices $P_{12}^1,P_{12}^2$ and the Williamson matrix $W_{12}$, with $K=3$, are all equivalent, to the following matrix:

{\normalsize\setlength{\arraycolsep}{3.5pt}$$P_{12}^1=\left(\begin{array}{cccccccccccc}+ & + & + & + & + & + & + & + & + & + & + & +\\ - & + & + & - & + & + & + & - & - & - & + & -\\ - & - & + & + & - & + & + & + & - & - & - & +\\ - & + & - & + & + & - & + & + & + & - & - & -\\ - & - & + & - & + & + & - & + & + & + & - & -\\ - & - & - & + & - & + & + & - & + & + & + & -\\ - & - & - & - & + & - & + & + & - & + & + & +\\ - & + & - & - & - & + & - & + & + & - & + & +\\ - & + & + & - & - & - & + & - & + & + & - & +\\ - & + & + & + & - & - & - & + & - & + & + & -\\ - & - & + & + & + & - & - & - & + & - & + & +\\ - & + & - & + & + & + & - & - & - & + & - & +\end{array}\right)$$}

\noindent At order 16, the 5 classes are $\mathrm{had}_{16.0}=W_{16}$, the Walsh matrix, and $\mathrm{had}_{16.1},\ldots,\mathrm{had}_{16.4}$. At order 20, the 3 classes are $P_{20}^1\sim\mathrm{had}_{20.2}$ with $q=19$, $P_{20}^2\sim W_{20}\sim\mathrm{had}_{20.0}$ with $q=9$ and $K=5$, and $\mathrm{had}_{20.1}$. At orders 24 and 28 we use first the Paley matrices $P_{24}^1$, with $q=23$, $P_{28}^1$, with $q=27$, and $P_{28}^2$, with $q=13$, and the Williamson matrix $W_{28}$, with $K=7$, which is equivalent to $P_{28}^2$, and then all the classes of \cite{mck}. The matrices of order $\geq16$ are not displayed, being given by the above constructions, and by \cite{mck}.

\medskip

\noindent\textbf{Transposition.} With an exact equivalence test, by row by row backtracking, which is complete, and with the equivalences found being checked, as done by the programs \texttt{equiv.py} and \texttt{preflight\_real.py}, all the above matrices are equivalent to their transposes, except for $\mathrm{had}_{16.3}$ and $\mathrm{had}_{16.4}$, which are transposes of each other, $\mathrm{had}_{16.3}^t\sim\mathrm{had}_{16.4}$. Thus $\mathrm{had}_{16.3},\mathrm{had}_{16.4}$ is the only pair to be computed separately.

\subsection{Results}\label{sec:realres}

For $H$ and $H^t$ alike, we find the following values, the first ones being computed by the programs \texttt{stage1.py} and \texttt{stage1b.py}, and the higher ones by the programs mentioned below:

\medskip

\begin{center}
\begin{tabular}{lll}
\toprule
matrix & $c_1,\ldots,c_4$ & higher levels\\
\midrule
$P_{12}^1$ & $1,2,5,14$ & $c_5=42$, $c_6=132$\\
$\mathrm{had}_{16.0}=W_{16}$ & $1,16,256,4096$ & \\
$\mathrm{had}_{16.1}$ & $1,7,58,568$ & $c_5=6496$, finite, of dimension 256\\
$\mathrm{had}_{16.2}$ & $1,4,21,142$ & \\
$\mathrm{had}_{16.3}$ & $1,3,11,46$ & $c_5=212$\\
$\mathrm{had}_{16.4}\sim\mathrm{had}_{16.3}^t$ & $1,3,11,46$ & $c_5=212$\\
$P_{20}^1$, $P_{20}^2$ & $1,2,5,14$ & $c_5=42$\\
$\mathrm{had}_{20.1}$ & $1,2,5,14$ & \\
$P_{24}^1$ & $1,2,5,14$ & $c_5=42$\\
$P_{28}^1$, $P_{28}^2$ & $1,2,5,14$ & \\
$P_{32}^1$, $P_{36}^2$, $P_{44}^1$, $P_{48}^1$ & $1,2,5,14$ & \\
$P_{52}^2$, $P_{100}^2$, $P_{104}^1$, $P_{108}^1$, $P_{108}^2$ & $1,2,5$ & \\
\bottomrule
\end{tabular}
\end{center}

\medskip

\noindent Here the values $c_5=42$ at orders 20 and 24, and all the values at the orders $N\geq32$, come from the Jones--Wenzl count of section \ref{sec:jw}, done in real arithmetic by the programs \texttt{bobcat.py} and \texttt{cormier.py}, which give $e_k=0$ in all cases. Level 5 at orders 20 and 24, in dimensions $3.2\cdot10^6$ and $8\cdot10^6$, takes 5 to 20 minutes, level 4 at order 48, in dimension $5.3\cdot10^6$, takes about 35 minutes, and level 3 at orders 100 to 108, in dimension about $1.2\cdot10^6$, takes about 2 minutes. For $P_{20}^1$, the value $c_5=42$ was obtained first by the program \texttt{puma.py}, with noncrossing seeds, in 75 minutes, and at orders 20 and 24 the converged norms computed by the program \texttt{sandcat.py} at level 5, which are $0.2919$, $0.5084$, $0.2396$ for $P_{20}^1$, $P_{20}^2$, $P_{24}^1$, confirm these values. Thus the matrices of orders 12, 20, 24, 28, 32, 36, 44, 48, 52, 100, 104, 108 tested are Temperley--Lieb as far as computed, while at order 16 none of them is, the moments of the 5 classes decreasing from $W_{16}$ to $\mathrm{had}_{16.3}\sim\mathrm{had}_{16.4}^t$, which is a comparison of moments, the corresponding inclusions of quantum groups being not proved, and with only $W_{16}$ having a commuting magic matrix. The Di\c t\u a block structure search shows that $\mathrm{had}_{16.1},\ldots,\mathrm{had}_{16.4}$ all have a Di\c t\u a block structure, in the general sense of section 1, while the other matrices above have no such structure. Among them, only $\mathrm{had}_{16.1}$ is a deformed Walsh matrix $F_{\mathbb Z_2^a}\otimes_QF_{\mathbb Z_2^b}$, as explained below.

\subsection{The matrices $\mathrm{had}_{16}$}\label{sec:had16}

The sequence $1,3,11,46$ of $\mathrm{had}_{16.3}$ and $\mathrm{had}_{16.4}$ starts as the little Schr\"oder sequence $1,3,11,45$ of $F_{8,5b}$, see section \ref{sec:cumul}, but leaves it at level 4, where 45 of the fixed vectors come from level 3, and 1 is new. The Hankel conditions give, with $c_1,\ldots,c_4=1,3,11,46$:
$$\det\begin{pmatrix}1&3&11\\3&11&46\\11&46&c_5\end{pmatrix}=2c_5-411\geq0$$
so $c_5\geq206$. For $\mathrm{had}_{16.3}$, the seeds give 207 fixed vectors, a search with random blocks on their orthogonal complement gives 5 more vectors, and a Lanczos computation of the top eigenvalues of $T_5$ on the orthogonal complement of these 212 vectors gives $0.875,0.8536,0.8248,0.8042$, the whole computation being redone, with a new random seed, by the program \texttt{margay.py}. Thus $c_5=212$, with $2c_5-411=13$, and the value $0.875$ is the same as the second eigenvalue at level 4. Since $\mathrm{had}_{16.4}\sim\mathrm{had}_{16.3}^t$, Theorem \ref{thm:transp} gives $c_5=212$ for $\mathrm{had}_{16.4}$ as well, and the direct computation by the program \texttt{ocelot.py}, which takes 40 minutes, gives indeed 212, with 207 vectors coming from the seeds. At the next level, the positivity of the Hankel matrix $(c_{i+j})_{i,j=0,\ldots,3}$, with $c_0=1$, amounts to $10c_6-10377\geq0$, and so gives $c_6\geq1038$, in agreement with the prediction $c_6\in\{1054,1055\}$ coming from the Boolean cumulants, discussed below.

\medskip

\noindent\textbf{OEIS.} The search for $1,3,11,46$ in the OEIS database returns 34 sequences, whose next terms are as follows, with the ones below 206 being excluded by the above Hankel bound:

\medskip

\begin{center}
\begin{tabular}{lc>{\raggedright\arraybackslash}p{7.5cm}}
\toprule
next term & number & sequences\\
\midrule
180--205 & 10 & excluded by $c_5\geq206$\\
206 & 2 & A271959, A151140\\
207 & 7 & A006605, A155134, A275757, A291090, A151141, A357233, A306822\\
208 & 2 & A193074, A287891\\
209 & 2 & A371428, A233389\\
210 & 3 & A281548, A086521, A225293\\
211 & 2 & A046996, A248426\\
213 & 2 & A337913, A129579\\
219, 230, 233, 234 & 4 & A393791, A030814, A030976, A264416\\
\bottomrule
\end{tabular}
\end{center}

\medskip

\noindent Thus $1,3,11,46,212$ is not in OEIS. The seeds alone give 207, the value of A006605, the number of modes of connections of $2n$ points, related to Baxter's generalization of the Temperley--Lieb operators. Since the candidates are consecutive integers, approximate values of $c_5$ would have been useless here, and the probabilistic count, which produces an integer, is the right tool.

\medskip

\noindent\textbf{The matrix $\mathrm{had}_{16.1}$.} An exhaustive search over the real deformed Fourier matrices $F_{\mathbb Z_2^a}\otimes_QF_{\mathbb Z_2^b}$ with $a+b=4$ and $Q$ having $\pm1$ entries, by using the counting of Proposition \ref{prop:nq}, done by the program \texttt{walshdita.py}, shows that only 2 classes appear, namely $W_{16}$, at $Q=1$, and $\mathrm{had}_{16.1}$, which is equivalent to $F_{\mathbb Z_2^2}\otimes_QF_{\mathbb Z_2^2}$, with $Q$ having a single entry $-1$, at the position $(3,3)$, the elements of $\mathbb Z_2^2$ being labeled $0,1,2,3$ in binary. Thus Theorem \ref{thm:cyclic} (2) applies, and shows that the quantum group of $\mathrm{had}_{16.1}$ is finite, of dimension 256, with $|\Gamma_{X,Y}/N_Q|=64$, and with
$$c_1,\ldots,c_7=1,7,58,568,6496,84352,1197568$$
which confirms the value $c_4=568$ found above, by the probabilistic count. The other 3 classes, $\mathrm{had}_{16.2},\mathrm{had}_{16.3},\mathrm{had}_{16.4}$, are thus of Di\c t\u a type only in the general sense $H\otimes_Q(K_1,\ldots,K_M)$ of section 1, and are not deformed Fourier matrices.

\medskip

\noindent\textbf{Cumulants.} For the order 16 matrices, with the notations of section \ref{sec:cumul}, we have:

\medskip

\begin{center}
\begin{tabular}{lll}
\toprule
matrix & $\kappa_1,\kappa_2,\ldots$ & $b_1,b_2,\ldots$\\
\midrule
$W_{16}$ & $1,15,210,2715$ & $1,15,225,3375$\\
$\mathrm{had}_{16.1}$ & $1,6,39,303$ & $1,6,45,423$\\
$\mathrm{had}_{16.2}$ & $1,3,11,61$ & $1,3,14,95$\\
$\mathrm{had}_{16.3}$, $\mathrm{had}_{16.4}$ & $1,2,4,9,26$ & $1,2,6,23,103$\\
\bottomrule
\end{tabular}
\end{center}

\medskip

\noindent The Boolean cumulants $1,2,6,23,103$ of $\mathrm{had}_{16.3}$ are the numbers of permutations avoiding a single pattern of length 4, and the passage from $F_{8,5b}$, with Boolean cumulants $1,2,6,22,90$ counting the permutations avoiding the 2 patterns $2413,3142$, amounts to dropping 1 of the 2 forbidden patterns, matching the single new fixed vector at level 4. The OEIS searches for the free cumulants $1,2,4,9,26$, for the moments $1,0,2,4,17,66$ of $\chi-1$, and for the binomial transform $1,2,6,24,113,588$, give nothing relevant. This suggests exploring the Boolean cumulants of compact quantum groups, as opposed to the Boolean quantum semigroups of Liu \cite{liu}. Since the 3 classes of patterns of length 4 give $|\mathrm{Av}_6(\tau)|=512$ or $513$, this Boolean picture predicts $c_6\in\{1054,1055\}$ for $\mathrm{had}_{16.3}$.

\subsection{All classes of orders 24 and 28}\label{sec:orders}

We computed $c_2,c_3$ for all the 60 classes of order 24, and $c_2$ for all the 487 classes of order 28, from the data of \cite{mck}, for both $H$ and $H^t$, which agree in all cases, with the programs \texttt{sweep\_real.py} and \texttt{sweep\_c2.py}. Then, for the 45 classes of order 24 with $c_2=2$, we computed $c_4$, for both $H$ and $H^t$, and for the 487 classes of order 28 we computed $c_3$, for both $H$ and $H^t$, and then $c_4$ for $H$, for the first 9 classes, these 2 overnight computations being done by the programs \texttt{lynx.py} and \texttt{fedor.py}. The results are as follows:

\medskip

\begin{center}
\begin{tabular}{llll}
\toprule
order & classes & $c_1,c_2,\ldots$ & Di\c t\u a structure\\
\midrule
24 & 45 & $1,2,5,14$ & no\\
24 & 14 & $1,3,11$ & yes, of size $2\times12$\\
24 & 1 & $1,4,20$ & yes, of size $2\times12$\\
28 & 487 & $1,2,5$, and $c_4=14$ for 9 of them & no\\
\bottomrule
\end{tabular}
\end{center}

\medskip

\noindent For the classes with $c_2=2$, these values are confirmed via Theorem \ref{thm:jw}, by computing the norms $\|\Phi_k\|$ of section \ref{sec:jw} with a converged Lanczos iteration, with residuals below $10^{-9}$, by the program \texttt{orders.py}, built on the module \texttt{gapcheck.py}: the largest values are $\|\Phi_4\|=0.7685$ at order 24, and $\|\Phi_3\|=0.5702$ at order 28, and $\|\Phi_4\|=0.3189$ for the 9 classes of order 28 computed at level 4. Thus, for all these 547 matrices, we have $c_2\geq3$ precisely for the matrices of Di\c t\u a type, and the other ones are Temperley--Lieb up to level 4 at order 24, and up to level 3 at order 28, and level 4 for the 9 classes computed. Observe that 14 of the Di\c t\u a classes of order 24 start, like $F_{8,5b}$ and $\mathrm{had}_{16.3}$, with $1,3,11$.

\subsection{Infinite depth}\label{sec:infdepth}

All the real matrices above with $c_2=2$ are Temperley--Lieb as far as computed, and the natural question is whether they are wild. A first step is to prove that $G$ is infinite, that is, that the subfactor has infinite depth, and this was done by Montgomery for the Paley II matrices and for $P_7$ \cite{mo2,mo3}, by using the angle operator. We use here instead the following integrality principle:

\begin{theorem}\label{thm:integral}
If $G_H$ is finite, there is a monic polynomial $p\in\mathbb Z[x]$ such that $p(NT_k)=0$, for any $k$. Thus the eigenvalues of all the matrices $NT_k$ are algebraic integers, from a finite set not depending on $k$, and if some $NT_k$ has an eigenvalue which is not an algebraic integer, then $G_H$ is infinite, and the subfactor has infinite depth.
\end{theorem}

\begin{proof}
Let $\pi:C(G)\to M_N(\mathbb C)$ be the factorizing representation, $u_{ij}\to P_{ij}$, and $\omega=Tr\circ\pi$, so that $NT_k=(id\otimes\omega)(u^{\otimes k})$, as in the proof of Theorem \ref{thm:jw}. The dual $C(G)^*$ is an algebra with the convolution $\varphi*\phi=(\varphi\otimes\phi)\Delta$, with unit $\varepsilon$, and since $U=u^{\otimes k}$ satisfies $(id\otimes\Delta)U=U_{12}U_{13}$, the map $\varphi\to(id\otimes\varphi)(U)$ is a unital algebra morphism $C(G)^*\to M_{N^k}(\mathbb C)$. Now assume that $C(G)=\bigoplus_\rho M_{d_\rho}(\mathbb C)$ is finite dimensional, and let $R\subset C(G)^*$ be the $\mathbb Z$-span of the traces $Tr_\rho=Tr\circ\rho$ of its irreducible representations. We have $\varepsilon=Tr_\varepsilon$, and $Tr_\rho*Tr_\sigma=Tr\circ(\rho\otimes\sigma)\Delta$ is a sum of traces $Tr_\tau$, with integer multiplicities, because $(\rho\otimes\sigma)\Delta$ is a representation. Thus $R$ is a unital ring, which is a finitely generated $\mathbb Z$-module, so all its elements are integral over $\mathbb Z$. Since $\omega=\sum_\rho m_\rho Tr_\rho$, with $m_\rho$ the multiplicities of the irreducible representations in $\pi$, we have $\omega\in R$, so $p(\omega)=0$ for some monic $p\in\mathbb Z[x]$, and applying the above morphism gives $p(NT_k)=0$. Finally, the vertices of the principal graph being the irreducible corepresentations of $G$, finite depth means that $G$ has finitely many of them, that is, that $G$ is finite.
\end{proof}

For instance, for $\mathrm{had}_{16.1}$, whose quantum group is finite, of dimension 256, the spectra of $NT_2$ and $NT_3$ are $\{16,8,0\}$ and $\{16,8,4,0\}$, and for $F_4^s$, whose quantum group is finite precisely when $s$ is a root of unity \cite{bic}, the spectrum of $NT_2$ is $\{4,0\}$ at $s=i$, and $\{4,3,1,0\}$ at $s=e(1/6)$, while at a generic point it contains the eigenvalues $2.0402$ and $1.9598$. More generally, for an analytic family $(H_q)$, if some function $q\to Tr((NT_k(H_q))^m)$ is not constant, then $G_{H_q}$ is infinite outside a countable union of proper analytic subsets, the algebraic numbers being countable, and this applies to $H_6^q$, $BN_6^q$, $P_7^q$, whose spectra vary, by the tables in section \ref{sec:jw}. Also, the eigenvalues of the angle operator found by Montgomery \cite{mo2,mo3} appear here as eigenvalues of $NT_2$, namely $16q/(q+1)^2$ for the Paley II matrix of order $2(q+1)$, which we checked at $q=5,9,13$, and $1/49$ for $P_7$. Since these are not algebraic integers, Theorem \ref{thm:integral} recovers the infinite depth results of \cite{mo2,mo3} in these cases. These checks are done by the program \texttt{certify.py}, which also certifies the rational eigenvalues exactly, in integer arithmetic.

Taking traces in Theorem \ref{thm:integral}, if $G_H$ is finite then the numbers $Tr((NT_k)^r)=N^rTr(T_k^r)$ are algebraic integers, for any $k,r$, and so they are integers when $H$ is real. In vertex model terms, since $N^{-k}Tr(T_k^r)$ is the partition function on the $k\times r$ torus, as explained after Theorem \ref{thm:transp}, finite depth forces these partition functions, suitably normalized, to be integers. This gives a quick test, implemented in the program \texttt{moments.py}, and for instance for the Paley I matrices below at $q=11,19,23,31,43$ we have $Tr((NT_2)^2)=1664/3,277696/125,9088/3,4868,13688704/1331$, so the test detects infinite depth at $q=11,19,23,43$, but not at $q=31$, where $N=32$ has no odd prime factor. In general, the denominators come from the odd primes dividing $q+1$.

\subsection{Paley II matrices}\label{sec:paley2}

For a real Hadamard matrix $H\in M_N(\pm1)$ consider the linear map $\Phi_H$ on $M_N(\mathbb R)\simeq\mathbb R^N\otimes\mathbb R^N$ given by $\Phi_H(Y)_{xy}=\sum_{ab}H_{ax}H_{ay}H_{bx}H_{by}Y_{ab}$. Since $(T_2)_{i_1i_2,j_1j_2}=N^{-3}(\sum_xH_{i_1x}H_{i_2x}H_{j_1x}H_{j_2x})^2$, expanding the square gives
$$N^3T_2=\Phi_{H^t}\Phi_H$$
and we will use this formula for computing eigenvectors. Let $C\in M_n(0,\pm1)$ be a symmetric conference matrix, that is, $C=C^t$, $C_{xx}=0$, $C_{xy}=\pm1$ for $x\neq y$, and $C^2=qI$, with $n=q+1$, so that $q\equiv1$ mod 4. The matrix
$$H=C\otimes\begin{pmatrix}1&1\\1&-1\end{pmatrix}+1_n\otimes\begin{pmatrix}1&-1\\-1&-1\end{pmatrix}$$
is then a symmetric Hadamard matrix of order $N=2n$, and for the Paley conference matrix, built from the quadratic character $\chi$ on $\{\infty\}\cup\mathbb F_q$, this is the Paley II matrix $P^2_{2(q+1)}$ \cite{pal}. We index $H$ by pairs $(x,a)$, with $x\in\{1,\ldots,n\}$ and $a\in\{0,1\}$, so that $H_{(x,a),(z,c)}=(-1)^{ac}(C_{xz}+\delta_{xz}(-1)^{a+c})$, and we consider the following 4 matrices in $M_N(0,1)$, indexed by pairs $((x,a),(y,b))$:
$$P_1=[x\neq y,a=b]\quad,\quad P_2=[x\neq y,a\neq b]\quad,\quad P_3=[x=y,a\neq b]\quad,\quad P_4=[x=y,a=b]$$
Thus $P_4=1_N$, and $P_1+P_2+P_3+P_4$ is the all-one matrix. We set $X=P_1-P_2$ and $Y=P_1+P_2$, so that $X((x,a),(y,b))=(1-\delta_{xy})(-1)^{a+b}$.

\begin{theorem}\label{thm:conference}
With $\Phi=\Phi_H$, the space $span(X,Y,P_3,P_4)$ is invariant under $\Phi$, with
\begin{eqnarray*}
\Phi(X)&=&2(q+1)X-2(q-3)Y+4q(q-3)P_3\\
\Phi(Y)&=&-2(q-3)X-2(q+1)Y+N(N-2)P_4\\
\Phi(P_3)&=&2(q-3)X-2(q+1)P_3+2(q+1)P_4\\
\Phi(P_4)&=&N(Y+P_3+P_4)
\end{eqnarray*}
and $N^3T_2X=\beta X$, with $\beta=4(q+1)^2+4(2q+1)(q-3)^2$. Thus $NT_2$ has the eigenvalue
$$\frac{\beta}{N^2}=1+\frac{(2q+1)(q-3)^2}{(q+1)^2}$$
which is not an algebraic integer for $q\geq5$, and $H$ has an infinite quantum group, and produces an infinite depth subfactor. In particular, this holds for all the Paley II matrices.
\end{theorem}

\begin{proof}
Since $H$ is symmetric, $N^3T_2=\Phi^2$. Given $u=(z,c)$ and $v=(z',c')$, set $r_i=H_{iu}H_{iv}$, so that $\Phi(Z)_{uv}=\sum_{ij}r_ir_jZ_{ij}$, and $\sigma=(-1)^c$, $\tau=(-1)^{c'}$. For the all-one matrix we have $\Phi(J)_{uv}=(\sum_ir_i)^2=N^2\delta_{uv}$, and $\Phi(P_4)_{uv}=\sum_ir_i^2=N$. Next, $\Phi(P_3)_{uv}=2\sum_xr_{(x,0)}r_{(x,1)}$, and
$$r_{(x,0)}r_{(x,1)}=\sigma\tau(C_{xz}^2-\delta_{xz})(C_{xz'}^2-\delta_{xz'})=\sigma\tau(1-2\delta_{xz})(1-2\delta_{xz'})$$
which sums to $\sigma\tau(q-3)$ for $z\neq z'$, and to $\sigma\tau(q+1)$ for $z=z'$, giving the formula of $\Phi(P_3)$. As for $X$, with $\varepsilon_{(x,a)}=(-1)^a$ we have $\Phi(X)_{uv}=(\sum_i\varepsilon_ir_i)^2-\sum_xm_x^2$, with $m_x=\sum_a(-1)^ar_{(x,a)}$. With $\kappa=-\sigma\tau$, we have
$$m_x=(1+\kappa)\left(C_{xz}C_{xz'}+\delta_{xz}\delta_{xz'}\sigma\tau\right)+(1-\kappa)\left(\delta_{xz}\sigma C_{zz'}+\delta_{xz'}\tau C_{zz'}\right)$$
and by using $C^2=qI$, the 4 cases $z=z'$ or not, $c=c'$ or not, give $\Phi(X)_{uv}=8,-4(q-1),4q(q-3),0$ on the supports of $P_1,P_2,P_3,P_4$, which is the formula of $\Phi(X)$. Finally, $\Phi(Y)=\Phi(J)-\Phi(P_3)-\Phi(P_4)$ gives the formula of $\Phi(Y)$. Now in $\Phi^2(X)$, the coefficients of $Y$, $P_3$, $P_4$ cancel, and the coefficient of $X$ is $4(q+1)^2+4(q-3)^2+8q(q-3)^2=\beta$. As for the last assertion, $2q+1$ is prime to $q+1$, and $(q+1)^2|(q-3)^2$ would imply $q+1|4$, so $\beta/N^2$ is a rational number which is not an integer.
\end{proof}

These formulae were checked numerically by \texttt{piiphi.py} and \texttt{piibeta.py}. The proof uses only the conference matrix axioms, so this applies as well to the symmetric conference matrices which are not of Paley type, such as the one of order 46 found by Mathon \cite{mat}, and gives new examples of infinite depth subfactors, of order 92. In matrix terms, $\Phi$ has on our 4-dimensional space the eigenvalues $\pm N^{3/2}$ and $\pm\sqrt{\beta}$, the eigenvalue $N^3$ of $\Phi^2$ corresponding to the Temperley--Lieb vectors $1_N$ and $J$. Montgomery's eigenvector \cite{mo1}, with eigenvalue $64q$ of $N^3T_2$, giving $16q/(q+1)^2$, is $S\otimes D$, with $S$ the matrix $(\chi(y-x))_{x,y\in\mathbb F_q}$ extended by 0 at $\infty$, and $D=diag(1,-1)$, and its computation uses the Paley structure, via $S^2=qI-J$. There are 2 more eigenvectors of the same type, as follows:

\begin{proposition}\label{prop:alpha}
For the Paley II matrix, the matrices $S\otimes1_2$ and $S\otimes\gamma$, with $\gamma=\begin{pmatrix}0&1\\1&0\end{pmatrix}$, satisfy $N^3T_2X=\alpha X$, with $\alpha=8T^2+32$, where
$$T=\sum_{s,t\in\mathbb F_q}\chi(s(s-1)t(t-1)(t-s))$$
When $q=p$ is prime, $p=a^2+b^2$ with $a$ odd, we have $T=4a^2-2p$, and so $\alpha=32((a^2-b^2)^2+1)$.
\end{proposition}

\begin{proof}
With $r_b(y)=H_{(y,b)u}H_{(y,b)v}$ for $y\in\mathbb F_q$, we have $\Phi(S\otimes E)_{uv}=\sum_{bb'}E_{bb'}\langle r_b,Sr_{b'}\rangle$. For $z\neq z'$ in $\mathbb F_q$ we have $r_b=(\sigma\tau)^b(g+(-1)^bC_{zz'}(\sigma e_z+\tau e_{z'}))$, with $g(y)=\chi((y-z)(y-z'))$. The affine change of variables $y=z+(z'-z)s$ gives $\langle g,Sg\rangle=\chi(z'-z)T$, and we have $(Sg)_z=(Sg)_{z'}=-\chi(z'-z)$, and $S_{zz'}=\chi(z'-z)$. This gives
$$\langle r_b,Sr_{b'}\rangle=(\sigma\tau)^{b+b'}\left[C_{zz'}(T+2(-1)^{b+b'}\sigma\tau)-(\sigma+\tau)((-1)^b+(-1)^{b'})\right]$$
and so $\Phi(S\otimes1_2)_{uv}=2C_{zz'}(T+2\sigma\tau)$, and $\Phi(S\otimes\gamma)_{uv}=2C_{zz'}(\sigma\tau T-2)$. When $z=z'$, or when $\infty\in\{z,z'\}$, the vectors $r_b$ are combinations of $1$, $e_z$ and $Se_z$, and by using the fact that $S$ has zero diagonal and zero row sums, and $S^2=qI-J$, a direct computation shows that the corresponding entries of $\Phi(S\otimes1_2)$ and $\Phi(S\otimes\gamma)$ vanish. Thus
$$\Phi(S\otimes1_2)=S\otimes((2T+4)1_2+(2T-4)\gamma)\quad,\quad\Phi(S\otimes\gamma)=S\otimes((2T-4)1_2-(2T+4)\gamma)$$
and since the matrix $\begin{pmatrix}2T+4&2T-4\\2T-4&-2T-4\end{pmatrix}$ squares to $(8T^2+32)1_2$, we obtain the first assertion. As for the value of $T$, the substitution $t\to1/t$ shows that $T$ is $q^2$ times the value at 1 of Greene's hypergeometric function ${}_3F_2(\chi,\chi,\chi;\varepsilon,\varepsilon\,|\,\cdot)$ over $\mathbb F_q$ \cite{gre}, whose value at $1$ for $q=p\equiv1$ mod 4 is $(4a^2-2p)/p^2$, see \cite{ono}. We have checked the formula $T=4a^2-2p$ numerically for all primes $p\equiv1$ mod 4 below 120, as well as $T=18,-14$ at $q=9,25$, with the programs \texttt{piialpha.py} and \texttt{piialphacheck.py}.
\end{proof}

Summarizing, the subspace of $\mathbb R^N\otimes\mathbb R^N$ invariant under the affine maps $x\to sx+t$ with $s$ a nonzero square, which has dimension 24, carries the eigenvalues $N^3$, $64q$, $\alpha$, $\beta$ of $N^3T_2$, and we have checked numerically, with the program \texttt{piiinv.py}, that for all primes $q\equiv1$ mod 4 below 200 the spectrum there is exactly $N^3,64q,\alpha,\beta,0$, with multiplicities $2,4,2,2,14$. Also, $\beta$ lives on the 8-dimensional subspace invariant under the larger group $PSL_2(\mathbb F_q)$, as found with \texttt{piiglob.py}. For Paley I, as we will see now, the affine-invariant subspace has dimension 6, and carries only the eigenvalues $N^3$ and $0$. This is why Paley I is harder.

\subsection{Paley I matrices}\label{sec:paley1}


\medskip

Let $q\equiv3$ mod 4 be a prime power, and $N=q+1$. We index $P^1_{q+1}$ by $\{\infty\}\cup\mathbb F_q$, so that $H_{\infty z}=1$ for any $z$, $H_{a\infty}=-1$ for $a\in\mathbb F_q$, and $H_{az}=h(z-a)$ for $a,z\in\mathbb F_q$, with $h(0)=1$ and $h(u)=\chi(u)$ for $u\neq0$. Since $H$ is real, $(T_2)_{i_1i_2,j_1j_2}=N^{-3}S_{i_1i_2j_1j_2}^2$, with $S_{i_1i_2j_1j_2}=\sum_xH_{i_1x}H_{i_2x}H_{j_1x}H_{j_2x}$, which is symmetric in the 4 indices. The affine maps $z\to sz+t$, with $s$ a nonzero square, fixing $\infty$, are automorphisms of $H$, and we use them as follows. Since $-1$ is not a square, $\mathbb F_q^*$ is the product of the squares and of $\{\pm1\}$, and the characters $\psi$ of $\mathbb F_q^*$ with $\psi(-1)=1$ are those of odd order, dividing $n=(q-1)/2$. For such a $\psi\neq1$, set $\psi(0)=0$, and consider the matrices $X_\psi$ indexed by $\{\infty\}\cup\mathbb F_q$, given by $X_\psi(x,y)=\psi(y-x)$ for $x,y\in\mathbb F_q$, and vanishing on the row and column $\infty$.

\begin{proposition}\label{prop:principal}
We have $N^3T_2X_\psi=m_\psi X_\psi$, with $m_\psi=\mu(\psi)\mu'(\psi)$, where
$$\mu(\psi)=\sum_{t\in\mathbb F_q}\psi(t)C(t)\quad,\quad\mu'(\psi)=\sum_{t\in\mathbb F_q}\psi(t)C'(t)$$
with $C(t)=\sum_ag(a)g(a+t)$, $C'(t)=\sum_ag'(a)g'(a+t)$, and $g(a)=h(-a)h(1-a)$, $g'(a)=h(a)h(a-1)$.
\end{proposition}

\begin{proof}
With the notations of section \ref{sec:paley2}, we have $N^3T_2=\Phi_{H^t}\Phi_H$. Let $M=\Phi_H(X_\psi)$. Since the maps $\sigma(z)=sz+t$ are automorphisms of $H$, and $X_\psi(\sigma x,\sigma y)=\psi(s)X_\psi(x,y)$, we have $M_{\sigma x,\sigma y}=\psi(s)M_{xy}$. Also, $M$ is symmetric, and its diagonal entries are $\sum_{ab}X_{ab}=0$. The entries $M_{\infty y}$ are constant in $y$, by translation invariance, and are multiplied by $\psi(s)$ under scaling, so they vanish, because $\psi$ is nontrivial on the squares. As for the entries $M_{xy}$ with $x\neq y$ in $\mathbb F_q$, they are of the form $f(y-x)$, with $f(sd)=\psi(s)f(d)$ for $s$ square, and $f(-d)=f(d)$, by symmetry, so $f=f(1)\psi$, and $M=\mu X_\psi$, with $\mu=M_{01}=\sum_{ab}g(a)g(b)\psi(b-a)=\sum_t\psi(t)C(t)$. The same argument applies to $\Phi_{H^t}$, which has the same automorphisms, and with $(H^t)_{a0}(H^t)_{a1}=g'(a)$ this gives $\Phi_{H^t}(X_\psi)=\mu'X_\psi$.
\end{proof}

\noindent Since $\Phi_{H^t}$ is the adjoint of $\Phi_H$, with respect to the scalar product of $M_N(\mathbb C)$, we have $\mu'(\psi)=\overline{\mu(\psi)}$, and so $m_\psi=|\mu(\psi)|^2\geq0$. These eigenvectors were found by the program \texttt{principal.py}, and the eigenvalues $m_\psi$ are computed, from a single entry, by \texttt{eigq.py}.

\medskip

\noindent\textbf{The action of $PSL_2(\mathbb F_q)$.} The automorphism group of $P^1_{q+1}$ contains $G=PSL_2(\mathbb F_q)$, acting on $\{\infty\}\cup\mathbb F_q=P^1(\mathbb F_q)$ by signed permutations, and since the entries of $N^3T_2$ are squares, $N^3T_2$ commutes with the corresponding permutation action of $G$ on pairs of points. We use the standard notations for the irreducible representations of $G$: the Steinberg representation $St$, of dimension $q$, the principal series $\pi(\psi)$, of dimension $q+1$, with $\psi\neq\psi^{-1}$ a character of the split torus, taken up to inversion, the discrete series $\sigma(\theta)$, of dimension $q-1$, with $\theta\neq\theta^{-1}$ a character of the nonsplit torus $C_{(q+1)/2}$, taken up to inversion, and the 2 representations $\rho_\pm$ of dimension $(q-1)/2$. Let $w$ be the element of order 2 of $C_{(q+1)/2}$, and recall that $G$ has order $q(q^2-1)/2$, its nontrivial elements being split, unipotent or elliptic, according to their number of fixed points on $P^1(\mathbb F_q)$, which is 2, 1 or 0, the split and nonsplit tori being cyclic, of orders $(q-1)/2$ and $(q+1)/2$. These facts, and the character table of $G$ used below, are standard, and can be found in the literature on the representations of the finite groups $SL_2(\mathbb F_q)$.

\begin{proposition}\label{prop:psl}
For the Paley I matrix, with $q\equiv3$ mod 4, the following hold:
\begin{enumerate}
\item $N^3T_2$ vanishes on the antisymmetric vectors, and on the vectors $\varphi(x)$, $\varphi(y)$, $\delta_{xy}\varphi(x)$ with $\sum\varphi=0$, and has eigenvalue $N^3$ on $1_N$ and on the all-one matrix.

\item The space of functions on unordered pairs of distinct points is, as a representation of $G$, the sum of $1$, $St$, all the $\pi(\psi)$, once each, the $\sigma(\theta)$ with $\theta(w)=-1$, twice each, and, when $q\equiv3$ mod 8, the representations $\rho_+,\rho_-$.

\item Thus $N^3T_2$ has eigenvalues $m_\psi$ on the $\pi(\psi)$, given by Proposition \ref{prop:principal}, 2 eigenvalues on each $\sigma(\theta)$ with $\theta(w)=-1$, and, when $q\equiv3$ mod 8, an eigenvalue $m_h$ on $\rho_\pm$, given by
$$m_h=\frac{1}{2|G|}\left((q-1)N^4-\sum_{u}\tau(u)-2\sum_t\theta_0(t)\tau(t)\right)$$
where $\tau(g)=Tr(N^3T_2\,g)$, the sums are over the nontrivial unipotent elements, and over the elliptic elements, and $\theta_0$ is the character of order 2 of the nonsplit torus.
\end{enumerate}
\end{proposition}

\begin{proof}
(1) The vanishing on antisymmetric vectors comes from the symmetry of $S$, and the other assertions follow from $S_{i_1i_2jj}=N\delta_{i_1i_2}$ for any $j$, which gives $\sum_jS_{i_1i_2jj}^2=N^3\delta_{i_1i_2}$, and from $\sum_{j_2}S_{i_1i_2j_1j_2}^2=N^2$, these formulae following from $HH^t=N1_N$.

(2) The space in question is $Ind_K^G(1)$, with $K$ the stabilizer of $\{\infty,0\}$, which is dihedral of order $q-1$, consisting of the split torus, of order $(q-1)/2$, and of $(q-1)/2$ involutions $x\to c/x$, which are elliptic, because $-1$ is not a square. Thus the multiplicity of $\rho$ is $\frac{1}{q-1}\sum_{g\in K}\chi_\rho(g)$, and the character table of $G$ gives $1$ for $1,St,\pi(\psi)$, then $1-\theta(w)$ for $\sigma(\theta)$, and $(1-\theta_0(w))/2$ for $\rho_\pm$, with $\theta_0(w)=(-1)^{(q+1)/4}$.

(3) The first assertions follow from (1) and (2). As for $m_h$, the projection onto the isotypic component of $\rho_+\oplus\rho_-$ is $e=\frac{q-1}{2|G|}\sum_g\chi(g)g$, with $\chi=\chi_{\rho_+}+\chi_{\rho_-}$, which is the character of $\sigma(\theta_0)$, equal to $q-1$, $0$, $-1$, $-2\theta_0(t)$ on the identity, the split, the unipotent and the elliptic elements. Since $Tr(N^3T_2e)=(q-1)m_h$, we obtain the formula.
\end{proof}

Accordingly, the spectrum of $N^3T_2$ is the multiset union of $N^3$ with multiplicity 2, of $0$ with multiplicity $(q+1)(q+4)/2-2$, of the numbers $m_\psi$, with multiplicity $q+1$, and of $(q+1)/4$ further numbers, with multiplicity $q-1$, the multiplicities adding up when some of these numbers coincide. For instance at $q=7$, where $P_8^1\sim W_8$, the eigenvalue $N^3$ has multiplicity 8, and the eigenvalue 0 has multiplicity 56. We checked this description at $q=19,23,31$, with the programs \texttt{fullspec.py} and \texttt{orbits.py}. The numbers $\tau(g)$ are easily computed, as $\tau(g)=\|U^tU\|^2$ with $U_{ix}=H_{ix}H_{gi,x}$, and using only one element in each unipotent subgroup and one nonsplit torus, the formula for $m_h$ takes a few seconds for $q\leq100$, with the program \texttt{halftrace.py}. We find that $m_h$ is always an integer, given by
$$m_h=(U-2)^2+16q\quad,\quad U=4a^2-2q$$
where $q=a^2+2b^2$ with $a$ odd, for all the 19 primes $q\equiv3$ mod 8 below 380, for instance $m_h=320,1600,832,8000$ at $q=11,19,43,59$. Here $U$ is the trace of the Frobenius over $\mathbb F_{q^2}$ of the elliptic curve $y^2=x^3+4x^2+2x$, which has complex multiplication by $\mathbb Z[\sqrt{-2}]$. Similarly, for Paley II, the number $T=4a^2-2q$ in Proposition \ref{prop:alpha} is the trace of the Frobenius over $\mathbb F_{q^2}$ of the curve $y^2=x^3-x$, which has complex multiplication by $\mathbb Z[i]$. Whether this formula of $m_h$ holds for all $q\equiv3$ mod 8 is open, and Proposition \ref{prop:psl} (3) might give a way of proving it. As for the affine-invariant subspace, it is spanned by the indicators of the 6 orbits of pairs, and carries only the eigenvalues $N^3$, twice, and $0$.

Thus $NT_2$ has the eigenvalue $m_\psi/N^2$. Up to boundary terms, $C(t)$ is the character sum $\sum_a\chi(f_t(a))$ with $f_t(a)=a(a-1)(a+t)(a+t-1)$, related to the number of points of the curve $y^2=f_t(x)$ over $\mathbb F_q$, and so $m_\psi$ has no simple closed form, contrary to the Paley II case. However, $m_\psi\in\mathbb Z[e^{2\pi i/d}]$, with $d$ the order of $\psi$, and the Galois group acts on the characters of order $d$, so the trace $P_d=\sum_{o(\psi)=d}m_\psi$ is an integer. If $m_\psi/N^2$ were an algebraic integer for some $\psi$ of order $d$, then so would be its conjugates, and so $N^2|P_d$. With $\mu(\psi)=2\sum_{s\ {\rm square}}\psi(s)C(s)$, the integers $P_d$ are computed exactly from the integers $C(t),C'(t)$, by using Ramanujan sums, and this is done by the program \texttt{trace\_cert.py}. For instance, with $m=m_\psi$ in the cases where all the characters of order $d$ give the same eigenvalue, we have:

\medskip

\begin{center}
\begin{tabular}{lcccccc}
\toprule
$q$ & 11 & 19 & 23 & 27 & 31 & 43\\
\midrule
$d$ & 5 & 3 & 11 & 13 & 3 & 3\\
$P_d$ & 1280 & 1152 & 11264 & 9984 & 4608 & 1152\\
$m_\psi$, and $m_\psi/N^2$ & $320$, $\frac{20}{9}$ & $576$, $\frac{36}{25}$ & -- & $832$, $\frac{52}{49}$ & $2304$, $\frac{9}{4}$ & $576$, $\frac{36}{121}$\\
\bottomrule
\end{tabular}
\end{center}

\medskip

\noindent At $q=23$ the 10 characters of order 11 give 5 different eigenvalues, none of them rational, but $P_{11}=11264$ is not divisible by $N^2=576$. At $q=7$ we have $P_3=0$, in agreement with the fact that $P_8^1\sim W_8$ has a finite quantum group. We ran this over all the prime powers $q\equiv3$ mod 4 with $11\leq q<10^5$, namely 4,806 primes and the 11 prime powers $27$, $243$, $343$, $1331$, $2187$, $6859$, $12167$, $16807$, $19683$, $29791$, $79507$. A first sweep, with the programs \texttt{trace\_fast.py} for the primes, and \texttt{trace\_gf.py} for the prime powers, computes the autocorrelations and the convolutions by fast Fourier transforms, rounded to integers, which is reliable, but not a proof. The final computation, by the program \texttt{trace\_exact.py}, uses integer arithmetic for all the quantities computed. The autocorrelations $C,C'$ on the additive group $\mathbb F_q\simeq\mathbb Z_p^r$, and the cyclic convolution on $\mathbb Z_n$ giving $P_d$ via Ramanujan sums, are computed exactly, by Kronecker substitution, that is, as products of big integers whose base $2^{64}$ digits are the coefficients, these products being done by the GMP library, through the module \texttt{gmpconv.py}. The field $\mathbb F_{p^r}$ is built from a primitive polynomial, found and checked by fast exponentiation, and the program also checks the Paley identity $\sum_a\chi(a)\chi(a+t)=-1$ for $t\neq0$, which makes $P^1_{q+1}$ Hadamard. Finally, at $q=7,11,19,23,27,31,43,47$, it recomputes the integers $P_d$ from the full matrix $N^3T_2$, as exact traces against Ramanujan sums, after checking exactly the following fact, with $s$ a generator of the squares, and $Y_j$ the indicator of the set $\{(x,y)\,|\,y-x\in\{s^j,-s^j\}\}$, extended by 0 on the row and column $\infty$: $N^3T_2Y_j$ is the sum of a cyclic translate, by $j$, of a fixed combination of the $Y_l$, and of a term which does not depend on $j$. Thus the span of the differences $Y_j-Y_0$, which contains the matrices $X_\psi$ with $\psi\neq1$, is invariant, and the $X_\psi$ are eigenvectors. The whole sweep takes about 40 minutes, and we always find a divisor $d$ of $(q-1)/2$ with $N^2\nmid P_d$, the certifying divisors, and the sums of all the traces $P_d$, being the same as in the first sweep, for all the primes. At $q=7$, excluded, there is no such divisor. For instance, for some large values of $q$, the smallest certifying divisor $d$, and the corresponding trace $P_d$, are as follows:

\medskip

\begin{center}
\small
\begin{tabular}{lccccc}
\toprule
$q$ & 10007 & 29791 & 50023 & 79507 & 99991\\
\midrule
$d$ & 5003 & 3 & 3 & 3 & 3\\
$P_d$ & 996,750,011,392 & 2,489,356,800 & 19,859,447,808 & 24,812,790,912 & 23,828,318,208\\
$P_d$ mod $N^2$ & 57,214,528 & 714,230,272 & 2,342,643,776 & 5,848,224,720 & 3,831,518,080\\
\bottomrule
\end{tabular}
\end{center}

\medskip

\noindent We obtain the following result:

\begin{theorem}\label{thm:paley}
For any prime power $q\equiv3$ mod 4 with $11\leq q<10^5$, the Paley I matrix $P^1_{q+1}$ has an infinite quantum group, and so produces an infinite depth subfactor. More precisely, there is a divisor $d$ of $(q-1)/2$ such that, for any character $\psi$ of order $d$, the eigenvalue $m_\psi/N^2$ of $NT_2$ is not an algebraic integer.
\end{theorem}

\begin{proof}
This follows from Theorem \ref{thm:integral}, Proposition \ref{prop:principal}, and the above computation, whose output is in the files \texttt{trace\_exact-primes.log}, \texttt{trace\_exact-powers.log} and \texttt{trace\_exact-brute.log}, with the corresponding json files listing, for each $q$, the certifying divisors $d$ and the first value $P_d$. Further checks are done by the programs \texttt{verify\_gf.py} and \texttt{exactcheck.py}, and the first sweep is in the files \texttt{trace\_small.log} and \texttt{trace\_full.log}.
\end{proof}

\section{Wild matrices, orders 6 and 7}\label{sec:wild}

We discuss here the complex Hadamard matrices of orders 6 and 7 which are not of Di\c t\u a type. At order 6 these are the known matrices $T_6$, $H_6^q$, $BF_6$, $BN_6^q$, and the Karlsson family $K_6$, which contains all of them except $T_6$, and at order 7 these are the Petrescu matrices. All of them have $c_2=2$, at generic parameters for the families, the family $K_6$ containing special points with $c_2=4,6$, and are Temperley--Lieb as far as computed, the higher levels being discussed in section \ref{sec:level5}.

\subsection{Order 6}

The Tao matrix $T_6$ \cite{tao}, see \cite[Proposition 6.16]{banicaHM}, is isolated, and regular:
$$T_6=\begin{pmatrix}
1&1&1&1&1&1\\
1&1&w&w&w^2&w^2\\
1&w&1&w^2&w^2&w\\
1&w&w^2&1&w&w^2\\
1&w^2&w^2&w&1&w\\
1&w^2&w&w^2&w&1
\end{pmatrix}$$

\noindent The Haagerup matrices $H_6^q$ \cite{ha1}, see \cite[Proposition 6.17]{banicaHM}, form a regular 1-parameter family, with $q\in\T$:
$$H_6^q=\begin{pmatrix}
1&1&1&1&1&1\\
1&-1&i&i&-i&-i\\
1&i&-1&-i&q&-q\\
1&i&-i&-1&-q&q\\
1&-i&\bar{q}&-\bar{q}&i&-1\\
1&-i&-\bar{q}&\bar{q}&-1&i
\end{pmatrix}$$

\noindent The Bj\"orck--Fr\"oberg matrix \cite{bjf}, see \cite[Proposition 5.17]{banicaHM}, is circulant, and non-regular, with $a$ a root of $a^2+(\sqrt3-1)a+1=0$, the 2 roots $a,\bar{a}$ giving 2 circulant matrices:
$$BF_6=\begin{pmatrix}
1&ia&-a&-i&-\bar{a}&i\bar{a}\\
i\bar{a}&1&ia&-a&-i&-\bar{a}\\
-\bar{a}&i\bar{a}&1&ia&-a&-i\\
-i&-\bar{a}&i\bar{a}&1&ia&-a\\
-a&-i&-\bar{a}&i\bar{a}&1&ia\\
ia&-a&-i&-\bar{a}&i\bar{a}&1
\end{pmatrix}$$

\noindent The Beauchamp--Nicoara matrices \cite{ben}, see \cite[Theorem 5.18]{banicaHM}, are self-adjoint and non-regular, and every self-adjoint $6\times6$ complex Hadamard matrix is of this form:
$$BN_6^q=\begin{pmatrix}
1&1&1&1&1&1\\
1&-1&\bar{x}&-y&-\bar{x}&y\\
1&x&-1&t&-t&-x\\
1&-\bar{y}&\bar{t}&-1&\bar{y}&-\bar{t}\\
1&-x&-\bar{t}&y&1&\bar{z}\\
1&\bar{y}&-\bar{x}&-t&z&1
\end{pmatrix}$$
$$x=\frac{1+2q+q^2-\sqrt{2}\sqrt{1+2q+2q^3+q^4}}{1+2q-q^2},\qquad t=\frac{1+2q+q^2-\sqrt{2}\sqrt{1+2q+2q^3+q^4}}{-1+2q+q^2}$$
$$y=q,\qquad z=\frac{1+2q-q^2}{q(-1+2q+q^2)}$$
with the principal branch of the square root. Note that this matrix is Hadamard only on the arc $Re(q)\leq(\sqrt3-1)/2$, that is, $|\arg q|\geq68.53^\circ$, and not on the whole circle, as stated in \cite{banicaHM}. At the endpoints we have $q^4+2q^3+2q+1=0$, and the constant $\sqrt3-1$ is that of the Bj\"orck--Fr\"oberg equation. The matrix $D_6$ appears at $q=i$.

\medskip

\noindent\textbf{$H_2$-reducible matrices.} Following Karlsson \cite{ka2}, a complex Hadamard matrix of order $2q$ is called $H_2$-reducible if, up to permutations of the rows and columns, its $q^2$ blocks of size $2\times2$ are all Hadamard. The basic construction of such matrices is by gluing 2 Hadamard matrices of order $q$, as follows. We use the Bloch sphere: to a unit vector $\xi\in\mathbb C^2$ we associate $r_\xi\in S^2$, with $|\langle\xi,\eta\rangle|^2=\frac{1}{2}(1+\langle r_\xi,r_\eta\rangle)$, and we use spherical coordinates on $S^2$, with polar angle measured from $r_{e_1}$.

\begin{proposition}\label{prop:glue}
Let $K,L\in M_q(\T)$ be Hadamard, and let $A_0,\ldots,A_{q-1}\in U_2$ and $B_0,\ldots,B_{q-1}\in\sqrt2\,U_2$. The matrix $H\in M_{2q}(\mathbb C)$ having as $2\times2$ blocks
$$H_{ij}=A_i\begin{pmatrix}K_{ij}&0\\0&L_{ij}\end{pmatrix}B_j$$
satisfies $HH^*=2q$, and it is Hadamard if and only if, for any $i,j$, we have
$$\cos\alpha_i\cos\beta_j+\sin\alpha_i\sin\beta_j\cos(\varphi_i-\psi_j-\gamma_{ij})=0$$
where $(\alpha_i,\varphi_i)$ and $(\beta_j,\psi_j)$ are the coordinates of the Bloch vectors of $A_i^*e_1$ and $\frac{1}{\sqrt2}B_je_1$, and $e^{i\gamma_{ij}}=L_{ij}/K_{ij}$. When $H$ is Hadamard, it is $H_2$-reducible.
\end{proposition}

\begin{proof}
With $A=diag(A_i)$ and $B=diag(B_j)$ we have $H=A(K\otimes e_{11}+L\otimes e_{22})B$, up to a permutation, and the middle matrix is $\sqrt q$ times a unitary, so $HH^*=2q$. Thus each block $H_{ij}$ is $\sqrt2$ times a unitary, and a $2\times2$ unitary $V$ has $|V_{11}|=|V_{22}|$ and $|V_{12}|=|V_{21}|$, with $|V_{11}|^2+|V_{12}|^2=1$. So $H_{ij}$ is Hadamard iff $|(H_{ij})_{11}|=1$. Now $(H_{ij})_{11}=\langle D_{ij}B_je_1,A_i^*e_1\rangle$ with $D_{ij}=diag(K_{ij},L_{ij})$, and since $D_{ij}$ acts on the Bloch sphere by the rotation of angle $\gamma_{ij}$ around $r_{e_1}$, we obtain $|(H_{ij})_{11}|^2=1+\cos\alpha_i\cos\beta_j+\sin\alpha_i\sin\beta_j\cos(\varphi_i-\psi_j-\gamma_{ij})$.
\end{proof}

\noindent Conversely, if $H$ is $H_2$-reducible, with the blocks $H_{ij}$ all Hadamard, we can write $H_{ij}=X_iW_{ij}Y_j$ with $Y_j=H_{0j}$, $X_i=H_{i0}H_{00}^{-1}\in U_2$, and $W_{ij}\in U_2$ with $W_{0j}=W_{i0}=1$. The block matrix $W$ is then $\sqrt q$ times a unitary, that is, a Hadamard matrix with entries in $U_2$, and $H$ is of the above form precisely when the $W_{ij}$ commute: in this case $W_{ij}=U\,diag(K_{ij},L_{ij})U^*$, and we can take $A_i=X_iU$, $B_j=U^*Y_j$. At order 6 this is automatic, and we recover in this way the result of Karlsson \cite{ka2,kar}, as follows:

\begin{theorem}\label{thm:karlsson}
Any $H_2$-reducible matrix of order 6 is equivalent to a matrix as in Proposition \ref{prop:glue}, with $K=F_3$ and $L=F_3$ or $L=\bar{F}_3$. In the first case $H$ is of Di\c t\u a type, and in the second case $H_{ij}=X_i\Omega^{ij}Y_j$, with $\Omega\in U_2$ satisfying $1+\Omega+\Omega^2=0$, and having eigenvalues $w,\bar{w}$.
\end{theorem}

\begin{proof}
With the above notations, the orthogonality of the block rows $0,1$, and of the block columns $0,1$, gives $1+W_{11}+W_{12}=0$ and $1+W_{11}+W_{21}=0$, and then the orthogonality of the block rows $0,2$ gives $W_{22}=-1-W_{21}=W_{11}$. Let $\Omega=W_{11}$. Since $-1-\Omega$ is unitary, the eigenvalues $\lambda$ of $\Omega$ satisfy $|1+\lambda|=1$, so $\lambda\in\{w,\bar{w}\}$, hence $\lambda^2=-1-\lambda$, and so $W_{ij}=\Omega^{ij}$. Diagonalizing $\Omega$, we obtain the form in Proposition \ref{prop:glue}, with $K,L\in\{F_3,\bar{F}_3\}$, and up to exchanging $K,L$, and permuting the block rows by $i\to-i$, we can assume $K=F_3$. If $L=F_3$ then $\gamma_{ij}=0$, and the conditions read $r_i\perp s_j$ for all $i,j$, with $r_i,s_j$ the Bloch vectors in Proposition \ref{prop:glue}. Thus either all $r_i$ are $\pm r$, or all $s_j$ are $\pm s$, for some $r,s\in S^2$. In the first case we can take $A_i$ diagonal, up to multiplying on the right by a fixed unitary, which can be moved into the $B_j$, and $H_{(ir),(jc)}=a_{ir}w^{ij}(B_j)_{rc}$ is of Di\c t\u a type, with the $B_j$ being Hadamard. The second case is similar.
\end{proof}

\noindent In the second case we obtain the Karlsson matrices \cite{kar}, a non-affine family with 3 real parameters, containing the previously known 1- and 2-parameter families. In the original coordinates of \cite{kar}, which we use for computations, they are as follows. Let $\Lambda=n\cdot\sigma$ be the Hermitian unitary associated to a unit vector $n\in S^2$, with polar coordinates $\theta,\phi$, explicitly $\Lambda=\begin{pmatrix}\cos\theta&e^{i\phi}\sin\theta\\e^{-i\phi}\sin\theta&-\cos\theta\end{pmatrix}$, and consider the matrix cube root of unity
$$\Omega=e^{2\pi i\Lambda/3}=-\frac{1}{2}+\frac{i\sqrt3}{2}\Lambda\quad,\quad\Omega^3=1\quad,\quad1+\Omega+\Omega^2=0$$
Then, with $F_2=\begin{pmatrix}1&1\\1&-1\end{pmatrix}$ and $D_z=diag(1,z)$, the matrix is
$$K_6=\begin{pmatrix}X_0\\&X_1\\&&X_2\end{pmatrix}\begin{pmatrix}1&1&1\\1&\Omega&\Omega^2\\1&\Omega^2&\Omega\end{pmatrix}\begin{pmatrix}Y_0\\&Y_1\\&&Y_2\end{pmatrix}$$
with $Y_j=D_{z_j}F_2$, where $z_0=1$, and with $X_0=1$ and $X_i=\frac{1}{2}F_2D_{z_{i+2}}F_2$ for $i=1,2$. Thus $K_6$ is a Fourier matrix $F_3$ with entries in $M_2(\mathbb C)$, the cube root of unity being replaced by $\Omega$, or, in quaternionic terms, by the cube root of unity $-\frac{1}{2}+\frac{\sqrt3}{2}u$, with $u=i\Lambda$ a pure unit quaternion. The relation $K_6K_6^*=6$ is automatic, because the middle matrix is $\sqrt3$ times a unitary, the $X_i$ are unitary, and $Y_jY_j^*=2$, and each $2\times2$ block is $\sqrt2$ times a unitary, so that it is Hadamard as soon as its entries are unimodular. The only condition is thus that the entries are unimodular, and this determines the phases from $z_1=e^{i\psi}$, as
$$z_3^2=\mathcal M_A(z_1^2),\quad z_4^2=\mathcal M_B(z_1^2),\quad z_2^2=\mathcal M_A^{-1}(\mathcal M_B(z_1^2)),\qquad \mathcal M_X(z)=\frac{X_{12}^2z-X_{11}^2}{\bar{X}_{11}^2z-\bar{X}_{12}^2}$$
with $A=F_2\Omega$ and $B=F_2\Omega^2$, with a choice of signs for the square roots, as in \cite{kar}. The degenerate points where $\mathcal M_A$ or $\mathcal M_B$ is undefined are excluded. At $\theta=0$ the matrix $\Omega=diag(w,\bar{w})$ is diagonal, and $K_6$ is a Di\c t\u a matrix, so $\theta$ is a normal coordinate to this Di\c t\u a locus, with $\phi$ giving the direction of the tilt of $\Omega$. Note that the square roots make the formula discontinuous at $\theta=0$.

\noindent In the coordinates of Proposition \ref{prop:glue}, where $K=F_3$, $L=\bar{F}_3$ and $\gamma_{ij}=2\pi ij/3$, the family $K_6$ is the variety of configurations of 6 points $r_0,r_1,r_2,s_0,s_1,s_2\in S^2$ subject to the 9 equations there, modulo the rotations around the axis. This is a 3-dimensional variety, as it should, the 9 equations on the 11 remaining coordinates having 1 relation, as seen numerically. The Di\c t\u a points are those where the $r_i$, or the $s_j$, are all at the poles, by the argument in the proof of Theorem \ref{thm:karlsson}, and $\theta=0$ above corresponds to the $r_i$ at the poles. The other matrices $H_6^q$, $BF_6$, $BN_6^q$ above are $H_2$-reducible \cite{ka2}, and so belong to $K_6$, while $T_6$ is not, and their configurations are special. For the Haagerup matrix $H_6^q$, in each of its 6 block structures, 4 of the 6 points make the magic angle with the axis, $\cos^2\alpha_i=\cos^2\beta_j=\frac{1}{3}$, and the remaining 2 satisfy $\cos^2\alpha_0+\cos^2\beta_0=\frac{2}{3}$, with suitable labels. For the self-adjoint matrices $BN_6^q$ and $BF_6$, we have $\beta_j=\pi-\alpha_j$ for all $j$. These facts were found numerically, at several values of the parameters, with the programs \texttt{bloch.py} and \texttt{h2red.py}.

\noindent The 4-parameter family of Sz\"oll\H{o}si \cite{sz4}, conjecturally generic at $N=6$, has no closed form, and is not treated here. A complete classification at $N=6$, as the union of this family, of $K_6$, and of $T_6$, with an algebraic reconstruction of the generic family, was recently announced in \cite{cwt}. For $H$ and $H^t$ alike, we find, with the program \texttt{stage3.py}:
\begin{enumerate}
\item $T_6$, and $BF_6$ for both roots $a,\bar{a}$: $c_1,\ldots,c_4=1,2,5,14$.

\item $H_6^q$, at all roots of unity $q$ of order $\leq12$ and 3 random points: $1,2,5,14$.

\item $BN_6^q$, at all roots of unity of order $\leq12$ on its arc, including $q=i$, which gives $D_6$, and $q=-1$, and at 3 random points of the arc: $1,2,5,14$.

\item $K_6$, at 3 random points: $1,2,5,14$. On the grid $\theta,\phi,\psi\in\frac{\pi}{6}\{0,\ldots,5\}$, at $k\leq3$: $1,2,5$ at 164 of the 185 valid points, and $1,4,18$ at the other 21, the 31 degenerate points being skipped. These are the generic values of the order 6 Di\c t\u a families, which are contained in $K_6$.
\end{enumerate}

\noindent The value $c_2^{gen}=2$ for $K_6$ can be proved exactly: at the Haagerup point $H_6^1\in K_6$, with entries in $\{\pm1,\pm i\}$, an exact computation in the Gaussian integers, by the program \texttt{c2exact.py}, shows that the graph of section \ref{sec:qg} has 2 components, of sizes 6 and 30, so $c_2(H_6^1)=2$, and by upper semicontinuity of $c_2=\dim\ker(T_2-1)$, we have $c_2=2$ near $H_6^1$, and so at the generic points of $K_6$.

\noindent Thus at order 6, every known non-Di\c t\u a matrix is Temperley--Lieb up to level 4, which extends the $c_2=2$ observations of \cite{cn}.

\subsection{Order 7}

Every regular complex Hadamard matrix of order 7 is equivalent to $F_7$, $P_7^q$ or $R_7^q$, with $q\in\T$ \cite{gt}, see also \cite{pet}. Here $R_7^q$ is a monomially normalized conjugate of $P_7^q$, kept because conjugation is not part of the equivalence, so that $c_k(R_7^q)=c_k(P_7^{q'})$ for some $q'$:
$$P_7^q=\begin{pmatrix}
-q&q&w&1&w&1&w\\
q&-q&w&1&1&w&w\\
w&w&-w&1&w&w&1\\
1&1&1&-1&w&w&w\\
w&1&w&w&-\bar{q}w&\bar{q}w&1\\
1&w&w&w&\bar{q}w&-\bar{q}w&1\\
w&w&1&w&1&1&-1
\end{pmatrix}$$
$$R_7^q=\begin{pmatrix}
1&1&1&1&1&1&1\\
1&1&1&w&-1&-1&w^2\\
1&-w^2&w^2&w^2&q&-q&w\\
1&w&-1&1&qw^2&-qw^2&w^2\\
1&-1&-w&-1&-q&q&-w^2\\
1&w^2&w&qw&q&-qw&-q\\
1&w^2&w&-qw&qw&-q&q
\end{pmatrix}$$

\noindent At 3 random points and at all roots of unity $q$ of order $\leq12$, for $H$ and $H^t$ alike, we find $c_1,\ldots,c_4=1,2,5,14$ everywhere, and in particular $c_3=5$ at $q=1$, as found in \cite{cn}, the program being \texttt{stage46.py}.

\section{Wild matrices, orders 8 to 11}\label{sec:wild8}

We continue with the non-Di\c t\u a matrices of orders 8 to 11, namely the regular matrices of order 8 from \cite{gt}, the isolated matrices of orders 9, 10, 11 found in \cite{ben}, and some new glued matrices of order 10. Then we discuss the higher levels, up to level 9, for all the wild candidates, as well as the special points of the families, and we summarize the findings.

\subsection{Order 8}

Among the 7 regular families of \cite{gt} in which only the cycle structures $A,B$ appear, the 3 families $F_{8,5}$ of section \ref{sec:f85} are of Di\c t\u a type, and the 4 which are not of Di\c t\u a type are the following, with independent parameters $a,b,c,d\in\T$, displayed as in \cite{gt}, in dephased form, the transposed families being displayed too, as they appear in \cite{gt}, and as they are used in the computations below:

\begin{center}\normalsize\setlength{\arraycolsep}{3.8pt}
\begin{adjustbox}{max width=\textwidth}$\displaystyle F_{8,4}(a,b,c,d)=\begin{pmatrix}
1 & 1 & 1 & 1 & 1 & 1 & 1 & 1 \\
1 & -1 & ad & -ad & -bd & bd & 1 & -1 \\
1 & -\frac{1}{c} & d & d & -d & -d & -1 & \frac{1}{c} \\
1 & \frac{1}{cd} & -ad & ad & -1 & -1 & 1 & -\frac{1}{cd} \\
1 & \frac{1}{c} & ad^{2} & -ad^{2} & bd^{2} & -bd^{2} & -1 & -\frac{1}{c} \\
1 & d & -d & -d & -bd^{2} & bd^{2} & -1 & d \\
1 & -\frac{1}{cd} & -1 & -1 & bd & -bd & 1 & \frac{1}{cd} \\
1 & -d & -ad^{2} & ad^{2} & d & d & -1 & -d
\end{pmatrix}$\end{adjustbox}
\end{center}

\begin{center}\normalsize\setlength{\arraycolsep}{3.8pt}
\begin{adjustbox}{max width=\textwidth}$\displaystyle F_{8,4}^{t}(a,b,c,d)=\begin{pmatrix}
1 & 1 & 1 & 1 & 1 & 1 & 1 & 1 \\
1 & -1 & -\frac{1}{c} & \frac{1}{cd} & \frac{1}{c} & d & -\frac{1}{cd} & -d \\
1 & ad & d & -ad & ad^{2} & -d & -1 & -ad^{2} \\
1 & -ad & d & ad & -ad^{2} & -d & -1 & ad^{2} \\
1 & -bd & -d & -1 & bd^{2} & -bd^{2} & bd & d \\
1 & bd & -d & -1 & -bd^{2} & bd^{2} & -bd & d \\
1 & 1 & -1 & 1 & -1 & -1 & 1 & -1 \\
1 & -1 & \frac{1}{c} & -\frac{1}{cd} & -\frac{1}{c} & d & \frac{1}{cd} & -d
\end{pmatrix}$\end{adjustbox}
\end{center}

\begin{center}\normalsize\setlength{\arraycolsep}{3.8pt}
\begin{adjustbox}{max width=\textwidth}$\displaystyle F_{8,2}(a,b)=\begin{pmatrix}
1 & 1 & 1 & 1 & 1 & 1 & 1 & 1 \\
1 & -1 & b & -b & -\frac{b}{a} & \frac{b}{a} & a & -a \\
1 & -a & \frac{b}{a} & b & -\frac{b}{a} & -b & -1 & a \\
1 & \frac{1}{a} & \frac{b}{a^{2}} & -\frac{b}{a} & \frac{b}{a} & -\frac{b}{a^{2}} & -\frac{1}{a} & -1 \\
1 & 1 & -1 & -1 & -1 & -1 & 1 & 1 \\
1 & a & -b & \frac{b}{a} & -\frac{b}{a} & b & -a & -1 \\
1 & -\frac{1}{a} & -\frac{b}{a} & -\frac{b}{a^{2}} & \frac{b}{a} & \frac{b}{a^{2}} & -1 & \frac{1}{a} \\
1 & -1 & -\frac{b}{a^{2}} & \frac{b}{a^{2}} & \frac{b}{a} & -\frac{b}{a} & \frac{1}{a} & -\frac{1}{a}
\end{pmatrix}$\end{adjustbox}
\end{center}

\begin{center}\normalsize\setlength{\arraycolsep}{3.8pt}
\begin{adjustbox}{max width=\textwidth}$\displaystyle F_{8,2}^{t}(a,b)=\begin{pmatrix}
1 & 1 & 1 & 1 & 1 & 1 & 1 & 1 \\
1 & -1 & -a & \frac{1}{a} & 1 & a & -\frac{1}{a} & -1 \\
1 & b & \frac{b}{a} & \frac{b}{a^{2}} & -1 & -b & -\frac{b}{a} & -\frac{b}{a^{2}} \\
1 & -b & b & -\frac{b}{a} & -1 & \frac{b}{a} & -\frac{b}{a^{2}} & \frac{b}{a^{2}} \\
1 & -\frac{b}{a} & -\frac{b}{a} & \frac{b}{a} & -1 & -\frac{b}{a} & \frac{b}{a} & \frac{b}{a} \\
1 & \frac{b}{a} & -b & -\frac{b}{a^{2}} & -1 & b & \frac{b}{a^{2}} & -\frac{b}{a} \\
1 & a & -1 & -\frac{1}{a} & 1 & -a & -1 & \frac{1}{a} \\
1 & -a & a & -1 & 1 & -1 & \frac{1}{a} & -\frac{1}{a}
\end{pmatrix}$\end{adjustbox}
\end{center}

\noindent The remaining regular families of \cite{gt}, involving other cycle structures, are the following 2, with $a,b\in\T$, which are not of Di\c t\u a type either, as explained below:

\begin{center}\normalsize\setlength{\arraycolsep}{3.8pt}
\begin{adjustbox}{max width=\textwidth}$\displaystyle K_{8,2}(a,b)=\begin{pmatrix}
1 & 1 & 1 & 1 & 1 & 1 & 1 & 1 \\
1 & a & -b^{2} & -\frac{b^{2}}{a} & -1 & -a & \frac{b^{2}}{a} & b^{2} \\
1 & \frac{a}{b} & -1 & \frac{b}{a} & 1 & -\frac{a}{b} & -\frac{b}{a} & -1 \\
1 & -b & b & \frac{b^{2}}{a} & -\frac{b}{a} & -1 & \frac{b}{a} & -\frac{b^{2}}{a} \\
1 & -a & -b & b & -\frac{a}{b} & \frac{a}{b} & -1 & a \\
1 & -1 & b & -\frac{b}{a} & \frac{b}{a} & -b & -\frac{b^{2}}{a} & \frac{b^{2}}{a} \\
1 & -\frac{a}{b} & -b & -1 & \frac{a}{b} & a & b & -a \\
1 & b & b^{2} & -b & -1 & b & -b & -b^{2}
\end{pmatrix}$\end{adjustbox}
\end{center}

\begin{center}\normalsize\setlength{\arraycolsep}{3.8pt}
\begin{adjustbox}{max width=\textwidth}$\displaystyle M_{8,1}(a)=\begin{pmatrix}
1 & 1 & 1 & 1 & 1 & 1 & 1 & 1 \\
1 & -wa^{2} & a & wa & -1 & wa^{2} & -a & -wa \\
1 & a & -a & a^{3} & -a^{3} & -a^{4} & a^{4} & -1 \\
1 & wa & -\frac{w}{a} & -wa & \frac{w}{a} & w^2 & -1 & -w^2 \\
1 & -1 & \frac{w}{a} & -a^{3} & wa^{2} & -wa^{2} & a^{3} & -\frac{w}{a} \\
1 & wa^{2} & w & -wa^{4} & -wa^{2} & -w^2a^{4} & -a^{4} & w^2 \\
1 & -a & -w & -1 & -\frac{w}{a} & w & a & \frac{w}{a} \\
1 & -wa & -1 & wa^{4} & a^{3} & -wa^{4} & -a^{3} & wa
\end{pmatrix}$\end{adjustbox}
\end{center}

\noindent Finally, there are 4 isolated regular matrices, which are Butson, given here in logarithmic form, with $\omega=e(1/12)$ for $I_{8,12}$ and $\omega=e(1/30)$ for the other 3, these logarithmic forms being also the ones used in the module \texttt{matrices.py}, which builds all the matrices of this paper, with the conventions used here, the 4 matrices themselves being taken from the classification in \cite{gt}:

{\setlength{\arraycolsep}{4pt}$$I'_{8,12}=\begin{pmatrix}0 & 0 & 0 & 0 & 0 & 0 & 0 & 0\\ 0 & 0 & 0 & 3 & 6 & 6 & 6 & 9\\ 0 & 1 & 5 & 9 & 1 & 5 & 9 & 6\\ 0 & 4 & 8 & 6 & 10 & 2 & 3 & 9\\ 0 & 6 & 2 & 11 & 9 & 8 & 3 & 5\\ 0 & 6 & 11 & 5 & 3 & 11 & 9 & 5\\ 0 & 7 & 6 & 1 & 7 & 3 & 9 & 1\\ 0 & 10 & 6 & 7 & 4 & 9 & 3 & 1\end{pmatrix}$$}
{\setlength{\arraycolsep}{4pt}$$K'_{8,30}=\begin{pmatrix}0 & 0 & 0 & 0 & 0 & 0 & 0 & 0\\ 0 & 0 & 6 & 10 & 12 & 18 & 20 & 24\\ 0 & 6 & 13 & 28 & 27 & 21 & 12 & 15\\ 0 & 10 & 28 & 20 & 16 & 4 & 10 & 22\\ 0 & 12 & 27 & 16 & 1 & 15 & 24 & 9\\ 0 & 18 & 21 & 4 & 15 & 19 & 6 & 3\\ 0 & 20 & 12 & 10 & 24 & 6 & 0 & 18\\ 0 & 24 & 15 & 22 & 9 & 3 & 18 & 7\end{pmatrix}$$}
{\setlength{\arraycolsep}{4pt}$$L'_{8,30}=\begin{pmatrix}0 & 0 & 0 & 0 & 0 & 0 & 0 & 0\\ 0 & 0 & 6 & 10 & 12 & 18 & 20 & 24\\ 0 & 6 & 21 & 2 & 19 & 17 & 4 & 15\\ 0 & 10 & 17 & 25 & 8 & 2 & 15 & 23\\ 0 & 12 & 4 & 23 & 27 & 15 & 19 & 8\\ 0 & 18 & 2 & 17 & 15 & 3 & 1 & 16\\ 0 & 20 & 19 & 15 & 4 & 16 & 5 & 1\\ 0 & 24 & 15 & 8 & 23 & 1 & 16 & 9\end{pmatrix}$$}
{\setlength{\arraycolsep}{4pt}$$L^{t\prime}_{8,30}=\begin{pmatrix}0 & 0 & 0 & 0 & 0 & 0 & 0 & 0\\ 0 & 0 & 6 & 10 & 12 & 18 & 20 & 24\\ 0 & 6 & 21 & 17 & 4 & 2 & 19 & 15\\ 0 & 10 & 2 & 25 & 23 & 17 & 15 & 8\\ 0 & 12 & 19 & 8 & 27 & 15 & 4 & 23\\ 0 & 18 & 17 & 2 & 15 & 3 & 16 & 1\\ 0 & 20 & 4 & 15 & 19 & 1 & 5 & 16\\ 0 & 24 & 15 & 23 & 8 & 16 & 1 & 9\end{pmatrix}$$}

\noindent For $H$ and $H^t$ alike, we find $c_1,\ldots,c_4=1,2,5,14$ for all these matrices, at generic parameters, with the program \texttt{stage46.py}. The value $c_2=2$ of $F_{8,2},F_{8,4}$ was checked independently with the profile graph, and the Di\c t\u a block structure search confirms that $F_{8,2},F_{8,4},K_{8,2},M_{8,1}$ are not of Di\c t\u a type. The special points found, at $k\leq3$, by the same procedure as in section \ref{sec:dita}, are as follows:
\begin{enumerate}
\item $F_{8,4},F_{8,4}^t$: all parameters $\pm1$ give $c_2,c_3=8,64$, and all parameters $\pm i$ give $5,28$.

\item $F_{8,2},F_{8,2}^t$: $c_2,c_3=5,28$ at $a=\pm1$, and the Fourier point.

\item $K_{8,2}$: $c_2,c_3=4,19$ at $b=\pm1$, and along the whole diagonal $a=b$, except at $a=b=\pm1$, where we get $8,64$, and at $a=b=\pm i$, where we get $6,40$.

\item $M_{8,1}$: $c_2,c_3=4,19$ at $a^6=1$.
\end{enumerate}

\noindent Note that $5,28$ are the generic values of $F_2\otimes_QF_4$, and $4,19$ those of $F_{8,5a}$, so all these families meet Di\c t\u a type loci, at special parameters.

\noindent\textbf{Block structure.} All the regular matrices of order 8 above, and the Di\c t\u a type ones of section \ref{sec:f85}, are $H_2$-reducible, in several ways, as found by a direct search, with the programs \texttt{types8.py} and \texttt{look8.py}. In each block structure, the normalized blocks $W_{ij}$ defined after Proposition \ref{prop:glue} are of one of 2 types. For $F_{8,4}$, $F_{8,4}^t$ and $F_{8,5a},F_{8,5b},F_{8,5c}$, in all the block structures the $W_{ij}$ commute, so these matrices are obtained by gluing 2 complex Hadamard matrices of order 4 as in Proposition \ref{prop:glue}, namely 2 matrices of the form $F_4^s$, up to permutations. For $F_{8,2}$, $K_{8,2}$, $M_{8,1}$, $I_{8,12}$, $K_{8,30}$, $L_{8,30}$ and their transposes, there are block structures in which the $W_{ij}$ do not commute, and in all of them $\det W_{ij}=1$, so that $W$ is a quaternionic Hadamard matrix of order 4, in the sense of \cite{cd}, where these matrices are classified. This dichotomy seems to be general: among 300 dephased Hadamard matrices of order 4 with entries in $U_2$, found by local minimization from random points, by the program \texttt{u2had.py}, all are commuting, or quaternionic. We do not know whether this dichotomy holds for all the Hadamard matrices of order 4 with entries in $U_2$, which would be a quaternionic analogue of the classification of the complex Hadamard matrices of order 4, recalled in section \ref{sec:numerics}. Thus all the matrices of order 8 considered here come either from a pair of complex Hadamard matrices of order 4, or from a quaternionic Hadamard matrix of order 4, and both types contain matrices with $c_2=2$.

\subsection{Orders 9 to 11}\label{sec:nine}

The following isolated matrices were found by computer in \cite[Section 3]{ben}. Here $y=\frac{1-i\sqrt{15}}{4}$ for $H_9$, $\epsilon=e(1/10)$ for $BN_9$, given in logarithmic form, and $a=\frac{-1+i\sqrt{15}}{4}$ for $BN_{10}$:
{\small$$H_9=\begin{pmatrix}
1&1&1&1&1&1&1&1&1\\
1&-1&-1&y&y^{-1}&-y^{-3}&y^{-3}&y&y^{-1}\\
1&y&y^3&-y^2&y^3&y&-1&-y&y\\
1&y&y&-1&y^{-1}&y^{-3}&-y^{-3}&-1&y^{-1}\\
1&-y&-y^2&y^3&y^3&-1&y&y&y\\
1&y&y^3&y^3&y&y^{-1}&y^{-1}&y&y^2\\
1&y^3&y&y&y^2&y^{-1}&y^{-1}&y^3&y\\
1&-y^2&-y&y&y&-1&y&y^3&y^3\\
1&y^3&y&-y&y&y&-1&-y^2&y^3
\end{pmatrix}$$}

{\small\setlength{\arraycolsep}{2.6pt}$$BN'_9=\begin{pmatrix}0 & 0 & 0 & 0 & 0 & 0 & 0 & 0 & 0\\ 0 & 5 & 3 & 3 & 5 & 9 & 8 & 7 & 1\\ 0 & 4 & 5 & 7 & 1 & 3 & 5 & 9 & 9\\ 0 & 3 & 7 & 5 & 1 & 8 & 9 & 3 & 5\\ 0 & 9 & 1 & 5 & 5 & 3 & 7 & 2 & 7\\ 0 & 9 & 5 & 1 & 3 & 5 & 1 & 7 & 6\\ 0 & 1 & 7 & 9 & 6 & 1 & 5 & 5 & 3\\ 0 & 7 & 9 & 4 & 9 & 5 & 3 & 5 & 1\\ 0 & 5 & 2 & 9 & 7 & 7 & 3 & 1 & 5\end{pmatrix}\quad\ BN_{10}=\begin{pmatrix}
1&1&1&1&1&1&1&1&1&1\\
1&1&a^{-2}&a^{-1}&a^{-2}&a^2&1&a^2&1&a\\
1&a^2&1&a^{-2}&a^{-2}&1&a&a^2&1&a^{-1}\\
1&1&1&1&a^{-2}&a^2&a^2&a&a^{-2}&a^{-1}\\
1&a^2&a^{-2}&1&1&a&1&a^2&a^{-2}&a^{-1}\\
1&1&a^{-1}&a^{-2}&1&1&a^2&a^2&a^{-2}&a\\
1&a^2&1&a^{-2}&a^{-1}&a^2&1&1&a^{-2}&a\\
1&a^2&a^{-2}&1&a^{-2}&1&a^2&1&a^{-1}&a\\
1&a&a^{-2}&a^{-2}&1&a^2&a^2&1&1&a^{-1}\\
1&a&a^{-1}&a^{-1}&a^{-1}&a&a&a&a^{-1}&1
\end{pmatrix}$$}

\noindent The 2 matrices $N_{11},N_{11}'$ have $x=\frac{3}{4}-\frac{i\sqrt7}{4}$, and $(a,b,y)=(-x,x,-x^2)$, respectively $(a,b,y)=(1,-1,\bar{x})$:
{\small\setlength{\arraycolsep}{3pt}$$N_{11},N_{11}'=\begin{pmatrix}
1&1&1&1&1&1&1&1&1&1&1\\
1&y&y&a&b&b&-1&x&-1&x&-1\\
1&y&b&b&a&y&x&x&-1&-1&-1\\
1&a&b&y&y&b&-1&-1&x&-1&x\\
1&b&a&y&b&y&x&-1&-1&-1&x\\
1&b&y&b&y&a&-1&-1&x&x&-1\\
1&-1&x&-1&x&-1&-\frac{x}{a}&-\frac{x}{b}&-\frac{x}{y}&-\frac{x}{y}&-\frac{x}{b}\\
1&x&x&-1&-1&-1&-\frac{x}{b}&-\frac{x}{y}&-\frac{x}{b}&-\frac{x}{y}&-\frac{x}{a}\\
1&-1&-1&x&-1&x&-\frac{x}{y}&-\frac{x}{b}&-\frac{x}{b}&-\frac{x}{a}&-\frac{x}{y}\\
1&x&-1&-1&-1&x&-\frac{x}{y}&-\frac{x}{y}&-\frac{x}{a}&-\frac{x}{b}&-\frac{x}{b}\\
1&-1&-1&x&x&-1&-\frac{x}{b}&-\frac{x}{a}&-\frac{x}{y}&-\frac{x}{b}&-\frac{x}{y}
\end{pmatrix}$$}

\noindent For these 5 matrices and their transposes, we find, with the program \texttt{stage46.py}, the values $c_1,\ldots,c_4=1,2,5,14$.

\noindent\textbf{Glued matrices of order 10.} The matrix $BN_{10}$ is not $H_2$-reducible. However, Proposition \ref{prop:glue} produces many matrices of order 10 with $c_2=2$, as follows. Take $K=F_5$ and $L_{ij}=\zeta^{\sigma(i)\tau(j)}$, with $\zeta=e(1/5)$ and $\sigma,\tau$ permutations of $\mathbb Z_5$ fixing 0. For each of the 576 choices of $(\sigma,\tau)$ we solved the equations of Proposition \ref{prop:glue} by local minimization from random points, and computed $c_2$, with the programs \texttt{build.py}, \texttt{classify10.py} and \texttt{ten.py}, the last one computing also the defect. When $\sigma,\tau$ are both multiplications by units, so that $W_{ij}=\Omega^{ij}$ with $\Omega^5=1$, as for $K_6$, all the solutions found have $c_2=6$, or $c_2=10$. For 500 of the remaining 560 choices, we found solutions with $c_2=2$. We then took 16 of these choices, including $\sigma=id$, $\tau=(34)$, and one solution with $c_2=2$ for each, and ran the overnight program \texttt{bones.py} on them. In all 16 cases we found $c_1,\ldots,c_6=1,2,5,14,42,132$, the values $c_5,c_6$ coming from the Jones--Wenzl count of section \ref{sec:jw}, which gives $e_5=e_6=0$, and the same values for $H^t$, by Theorem \ref{thm:transp}. For $\sigma=id$, $\tau=(34)$, the count at level 7, in dimension $10^7$, gives $e_7=0$ as well, so that $c_7=429$. Within the construction, the solution sets have local dimension 2 for 12 of the 16 choices, and 1 for the other 4, and the defect is 8, except for 2 choices, where it is 4. The largest eigenvalue found on the Jones--Wenzl pieces at levels $5,6$ is about $0.1$, except for 1 choice, where it is about $0.96$. Similarly, at order 8, the quaternionic Fourier matrix $W_{ij}=u^{ij}$, with $u$ a pure unit quaternion, only gives matrices with $c_2\geq3$, the generic values being $c_2,c_3=3,11$, as for $F_{8,5b}$. Thus the Fourier pattern of $K_6$ gives no wild candidates at orders 8 and 10, while gluing $F_5$ with a permuted copy of itself gives them at order 10. We have not checked whether these matrices belong to the known families of order 10 \cite{szp,tz1}. The matrix with $\sigma=id$, $\tau=(34)$, which was used at level 7, is given numerically in the file \texttt{example.npy} accompanying \texttt{bones.py}, and the 16 solutions are written by this program in the files \texttt{pair-a.npy} to \texttt{pair-p.npy}, from fixed random seeds. These solutions are approximate, the equations of Proposition \ref{prop:glue} being solved with residuals below $10^{-11}$, and the values above are those of these approximate matrices: we have no exact representatives, and the statements about them should be understood as numerical observations, and not as existence or classification results.

\subsection{Higher levels}\label{sec:level5}

We computed $c_5$ for all the matrices of this section, at a random point for the families, plus a few special points, for $H$ and $H^t$ alike, by the seeded method of section \ref{sec:conv} with a deflated Lanczos verification, in dimension up to $11^5=161051$, with the program \texttt{level5.py}. Then we computed $c_6$ and $c_7$ for some of them, by the probabilistic count with noncrossing seeds, with the programs \texttt{puma.py} and \texttt{gsp.py} at level 6, and \texttt{caracal.py} at level 7, and $c_6,\ldots,c_9$, by the Jones--Wenzl count of section \ref{sec:jw}, with the programs \texttt{jaguar.py}, \texttt{aerts.py} and \texttt{lennox.py}, for $H$, which suffices, by Theorem \ref{thm:transp}. The values $c_6,\ldots,c_9$ in the tables below come from probabilistic counts, but they are confirmed by the converged norms of section \ref{sec:jw}, all below 1, which is why they are written without star, except for $BF_6$ at level 9. A blank in the tables means that the corresponding value was not computed. At orders 6 and 7, we find:

\medskip

\begin{center}
\begin{tabular}{lcccccc}
\toprule
matrix & $c_5$ & next eigenvalue & $c_6$ & $c_7$ & $c_8$ & $c_9$\\
\midrule
$T_6$ & 42 & 0.619 & 132 & 429 & 1430 & 4862\\
$BF_6$ & 42 & 0.667 &  & 429 & 1430 & $4862^*$\\
$H_6^q$, $q$ random & 42 & 0.673 & 132 & 429 & 1430 & 4862\\
$H_6^q$, $q=1$ & 42 & 0.667 &  & 429 & 1430 & 4862\\
$BN_6^q$, $q$ random & 42 & 0.692 & 132 & 429 & 1430 & 4862\\
$BN_6^q$, $q=i$ ($D_6$) & 42 & 0.667 &  & 429 & 1430 & 4862\\
$K_6$, random & 42 & 0.795 & 132 & 429 & 1430 & 4862\\
$P_7^q$, $q$ random & 42 & 0.587 & 132 & 429 & 1430 & \\
$P_7^q$, $q=1$ & 42 & 0.570 & 132 & 429 & 1430 & \\
$R_7^q$, $q$ random & 42 & 0.575 & 132 & 429 & 1430 & \\
\bottomrule
\end{tabular}
\end{center}

\medskip

\noindent At orders 8 to 11, where $c_8$ and $c_9$ were not computed, we find:

\medskip

\begin{center}
\begin{tabular}{lcccc}
\toprule
matrix & $c_5$ & next eigenvalue & $c_6$ & $c_7$\\
\midrule
$F_{8,4}$, random & 42 & 0.893 & 132 & 429 \\
$F_{8,2}$, random & 42 & 0.638 & 132 & 429 \\
$K_{8,2}$, random & 42 & 0.566 & 132 & 429 \\
$M_{8,1}$, random & 42 & 0.859 & 132 & 429 \\
$I_{8,12}$ & 42 & 0.553 & 132 & 429 \\
$K_{8,30}$ & 42 & 0.541 & 132 & 429 \\
$L_{8,30}$, $L_{8,30}^t$ & 42 & 0.946 & 132 & 429 \\
$H_9$ & 42 & 0.442 & 132 & 429 \\
$BN_9$ & 42 & 0.747 & 132 & 429 \\
$BN_{10}$ & 42 & 0.648 & 132 & 429 \\
$N_{11}$, $N_{11}'$ & 42 & 0.382 & 132 &  \\
\bottomrule
\end{tabular}
\end{center}

\medskip

\noindent Also, $P_{12}^1$ has $c_6=132$, by the same method. The second eigenvalues at these levels are given in section \ref{sec:jw}, and are at least those at level 5, by Theorem \ref{thm:jw}. Thus all these matrices are Temperley--Lieb as far as computed, and so are the real matrices of section \ref{sec:real} with $c_2=2$.

\subsection{Special parameters}

The non-Di\c t\u a families $H_6^q$, $BN_6^q$, $K_6$, $P_7^q$, $R_7^q$, $F_{8,4}$, $F_{8,2}$, $K_{8,2}$, $M_{8,1}$ were scanned on fine grids, together with all the points having as parameters roots of unity of order $\leq24$ for 1 parameter, $\leq12$ for 2 parameters, and $\leq8$ for 4 parameters, plus random points, for a total of about 13,600 points with $c_2=2$, with the programs \texttt{scan.py} and \texttt{scan4.py}. At all of them we found $c_3=5$ and $c_4=14$, the 15 points where the probabilistic count gave $c_4>14$ being, by a full recomputation with dense linear algebra, artifacts of points close to a special locus, with eigenvalues $1-10^{-5}$ or $1-10^{-6}$. The special loci found are as follows:

\medskip

\begin{center}
\begin{tabular}{lll}
\toprule
family & special locus & $c_2,c_3,c_4$\\
\midrule
$F_{8,2}$ & line $a=1$ & $5,28,167$, generic $2\times4$ Di\c t\u a\\
$K_{8,2}$ & lines $a=b$ and $b=1$ & $4,19,98$, as for $F_{8,5a}$\\
$K_{8,2}$ & isolated points & $6,40$\\
$F_{8,4}$ & 69 points & $3,11$\\
$F_{8,4}$ & other points & $5,28$\\
$M_{8,1}$ & 12 points & $4,19$\\
$K_6$ & 45 points & $4,18$, generic $2\times3$ Di\c t\u a\\
$F_{8,2},K_{8,2},F_{8,4}$ & a few points & $8,64$, Fourier type\\
$H_6^q,BN_6^q,P_7^q,R_7^q$ & none & --\\
\bottomrule
\end{tabular}
\end{center}

\medskip

\noindent In all cases the jump starts at level 2.

\subsection{Wild matrices}\label{sec:nondita}

Let us call wild the complex Hadamard matrices $H$ with $G_H=S_N^+$. For such matrices the free cumulants of $\chi$ are all equal to 1, $\chi$ being free Poisson, and the Boolean cumulants are the shifted Catalan numbers $1,1,2,5,14,\ldots$ Summarizing, all the matrices of sections \ref{sec:real}, \ref{sec:wild} and \ref{sec:wild8} with $c_2=2$ are Temperley--Lieb as far as computed: the 24 matrices of section \ref{sec:level5} up to $c_5$, without star, and up to $c_6=132$, 20 of them up to $c_7=429$, the 3 of order 7 up to $c_8=1430$, and the 7 of order 6 up to $c_9=4862$, with a star for $BF_6$ at level 9, the 13,600 special points up to $c_4$, and the 532 real matrices of orders 24 and 28 with $c_2=2$ up to $c_4$, respectively $c_3$. The Di\c t\u a block structure search was run directly on $F_{8,2},F_{8,4},K_{8,2},M_{8,1},P_7^q$, and for the other matrices the absence of a Di\c t\u a structure follows from $c_2=2$, a Di\c t\u a structure producing a block partition which is invariant under $G$, and so $c_2\geq3$. In addition, at all the special points of the non-Di\c t\u a families, the jump from the Temperley--Lieb values happens at level 2. Thus, for all the matrices tested with $c_2=2$, meaning that $G_H$ is quantum 2-transitive, we find $c_k=Cat_k$ at all the levels computed, which extends the observations of \cite{cn} at levels 2 and 3, and raises the question whether all these matrices are wild.

Note that the matrices of Di\c t\u a type have $c_2\geq3$, but that the converse fails. Consider the point $(a,b,c,d,e)=(z,z^2,z^4,z^8,z^{16})$ of $F_{8,5b}$, with $z=e^{2\pi i/101}$. All its entries are roots of unity of order 202, and an exact computation in $\mathbb Q(e^{2\pi i/202})$, by reduction modulo the cyclotomic polynomial, done by the program \texttt{f85bpoint.py}, shows that $c_2=3$, that the only Di\c t\u a block structure is the one of size $2\times4$ from Proposition \ref{prop:f85b}, and that the inner blocks $K,L$ have different Haagerup sets $\{H_{ij}\bar{H}_{kj}H_{kl}\bar{H}_{il}\}$, namely $\{\pm1,\pm z^{\pm10}\}$ and $\{\pm1,\pm z^{\pm20}\}$. Thus this matrix is of the form $F_2\otimes_Q(K,L)$, but not of the form $H\otimes_QK$. As for its transpose, which has $c_2=3$ as well, by Theorem \ref{thm:transp}, the same program shows that its only candidate block structure, of size $4\times2$, has outer blocks with the same 2 different Haagerup sets, while in a matrix $H\otimes_Q(K_1,\ldots,K_M)$ the outer blocks, obtained by restricting one row of each group to one class of columns, are all equivalent to $H$. Thus the transpose is not of Di\c t\u a type, even in the general sense of section 1, and has $c_2=3$. The same happens for the family $F_{8,5c}$ itself, at the same point: the program \texttt{f85cpoint.py} shows that $c_2=3$, and that the only candidate structure, of size $4\times2$, has outer blocks with different Haagerup sets, $\{\pm1,\pm z^{\pm4}\}$ and $\{\pm1,\pm z^{\pm1}\}$, while its transpose has the form $F_2\otimes_Q(K,L)$. Thus the relevant condition here is $c_2=2$, and not the absence of a Di\c t\u a structure.

\section{The Jones--Wenzl decomposition}\label{sec:jw}

\subsection{Jones--Wenzl pieces}\label{sec:jwpieces}

For $N\geq4$ the irreducible corepresentations of $S_N^+$ are $r_0=1,r_1,r_2,\ldots$, with $u=r_0+r_1$ and $r_1\otimes r_s=r_{s-1}+r_s+r_{s+1}$ for $s\geq1$, so that $u\otimes r_s=r_{s-1}+2r_s+r_{s+1}$. The dimensions are $d_0=1$, $d_1=N-1$ and $d_{s+1}=(N-2)d_s-d_{s-1}$, and $r_s\subset u^{\otimes s}$ is the range of the Jones--Wenzl projection $p_s\in End(u^{\otimes s})$. We denote by $m_k(s)$ the multiplicity of $r_s$ in $u^{\otimes k}$, with $m_k(0)=Cat_k$.

\begin{theorem}\label{thm:jw}
Let $H\in M_N(\T)$ be complex Hadamard, with $N\geq4$, let $\psi=tr\circ\pi:C(S_N^+)\to\mathbb C$, and set $\Phi_s=(id\otimes\psi)(r_s)\in M_{d_s}(\mathbb C)$. Then:
\begin{enumerate}
\item $T_k\simeq\bigoplus_{s=0}^km_k(s)\Phi_s$, with $\Phi_0=1$, $\Phi_1=0$, and $0\leq\Phi_s\leq1$.

\item $c_k=\sum_sm_k(s)e_s$, where $e_s$ is the multiplicity of the eigenvalue 1 of $\Phi_s$. In particular $G_H=S_N^+$ iff $\|\Phi_s\|<1$ for all $s\geq2$.

\item The spectrum of $T_k$ is contained in that of $T_{k+1}$, with multiplicities. If $c_k=Cat_k$, the largest eigenvalue of $T_k$ below 1 is $\max_{2\leq s\leq k}\|\Phi_s\|$, which is nondecreasing in $k$.

\item With respect to $\psi$, the law of $\chi$ is $(1-\frac{1}{N})\delta_0+\frac{1}{N}\delta_N$, for any $H$. In particular $Tr(\Phi_s)$ depends only on $N$ and $s$.
\end{enumerate}
\end{theorem}

\begin{proof}
By definition $(T_k)_{ij}=\psi(u_{i_1j_1}\ldots u_{i_kj_k})$, so $T_k=(id\otimes\psi)(u^{\otimes k})$, and decomposing $u^{\otimes k}$ over $S_N^+$, with a unitary intertwiner, gives (1). Here $\Phi_0=\psi(1)=1$, and $(id\otimes\psi)(u)=(tr(P_{ij}))_{ij}=\frac{1}{N}\mathbb I$ is the projection onto the constant vector, which spans $r_0\subset u$, so $\Phi_1=0$. The bounds $0\leq\Phi_s\leq1$ come from $0\leq T_k\leq1$. Now (2) follows from (1), and from $c_k=\dim\ker(T_k-1)$, and (3) follows from $u^{\otimes k}\subset u^{\otimes k}\otimes u=u^{\otimes k+1}$. As for (4), since $P_{ii}=Proj(1)$ for any $i$, we have $Tr(T_p)=tr((\sum_iP_{ii})^p)=N^{p-1}$ for any $p$, and the characters of the $r_s$ being polynomials in $\chi$, this gives the last assertion.
\end{proof}

Thus $c_k$ only sees the eigenvalue 1, while the whole spectrum of $T_k$ is encoded in the matrices $\Phi_s$, whose traces are universal, and whose norms are the relevant quantities. For the Di\c t\u a matrices the integers $e_s$ give a sparser encoding of the $c_k$. For instance $e_2,e_3,\ldots=1,0,1,0$ for $F_4^s$, $1,1,4,9,27$ for $F_{8,5b}$ and $F_{8,5c}$, $2,4,16,55,207^*$ for $F_{8,5a}$, and $1,1,5,15$ for $\mathrm{had}_{16.3}$. For the wild matrices, on the other hand, all the integers $e_s$ with $s\geq2$ vanish, as far as computed, and the relevant information is in the norms $\|\Phi_s\|$, computed as follows, the spectral gap discussed at the end of this section being defined in terms of these norms.

\medskip

\noindent\textbf{Computation.} The range of $p_k$ is the intersection of the kernels of the lowering maps, namely the merging of 2 adjacent legs, and the summation over a leg, which are the adjoints of the generating intertwiners of $S_N^+$. Equivalently, it is the intersection of the space of vectors whose leg sums all vanish, given by the projection $(1-\frac{1}{N}\mathbb I)^{\otimes k}$, with the space of vectors vanishing on the words having 2 equal adjacent letters. Both these projections are explicit, and commute with $T_k$, and alternating between them converges geometrically, at a rate of about $0.63$ per sweep, as it must by von Neumann's alternating projection theorem, the limit being the projection onto the intersection of the 2 ranges. This range is small, of dimension $d_k$, for instance $d_k=19,71,265,989,3691,13775,51409$ at $N=6$ and $k=2,\ldots,8$, versus $6^8\simeq1.7\cdot10^6$ for the ambient space, but it has no simple explicit basis, and this is why we work in the ambient space. Note also that this range is not invariant under the rotation of the legs, the merging of the last leg with the first one being not among the lowering maps, so that, unlike for the annular pieces of section \ref{sec:jw}, the rotation sectors cannot be used here. We then compute $\|\Phi_k\|$ by restarted block Krylov on the range of $p_k$, followed by an accurate final projection and an exact Rayleigh--Ritz step. In practice, the Krylov blocks have $N$ vectors, with 3 to 5 steps between the restarts, and at most 30 restarts, the iteration stopping when the residuals of the 4 top Ritz vectors are below $10^{-7}$, while for $N^k\leq2000$ the operator is simply diagonalized. The tables below were computed by \texttt{crocop.py} and \texttt{serval.py}. The vectors still have dimension $N^k$, but no fixed vectors need to be stored, and level 8 at $N=6$, in dimension $6^8\simeq1.7\cdot10^6$, takes about 15 minutes on a laptop. At low levels, the values agree with the full diagonalization of $T_k$ in the rotation sectors, done by the programs \texttt{spectro.py} and \texttt{hw.py}, which confirms the decomposition (1) to precision $10^{-15}$. Also, the norms at the various levels are computed independently, and since $\|\Phi_k\|$ is nondecreasing in $k$, by Theorem \ref{thm:annular} below, this provides a check on the numerics, all the unstarred values being indeed nondecreasing. Note also that the top eigenvalues are often degenerate, as for instance $\|\Phi_2(T_6)\|=\frac{1}{4}$, of multiplicity at least 4, and this is one reason for using blocks of vectors, rather than single vectors. For the converged norms explained below, each Ritz value $\theta$ comes with its residual $\|\Phi_kv-\theta v\|$, which bounds the distance from $\theta$ to the spectrum of $\Phi_k$, and with its leak, the norm of the component of $v$ outside the range of $p_k$, which measures the accuracy of the projections, and stays below $10^{-12}$ in all cases. A Ritz value is a lower bound for $\|\Phi_k\|$, up to these small errors, but a Krylov method alone gives no upper bound, an eigenvector nearly orthogonal to the starting block being possible in principle, so the values below, although reliable, are not proofs. At orders 6 and 7, we find:

\medskip

\begin{center}
\small
\begin{tabular}{lcccccccc}
\toprule
matrix & $\|\Phi_2\|$ & $\|\Phi_3\|$ & $\|\Phi_4\|$ & $\|\Phi_5\|$ & $\|\Phi_6\|$ & $\|\Phi_7\|$ & $\|\Phi_8\|$ & $\|\Phi_9\|$\\
\midrule
$T_6$ & 0.2500 & 0.4284 & 0.6189 & 0.6189 & 0.6510 & 0.6510 & 0.6650 & 0.6657\\
$BF_6$ & 0.6335 & 0.6335 & 0.6667 & 0.6667 & 0.6667 & 0.6667 & 0.6667 & $0.6667^*$\\
$H_6^q$, $q$ random & 0.4019 & 0.4019 & 0.6733 & 0.6733 & 0.6733 & 0.6733 & 0.6812 & 0.6812\\
$H_6^q$, $q=1$ & 0.3704 & 0.3704 & 0.6667 & 0.6667 & 0.6667 & 0.6667 & 0.6667 & 0.6667\\
$BN_6^q$, $q$ random & 0.5008 & 0.5008 & 0.6922 & 0.6922 & 0.6922 & 0.6922 & 0.6999 & 0.6999\\
$BN_6^q$, $q=i$ ($D_6$) & 0.3704 & 0.3704 & 0.6667 & 0.6667 & 0.6667 & 0.6667 & 0.6667 & 0.6667\\
$K_6$, random & 0.7416 & 0.7416 & 0.7946 & 0.7946 & 0.7946 & 0.7946 & 0.7946 & 0.7946\\
$P_7^q$, $q$ random & 0.3528 & 0.5303 & 0.5619 & 0.5870 & 0.5942 & 0.5942 & 0.5942 &\\
$P_7^q$, $q=1$ & 0.3528 & 0.5702 & 0.5702 & 0.5702 & 0.5706 & 0.5706 & 0.5710 &\\
$R_7^q$, $q$ random & 0.3528 & 0.5469 & 0.5533 & 0.5746 & 0.5927 & 0.5927 & 0.5927 &\\
\bottomrule
\end{tabular}
\end{center}

\medskip

\noindent All the values in these tables are Ritz values, and so lower bounds for the norms. At order 6, all the values come from the converged computation of the program \texttt{sefo.py}, explained below, with residuals below $10^{-9}$, except for $BF_6$ at level 8, with residual $6\cdot10^{-8}$, and at level 9, where the iteration was stopped with residual $5\cdot10^{-5}$, this value being starred. At orders 7 to 12, all the values come from the converged computation of the program \texttt{sandcat.py}, explained below, with residuals below $3\cdot10^{-9}$. At orders 8 to 12, where the norms were computed up to level 7 or 6, we find:

\medskip

\begin{center}
\small
\begin{tabular}{lcccccc}
\toprule
matrix & $\|\Phi_2\|$ & $\|\Phi_3\|$ & $\|\Phi_4\|$ & $\|\Phi_5\|$ & $\|\Phi_6\|$ & $\|\Phi_7\|$\\
\midrule
$F_{8,4}$, random & 0.8934 & 0.8934 & 0.8934 & 0.8934 & 0.8934 & 0.8934\\
$F_{8,2}$, random & 0.6218 & 0.6218 & 0.6377 & 0.6377 & 0.6377 & 0.6377\\
$K_{8,2}$, random & 0.4519 & 0.4519 & 0.5658 & 0.5658 & 0.5658 & 0.5658\\
$M_{8,1}$, random & 0.8585 & 0.8585 & 0.8585 & 0.8585 & 0.8585 & 0.8585\\
$I_{8,12}$ & 0.4648 & 0.4648 & 0.5529 & 0.5529 & 0.5529 & 0.5529\\
$K_{8,30}$ & 0.4771 & 0.4771 & 0.5412 & 0.5412 & 0.5412 & 0.5412\\
$L_{8,30}$ & 0.9410 & 0.9410 & 0.9457 & 0.9457 & 0.9457 & 0.9457\\
$H_9$ & 0.4174 & 0.4205 & 0.4415 & 0.4415 & 0.4442 & 0.4444\\
$BN_9$ & 0.6187 & 0.7466 & 0.7466 & 0.7466 & 0.7466 & 0.7466\\
$BN_{10}$ & 0.5641 & 0.5641 & 0.6484 & 0.6484 & 0.6484 & 0.6484\\
$N_{11}$ & 0.2805 & 0.3819 & 0.3819 & 0.3819 & 0.3819 & \\
$P_{12}^1$ & 0.1852 & 0.3086 & 0.3992 & 0.3992 & 0.4326 & \\
\bottomrule
\end{tabular}
\end{center}

\medskip

\noindent\textbf{Counting.} By Theorem \ref{thm:jw} (2), if $c_{k-1}=Cat_{k-1}$, then $c_k=Cat_k+e_k$. Thus, at the next level, we only need to know whether $\Phi_k$ has the eigenvalue 1, and the Krylov norms, which are only lower bounds, can be replaced by the probabilistic count of section \ref{sec:conv}, run on the range of $p_k$ instead of $S^\perp$, with the projection of the computation above. This is the module \texttt{jwcount.py}, requiring neither seeds nor fixed vectors, and a block of $b=8$ vectors, or even $b=2$, suffices, a fixed vector, if any, being detected with high probability by any random block, after which the block is doubled until it exceeds twice the count. As a test, we recover $e_2,\ldots,e_5=1,1,4,9$ for $F_{8,5b}$, with blocks of 2 vectors already detecting the fixed vectors at levels 4 and 5. The cost is about 5 minutes for level 8 at $N=6$, and 10 minutes for level 9, in dimension $6^9\simeq10^7$, with $b=2$. With the programs \texttt{jaguar.py}, \texttt{aerts.py} and \texttt{lennox.py}, we find $e_k=0$ in all cases, namely $e_6$ for $P_7^1$, $I_{8,12}$, $K_{8,30}$, $L_{8,30}$, $BN_9$, $N_{11}$ and $P_{12}^1$, then $e_7$ for $P_7^q$, $P_7^1$, $R_7^q$, the 7 matrices of order 8, $H_9$, $BN_9$ and $BN_{10}$, then $e_8$ for $P_7^q$, $P_7^1$, $R_7^q$, and $e_8$, $e_9$ for the 7 matrices of order 6. This gives the probabilistic values $c_6=132^*$, $c_7=429^*$, then $c_8=1430^*$ at orders 6 and 7, and $c_9=4862^*$ at order 6, which are confirmed by the converged norms below, except for $BF_6$ at level 9, and are therefore written without stars in section \ref{sec:level5}. Level 7 at $N=10$, in dimension $10^7$, takes about 30 minutes on a laptop, and level 8 at $N=7$, in dimension $7^8\simeq5.8\cdot10^6$, about 6 minutes, while level 8 at $N=8$, in dimension $8^8\simeq1.7\cdot10^7$, is beyond the reach of a laptop with 8GB of memory.

\medskip

\noindent\textbf{Converged norms.} In order to certify the values of the tables, and so remove the stars of section \ref{sec:level5}, we use as well a Lanczos iteration run to convergence on the range of $p_k$, with the Paige bound on the residual. This is implemented in the module \texttt{gapcheck.py}, which also computes the multiplicities $e_k$ by a block method of LOBPCG type, with locking, and was used for the orders 24 and 28 in section \ref{sec:orders}. The program \texttt{sefo.py} applies it to the 7 matrices of order 6, up to level 9, in about 1.7 hours for level 9, and its output, in the files \texttt{sefo.log} and \texttt{sefo.json}, gives the order 6 rows of the above tables, with 8 digits, for instance $\|\Phi_9(T_6)\|=0.66574293$, $\|\Phi_8(H_6^q)\|=\|\Phi_9(H_6^q)\|=0.68118741$ and $\|\Phi_k(K_6)\|=0.79456168$ for $k=4,\ldots,9$. The program \texttt{sandcat.py} applies the same method to the matrices of orders 7 to 12, up to levels 8, 7, 6 at orders 7, 8 to 10, 11 to 12, in about 20 minutes for $P_7^1$ at level 8, and 7 minutes for $BN_{10}$ at level 7, in dimension $10^7$, and gives the corresponding rows of the above tables, its output being in the files \texttt{sandcat.log} and \texttt{sandcat.json}. It also computes the norms at level 5 for $P_{20}^1,P_{20}^2,P_{24}^1$, and the multiplicities of section \ref{sec:f85}.

\medskip

\noindent Several observations are in order:
\begin{enumerate}
\item All the norms are below 1. For the unstarred values, this gives $c_k=Cat_k$, in the same sense as the Lanczos computations of section \ref{sec:conv}. For the starred value, that of $BF_6$ at level 9, the equality $c_9=Cat_9$ comes from the count above. Thus all the entries of the tables come with $c_k=Cat_k$, and this extends to level 8 at order 7, and to level 9 at order 6.

\item For 13 of the 22 matrices the profile becomes constant from level 3 or 4 on, as far as computed, and for $F_{8,4}$ and $M_{8,1}$ already from level 2. However, the generic points of $H_6^q$ and $BN_6^q$, constant from level 4 to level 7, jump at level 8, to $0.6812$ and $0.6999$, the previous values $0.6733$ and $0.6922$ remaining eigenvalues of $\Phi_8$, and staying constant at level 9, while $T_6$ jumps at the levels 3, 4, 6, 8 and 9, and $BF_6$, constant at $0.666660$ from level 4 to level 7, increases slightly at level 8, to $0.666664$, staying below $\frac{2}{3}$. Thus the profiles are staircases, and a constant stretch does not mean that $\rho(H)$ is reached. The explanation is given by Theorem \ref{thm:annular} below. Some values are rational, such as $\frac{10}{27},\frac{2}{3}$ for $H_6^1$ and for the $D_6$ point of $BN_6^q$, and $\frac{5}{27},\frac{25}{81},\frac{97}{243}$ for $P_{12}^1$.

\item The profiles of $P_7^q$ and $R_7^q$ increase up to level 6, and then $\|\Phi_8\|=\|\Phi_7\|=\|\Phi_6\|$, to 8 digits, while $\|\Phi_5\|>\|\Phi_4\|$, and $P_7^1$ is constant at levels 3 to 5, and at levels 6, 7, and increases slightly at level 8. Thus the order 6 pattern of odd levels repeating the even ones appears at order 7 too, but only from level 7 on. In the notations below, this means $\|\Psi_7\|\leq\max_{2\leq t\leq6}\|\Psi_t\|=\|\Phi_6\|$, and this maximum is $\|\Psi_6\|$ when $\|\Phi_6\|>\|\Phi_5\|$, as is the case numerically. As for $H_9$, its profile is still slowly increasing at level 7.
\end{enumerate}

\medskip

\subsection{Annular pieces}

The staircases come from the rotation $\rho$ of the legs of $(\mathbb C^N)^{\otimes k}$. By traciality $T_k$ commutes with $\rho$, and also with the $S_N^+$ intertwiners, and so with the whole annular category \cite{jo2}. Since $\rho$ is not an $S_N^+$ intertwiner, it mixes the isotypic components: with $D$ the doubling of the last leg, the map $\Gamma_s=p_{s+1}\rho D:r_s\to r_{s+1}$ satisfies $\Phi_{s+1}\Gamma_s=\Gamma_s\Phi_s$, for any tracial state. For $N=5,6,7$ this map is injective for $s=2,4$, while for $s=3$ it has, together with the similar maps $p_4\rho^jD_i$, with $D_i$ doubling the leg $i$, a common kernel of dimension $N-1$, on which $\Phi_3=0$. Now let $\Lambda_t\subset(\mathbb C^N)^{\otimes t}$ be the space of vectors with all leg sums 0, vanishing on the words with cyclically adjacent equal letters, this latter condition being void for $t=1$, so that $\dim\Lambda_1=N-1$. This is the annular lowest weight space, invariant under $T_t$ and $\rho$, and we denote by $\Psi_t$ the restriction of $T_t$ to it, and by $f_t$ the multiplicity of its eigenvalue 1.

\begin{theorem}\label{thm:annular}
For $N\geq5$, and any tracial state $\psi$ with $\psi(u_{ij})=\frac{1}{N}$, as for instance the one coming from a complex Hadamard matrix, we have, for any $s\geq1$,
$$\Phi_s\simeq\Psi_2\oplus\Psi_3\oplus\ldots\oplus\Psi_s\oplus[s\ {\rm odd}]\cdot0_{N-1}$$
and in particular $d_s=\sum_{t=2}^s\dim\Lambda_t+[s\ {\rm odd}](N-1)$, while $\Phi_0=1$. Thus $\|\Phi_s\|=\max_{2\leq t\leq s}\|\Psi_t\|$ for $s\geq2$, and $e_s=f_2+\ldots+f_s$, which is nondecreasing. If $\psi$ comes from a complex Hadamard matrix $H$, then $G_H=S_N^+$ iff $\|\Psi_t\|<1$ for all $t\geq2$.
\end{theorem}

\begin{proof}
The annular category acts on $V=\bigoplus_k(\mathbb C^N)^{\otimes k}$ in a way compatible with the adjoints, so this representation is completely reducible, and since lowest weight vectors of the same weight $t$ and rotation eigenvalue $\omega$ have the same inner products with all annular diagrams, we obtain $V=\bigoplus_{t,\omega}W_{t,\omega}\otimes\Lambda_t^\omega$, with $W_{t,\omega}$ irreducible annular modules, plus the Temperley--Lieb part at $t=0$. Since $T_k$ commutes with the annular category, it acts as $1\otimes\Psi_{t,\omega}$. Next, restricting to the rectangular category, $W_{t,\omega}(s)$ contains $n_t(s)$ rectangular lowest weight vectors at level $s$, and these generate copies of $r_s$, so $\Phi_s\simeq\bigoplus_{t,\omega}\Psi_{t,\omega}^{\oplus n_t(s)}$, with the numbers $n_t(s)$ determined by
$$\dim W_{t,\omega}(s)=\sum_{j\leq s}n_t(j)m_s(j)\quad,\quad m_s(j)=\binom{2s}{s-j}-\binom{2s}{s-j-1}$$
the formula for $m_s(j)$ coming from the fattening $u\leftrightarrow\frac{1}{2}\otimes\frac{1}{2}$. For $t\geq2$ we have $\dim W_{t,\omega}(s)=\binom{2s}{s-t}$, for the modulus $\sqrt N>2$, by \cite[Proposition 5.3]{jo2}, and by telescoping $n_t(s)=1$ for $s\geq t$. For $t=1$, the module generated by $v\in\Lambda_1$ is spanned by the vectors $\xi_{\pi,B}(v)$, with $\pi\in NC(s)$ and $B$ a marked block carrying $v$, so $\dim W_1(s)\leq\binom{2s-1}{s-1}$, the total number of blocks of the noncrossing partitions of $\{1,\ldots,s\}$. For the equality, we use the shaded picture: the spin planar algebra is the planar algebra of the star graph with one central vertex and $N$ leaves, and the space $\Lambda_1$ of the vectors with sum 0 sits in the zero-box space of the shading opposite to the one of the central vertex, in the kernel of the adjacency summation back to the center. Thus these vectors are lowest weight vectors of weight 0, of the exceptional type, and the corresponding irreducible annular modules have dimension $\frac{1}{2}\binom{2s}{s}=\binom{2s-1}{s-1}$ at level $s\geq1$, by \cite[Proposition 5.21, Theorem 5.23]{jo2}, for the modulus $\sqrt N>2$. Thus the vectors $\xi_{\pi,B}(v)$ are linearly independent, $\dim W_1(s)=\binom{2s-1}{s-1}$, and since $\binom{2s-1}{s-1}=\sum_{j\ odd}m_s(j)$, telescoping gives $n_1(s)=[s\ {\rm odd}]$. On these modules $T_1$ acts on $\Lambda_1$ as $(\psi(u_{ij}))_{ij}=\frac{1}{N}\mathbb I$, which vanishes there, and this gives the term $[s\ {\rm odd}]\cdot0_{N-1}$. Finally, if $\psi$ comes from $H$, then $T_t\geq0$, by Theorem \ref{thm:transp}, so $\|\Psi_t\|=1$ means that 1 is an eigenvalue of $\Psi_t$, and this gives the last assertion. Positivity is needed here: the uniform measure on the odd permutations of $S_5$ gives a tracial state with $\psi(u_{ij})=\frac{1}{5}$, which acts by $-1$ on the copy of $\Lambda^4(\mathbb C^5\ominus\mathbb C)$ in level 4, so that a block of norm 1 need not produce the eigenvalue 1 for an arbitrary tracial state.
\end{proof}

We checked the dimensions $\dim W_t(s)$ directly, by generating the modules at $N=6$, with the program \texttt{wdims.py}, obtaining $1,3,10,35,126$ for $t=1$, $1,6,28,120$ for $t=2$ and $1,8,45$ for $t=3$, at $s=t,\ldots,5$, and the decomposition itself numerically, with the programs \texttt{annular.py} and \texttt{annular2.py}, with all the spectra matched to $10^{-8}$ for $s\leq4$, and the dimension formula for $d_s$ for $N=5,6,7,8$ and $s\leq5$. At $N=4$ the statement fails, as it must, since $F_4^s$ has $e_2,e_3=1,0$. For the Haar state of a quantum group $G\subset S_N^+$, where $T_k$ is the projection onto the fixed vectors, Theorem \ref{thm:annular} reduces to the formula $e_s=f_2+\ldots+f_s$, which is the annular decomposition of the planar algebra $Fix(G)$ due to Jones \cite{jo2}, written in $S_N^+$ terms. The new point is the spectral statement for the tracial states with $\psi(u_{ij})=\frac{1}{N}$, coming from the annular decomposition of the whole spin planar algebra. The integers $f_t$ are the dimensions of the lowest weight spaces of the planar algebra \cite{jo2}, for instance $f_2,f_3,f_4=1,0,3$ for $F_{8,5b}$, $2,1,8$ for $F_6(r,s)$, and $0,0,1$ for $S_6$, the crossing. Note also that by Theorem \ref{thm:jw} (4) we have $Tr(\Phi_s)=\frac{1}{N}(d_s+(-1)^s(N-1))$, which gives $4,11,45,164$ at $N=6$. The annular pieces, computed by the program \texttt{psinorms.py}, are as follows:

\medskip

\begin{center}
\begin{tabular}{lcccc}
\toprule
matrix & $\|\Psi_2\|$ & $\|\Psi_3\|$ & $\|\Psi_4\|$ & $\|\Psi_5\|$\\
\midrule
$T_6$ & 0.2500 & 0.4284 & 0.6189 & 0.5915\\
$BF_6$ & 0.6335 & 0.5095 & 0.6667 & 0.6272\\
$H_6^q$, $q$ random & 0.4019 & 0.3437 & 0.6733 & 0.6123\\
$BN_6^q$, $q$ random & 0.5008 & 0.3986 & 0.6922 & 0.6770\\
$K_6$, random & 0.7416 & 0.5835 & 0.7946 & 0.7093\\
\bottomrule
\end{tabular}
\end{center}

\medskip

\noindent At order 6 the odd pieces are smaller than the even ones, except for $\Psi_3(T_6)$, which explains the staircases. Computationally, $\Lambda_t$ is smaller than $r_t$, as for instance $47$ versus $71$ at $t=3$, $N=6$, and splits further into rotation sectors. 

\subsection{The spectral gap}

The annular pieces suggest a quantitative, and therefore more stable, version of the wildness question of section \ref{sec:nondita}. For $H\in M_N(\T)$ with $N\geq4$ and $c_2=2$, consider the spectral gap $\rho(H)=\sup_{s\geq2}\|\Phi_s\|$. The question is then whether $\rho(H)<1$, which would mean that the truncated integrals of Theorem \ref{thm:trunc} converge to the Haar integral of $S_N^+$ exponentially fast, uniformly over the irreducible corepresentations. Note that a finite number of levels is not known to decide wildness by itself, a bound on the levels at which the proper quantum subgroups of $S_N^+$ differ from $S_N^+$ being not available. A Brannan type estimate \cite{bra} of $\|\Psi_{t+1}\|$ in terms of the previous annular pieces, combined with computations up to some level, would on the other hand give $\rho(H)<1$ for explicit matrices, and so wildness.

\begin{thebibliography}{999}
\setlength{\itemsep}{0pt plus 1pt}
\raggedright\hypersetup{urlcolor=black,linkcolor=black,citecolor=black}

\bibitem{ade} N. Andruskiewitsch and J. Devoto, Extensions of Hopf algebras, {\em St. Petersburg Math. J.} {\bf 7} (1996), 17--52.

\bibitem{ba3x} T. Banica, Representations of compact quantum groups and subfactors, {\em J. Reine Angew. Math.} {\bf 509} (1999), 167--198.

\bibitem{ba1x} T. Banica, Symmetries of a generic coaction, {\em Math. Ann.} {\bf 314} (1999), 763--780.

\bibitem{ba0} T. Banica, Compact Kac algebras and commuting squares, {\em J. Funct. Anal.} {\bf 176} (2000), 80--99.

\bibitem{ba2} T. Banica, Quantum automorphism groups of homogeneous graphs, {\em J. Funct. Anal.} {\bf 224} (2005), 243--280.

\bibitem{ba1} T. Banica, The defect of generalized Fourier matrices, {\em Linear Algebra Appl.} {\bf 438} (2013), 3667--3688.

\bibitem{ba5} T. Banica, Truncation and duality results for Hopf image algebras, {\em Bull. Pol. Acad. Sci. Math.} {\bf 62} (2014), 161--179.

\bibitem{banicaHM} T. Banica, Invitation to Hadamard matrices, \arx{1910.06911} (2024).

\bibitem{bb4} T. Banica and J. Bichon, Free product formulae for quantum permutation groups, {\em J. Inst. Math. Jussieu} {\bf 6} (2007), 381--414.

\bibitem{bb2} T. Banica and J. Bichon, Quantum groups acting on 4 points, {\em J. Reine Angew. Math.} {\bf 626} (2009), 75--114.

\bibitem{bb1} T. Banica and J. Bichon, Hopf images and inner faithful representations, {\em Glasg. Math. J.} {\bf 52} (2010), 677--703.

\bibitem{bb3} T. Banica and J. Bichon, Random walk questions for linear quantum groups, {\em Int. Math. Res. Not.} {\bf 2015} (2015), no. 24, 13406--13436.

\bibitem{bbc} T. Banica, J. Bichon and B. Collins, The hyperoctahedral quantum group, {\em J. Ramanujan Math. Soc.} {\bf 22} (2007), 345--384.

\bibitem{bbs} T. Banica, J. Bichon and J.-M. Schlenker, Representations of quantum permutation algebras, {\em J. Funct. Anal.} {\bf 257} (2009), 2864--2910.

\bibitem{bc} T. Banica and B. Collins, Integration over the Pauli quantum group, {\em J. Geom. Phys.} {\bf 58} (2008), 942--961.

\bibitem{bsp} T. Banica and R. Speicher, Liberation of orthogonal Lie groups, {\em Adv. Math.} {\bf 222} (2009), 1461--1501.

\bibitem{ben} K. Beauchamp and R. Nicoara, Orthogonal maximal abelian $*$-subalgebras of the $6\times 6$ matrices, {\em Linear Algebra Appl.} {\bf 428} (2008), 1833--1853.

\bibitem{bia} P. Biane, Some properties of crossings and partitions, {\em Discrete Math.} {\bf 175} (1997), 41--53.

\bibitem{bi0} J. Bichon, Free wreath product by the quantum permutation group, {\em Algebr. Represent. Theory} {\bf 7} (2004), 343--362.

\bibitem{bi1} J. Bichon, Algebraic quantum permutation groups, {\em Asian-Eur. J. Math.} {\bf 1} (2008), 1--13.

\bibitem{bic} J. Bichon, Quotients and Hopf images of a smash coproduct, {\em Tsukuba J. Math.} {\bf 39} (2015), 285--310.

\bibitem{bha} D. Bisch and U. Haagerup, Composition of subfactors: new examples of infinite depth subfactors, {\em Ann. Sci. \'Ecole Norm. Sup.} {\bf 29} (1996), 329--383.

\bibitem{bjo} D. Bisch and V. Jones, Algebras associated to intermediate subfactors, {\em Invent. Math.} {\bf 128} (1997), 89--157.

\bibitem{bjl} D. Bisch, V. Jones and Z. Liu, Singly generated planar algebras of small dimension, Part III, {\em Trans. Amer. Math. Soc.} {\bf 369} (2017), 2461--2476.

\bibitem{bjf} G. Bj\"orck and R. Fr\"oberg, A faster way to count the solutions of inhomogeneous systems of algebraic equations, with applications to cyclic $n$-roots, {\em J. Symbolic Comput.} {\bf 12} (1991), 329--336.

\bibitem{bra} M. Brannan, Reduced operator algebras of trace-preserving quantum automorphism groups, {\em Doc. Math.} {\bf 18} (2013), 1349--1402.

\bibitem{bcf} M. Brannan, A. Chirvasitu and A. Freslon, Topological generation and matrix models for quantum reflection groups, {\em Adv. Math.} {\bf 363} (2020), 106982.

\bibitem{bghp} M. Brannan, D. Gromada, R. Hern\'andez Palomares and N. Priebe, Complex Hadamard matrices: quantum symmetries, equivalence and non-local games, \arx{2608.14858}.

\bibitem{bur} R.D. Burstein, Group-type subfactors and Hadamard matrices, {\em Trans. Amer. Math. Soc.} {\bf 367} (2015), 6783--6807.

\bibitem{cn} W. Camp and R. Nicoara, Subfactors and Hadamard matrices, \arx{0704.1128}.

\bibitem{cwt} M. C\'ardenes Wuttig and J. Tindall, A complete classification of complex Hadamard matrices of order six, \arx{2608.18053} (2026).

\bibitem{cd} O. Chterental and D.\v{Z}. \DJ okovi\'c, On orthostochastic, unistochastic and qustochastic matrices, {\em Linear Algebra Appl.} {\bf 428} (2008), 1178--1201.

\bibitem{bcm} Claude and Teo, On the maximality conjecture for quantum permutation groups, preprint 2026, \href{https://zenodo.org/records/23858973}{zenodo:23858973}.

\bibitem{dit} P. Di\c t\u a, Some results on the parametrization of complex Hadamard matrices, {\em J. Phys. A} {\bf 37} (2004), 5355--5374.

\bibitem{dfk} J. van Dobben de Bruyn, A. Freslon, P.N. Kar, D.E. Roberson and P. Zeman, Free inhomogeneous wreath product of quantum groups, \arx{2504.13826}.

\bibitem{dkr} J. van Dobben de Bruyn, P.N. Kar, D.E. Roberson, S. Schmidt and P. Zeman, Quantum automorphism groups of trees, {\em J. Noncommut. Geom.} {\bf 20} (2026), 489--505.

\bibitem{en} D. Enders, A characterization of semiprojectivity for subhomogeneous $C^*$-algebras, {\em Doc. Math.} {\bf 21} (2016), 987--1049.

\bibitem{ftr} P. Fima and A. Troupel, Generalized free wreath products and their operator algebras, \arx{2504.00596}.

\bibitem{fnw} I. Francis, R. Nicoara and C. Worley, Bismash product commuting squares, preprint.

\bibitem{fre} A. Freslon, Free wreath products with amalgamation, {\em Comm. Algebra} {\bf 51} (2023), 72--94.

\bibitem{gt} GPT and Teo, Regular complex Hadamard matrices of order 7,8, preprint 2026, \href{https://zenodo.org/records/23934849}{zenodo:23934849}.

\bibitem{gre} J. Greene, Hypergeometric functions over finite fields, {\em Trans. Amer. Math. Soc.} {\bf 301} (1987), 77--101.

\bibitem{gro} D. Gromada, Quantum symmetries of Hadamard matrices, {\em Trans. Amer. Math. Soc.} {\bf 377} (2024), 6341--6377.

\bibitem{gj} P. Grossman and V.F.R. Jones, Intermediate subfactors with no extra structure, {\em J. Amer. Math. Soc.} {\bf 20} (2007), 219--265.

\bibitem{ha1} U. Haagerup, Orthogonal maximal abelian $*$-subalgebras of the $n\times n$ matrices and cyclic $n$-roots, in ``Operator algebras and quantum field theory'', International Press (1997), 296--322.

\bibitem{hw} L.M. Higginbotham and C.T. Worley, On some quaternionic Hadamard matrices of small order, \arx{2112.09613}.

\bibitem{ho} S. Hopkins, Sequential deformations of Hadamard matrices and commuting squares, PhD thesis, Univ. of Tennessee (2022).

\bibitem{hn2} S. Hopkins and R. Nicoara, A finiteness result for complex Hadamard matrices in maximal abelian self-adjoint subalgebras of $M_n(\mathbb C)$, preprint.

\bibitem{hn1} S. Hopkins and R. Nicoara, On the structure of the moduli space of commuting squares around the Fourier spin model, preprint.

\bibitem{imp} M. Izumi, S. Morrison and D. Penneys, Quotients of $A_2*T_2$, {\em Canad. J. Math.} {\bf 68} (2016), 999--1022.

\bibitem{jmn} F. Jaeger, M. Matsumoto and K. Nomura, Bose--Mesner algebras related to type II matrices and spin models, {\em J. Algebraic Combin.} {\bf 8} (1998), 39--72.

\bibitem{jo1} V.F.R. Jones, Index for subfactors, {\em Invent. Math.} {\bf 72} (1983), 1--25.

\bibitem{jon} V.F.R. Jones, Planar algebras I, {\em New Zealand J. Math.} {\bf 52} (2021), 1--107.

\bibitem{jo2} V.F.R. Jones, The annular structure of subfactors, {\em Monogr. Enseign. Math.} {\bf 38} (2001), 401--463.

\bibitem{jsu} V.F.R. Jones and V.S. Sunder, Introduction to subfactors, {\em London Math. Soc. Lect. Notes} {\bf 234}, Cambridge Univ. Press (1997).

\bibitem{ka2} B.R. Karlsson, $H_2$-reducible complex Hadamard matrices of order 6, {\em Linear Algebra Appl.} {\bf 434} (2011), 239--246.

\bibitem{kar} B.R. Karlsson, Three-parameter complex Hadamard matrices of order 6, {\em Linear Algebra Appl.} {\bf 434} (2011), 247--258.

\bibitem{liu} W. Liu, A noncommutative de Finetti theorem for boolean independence, {\em J. Funct. Anal.} {\bf 269} (2015), 1950--1994.

\bibitem{mat} R. Mathon, Symmetric conference matrices of order $pq^2+1$, {\em Canad. J. Math.} {\bf 30} (1978), 321--331.

\bibitem{mck} B.D. McKay, Hadamard matrices, data library, \url{https://users.cecs.anu.edu.au/~bdm/data/hadamard.html}.

\bibitem{mwe} D. McNulty and S. Weigert, Isolated Hadamard matrices from mutually unbiased product bases, {\em J. Math. Phys.} {\bf 53} (2012), 032202.

\bibitem{mo2} M. Montgomery, Angle operators of commuting square subfactors, PhD thesis, Vanderbilt Univ. (2022).

\bibitem{mo3} M. Montgomery, Colored planar algebras for commuting squares and applications to Hadamard subfactors, {\em J. Funct. Anal.} {\bf 285} (2023), 110010.

\bibitem{mo1} M. Montgomery, The spectrum of spin model angle operators, \arx{2110.06720}.

\bibitem{nsp} A. Nica and R. Speicher, Lectures on the combinatorics of free probability, {\em London Math. Soc. Lect. Notes} {\bf 335}, Cambridge Univ. Press (2006).

\bibitem{ni} R. Nicoara, Subfactors and Hadamard matrices, {\em J. Operator Theory} {\bf 64} (2010), 453--468.

\bibitem{nw1} R. Nicoara and J. White, Twisted commuting squares, {\em Internat. J. Math.} {\bf 24} (2013), 1350036.

\bibitem{nw2} R. Nicoara and J. White, The defect of a group-type commuting square, {\em Rev. Roumaine Math. Pures Appl.} {\bf 59} (2014), 245--254.

\bibitem{nwh} R. Nicoara and J. White, Analytic deformations of group commuting squares and complex Hadamard matrices, {\em J. Funct. Anal.} {\bf 272} (2017), 3486--3505.

\bibitem{nw4} R. Nicoara and J. White, Second order deformations of group commuting squares and Hadamard matrices, {\em Proc. Amer. Math. Soc.} {\bf 148} (2020), 3967--3974.

\bibitem{nwo} R. Nicoara and C. Worley, A finiteness result for circulant core complex Hadamard matrices, {\em Linear Algebra Appl.} {\bf 571} (2019), 143--153.

\bibitem{nom} K. Nomura, An algebra associated with a spin model, {\em J. Algebraic Combin.} {\bf 6} (1997), 53--58.

\bibitem{ocn} A. Ocneanu, Quantized groups, string algebras and Galois theory for algebras, {\em London Math. Soc. Lect. Notes} {\bf 136} (1988), 119--172.

\bibitem{ono} K. Ono, Values of Gaussian hypergeometric series, {\em Trans. Amer. Math. Soc.} {\bf 350} (1998), 1205--1223.

\bibitem{pal} R. Paley, On orthogonal matrices, {\em J. Math. Phys.} {\bf 12} (1933), 311--320.

\bibitem{pet} M. Petrescu, Existence of continuous families of complex Hadamard matrices of certain prime dimensions and related results, PhD thesis, UCLA (1997).

\bibitem{pop} S. Popa, Orthogonal pairs of $*$-subalgebras in finite von Neumann algebras, {\em J. Operator Theory} {\bf 9} (1983), 253--268.

\bibitem{swa} T. Sano and Y. Watatani, Angles between two subfactors, {\em J. Operator Theory} {\bf 32} (1994), 209--241.

\bibitem{sch} G. Schmidt, Classification results of Hadamard matrices, MS thesis, Univ. of Tennessee (2017).

\bibitem{schn} H.-J. Schneider, Some remarks on exact sequences of quantum groups, {\em Comm. Algebra} {\bf 21} (1993), 3337--3357.

\bibitem{szp} F. Sz\"oll\H{o}si, Parametrizing complex Hadamard matrices, {\em European J. Combin.} {\bf 29} (2008), 1219--1234.

\bibitem{szo} F. Sz\"oll\H{o}si, Exotic complex Hadamard matrices and their equivalence, {\em Cryptogr. Commun.} {\bf 2} (2010), 187--198.

\bibitem{sz4} F. Sz\"oll\H{o}si, Complex Hadamard matrices of order 6: a four-parameter family, {\em J. Lond. Math. Soc.} {\bf 85} (2012), 616--632.

\bibitem{tz1} W. Tadej and K. \.Zyczkowski, A concise guide to complex Hadamard matrices, {\em Open Syst. Inf. Dyn.} {\bf 13} (2006), 133--177.

\bibitem{tz2} W. Tadej and K. \.Zyczkowski, Defect of a unitary matrix, {\em Linear Algebra Appl.} {\bf 429} (2008), 447--481.

\bibitem{tao} T. Tao, Fuglede's conjecture is false in 5 and higher dimensions, {\em Math. Res. Lett.} {\bf 11} (2004), 251--258.

\bibitem{twa} P. Tarrago and J. Wahl, Free wreath product quantum groups and standard invariants of subfactors, {\em Adv. Math.} {\bf 331} (2018), 1--57.

\bibitem{tle} H.N.V. Temperley and E.H. Lieb, Relations between the ``percolation'' and ``colouring'' problem and other graph-theoretical problems associated with regular planar lattices: some exact results for the ``percolation'' problem, {\em Proc. Roy. Soc. London Ser. A} {\bf 322} (1971), 251--280.

\bibitem{wa0} S. Wang, Free products of compact quantum groups, {\em Comm. Math. Phys.} {\bf 167} (1995), 671--692.

\bibitem{wa1} S. Wang, Quantum symmetry groups of finite spaces, {\em Comm. Math. Phys.} {\bf 195} (1998), 195--211.

\bibitem{wen} H. Wenzl, On sequences of projections, {\em C. R. Math. Rep. Acad. Sci. Canada} {\bf 9} (1987), 5--9.

\bibitem{whi} J. White, Isolation and deformation results for commuting squares of finite dimensional $*$-algebras, PhD thesis, Univ. of Tennessee (2014).

\bibitem{wor} C. Worley, Construction and classification results for commuting squares of finite dimensional $*$-algebras, PhD thesis, Univ. of Tennessee (2017).

\bibitem{wo1} S.L. Woronowicz, Compact matrix pseudogroups, {\em Comm. Math. Phys.} {\bf 111} (1987), 613--665.

\bibitem{wo2} S.L. Woronowicz, Tannaka--Krein duality for compact matrix pseudogroups. Twisted SU(N) groups, {\em Invent. Math.} {\bf 93} (1988), 35--76.

\end{thebibliography}

\end{document}
